% Identifying different Transmission Types — ICT Upper Sixth — Cours signé OnBuch+
% Official MINESEC syllabus. Compiler avec : tectonic ICT08-identifying-different-transmission-types.tex
\newcommand{\DOCMATIERE}{ICT}
\newcommand{\DOCNIVEAU}{Upper Sixth}
\newcommand{\DOCMODULE}{Upper Sixth — ICT}
\newcommand{\DOCLECON}{8}
\newcommand{\DOCTITRE}{Identifying different Transmission Types}
% OnBuch+ — préambule « cours signés » ENRICHI (compilé avec tectonic / XeLaTeX)
% Système visuel « Pop » d'OnBuch+ : fond crème (jamais blanc), bordures encre
% épaisses, ombres « dures » décalées, titres Archivo Black en capitales.
% Version MATHÉMATIQUES : graphes (pgfplots/tikz), boîtes pédagogiques
% (définition, propriété, méthode, attention, le savais-tu, exercice résolu,
% exercice, TP, bilan), page de garde et signature OnBuch+ ancrée en bas.
%
% Macros attendues dans main.tex AVANT \input{preamble} :
%   \DOCMATIERE  \DOCNIVEAU  \DOCMODULE  \DOCLECON (numéro)  \DOCTITRE
\documentclass[11pt]{article}

\usepackage{fontspec}
\usepackage[english]{babel}
\usepackage[a4paper,margin=2cm,top=3.4cm,bottom=2.8cm,headheight=1.4cm,footskip=1.3cm]{geometry}
\usepackage{amsmath,amssymb,amsfonts}
\usepackage{mathtools} % \xmapsto, \coloneqq... souvent utilisés par les modèles
\usepackage{cancel} % simplifications barrées (bilans de forces)
% \abs{z} (module, valeur absolue) et \norm{u} (norme), utilisés par les modèles
\providecommand{\abs}{}\let\abs\relax\DeclarePairedDelimiter{\abs}{\lvert}{\rvert}
\providecommand{\norm}{}\let\norm\relax\DeclarePairedDelimiter{\norm}{\lVert}{\rVert}
\usepackage[table]{xcolor} % [table] : fournit aussi \rowcolors (alternance de lignes)
\usepackage{pagecolor}
\usepackage{tikz}
\usetikzlibrary{shapes.symbols,circuits.logic.US,circuits.logic.IEC,arrows.meta,positioning,calc,shapes.geometric,shapes.misc,shapes.arrows,decorations.pathmorphing,decorations.pathreplacing,decorations.markings,patterns,shadows,mindmap,trees,fit,backgrounds,matrix,intersections,angles,quotes}
\usepackage{pgfplots}
\usepackage[version=4]{mhchem} % équations nucléaires \ce{^{235}_{92}U}
\usepackage[european,siunitx]{circuitikz} % schémas électriques
\pgfplotsset{compat=1.18}
\usepgfplotslibrary{fillbetween,groupplots} % aires entre courbes (calcul intégral)
\usepackage{enumitem}
\usepackage{fancyhdr}
% [most] charge toutes les bibliothèques tcolorbox (listings, minted, raster,
% documentation...) et gonfle inutilement la mémoire TeX sur un document long
% (~300 boîtes) : on ne charge que ce qui est réellement utilisé.
\usepackage[skins,breakable]{tcolorbox}
\usepackage{graphicx}
\usepackage{colortbl}
\usepackage{booktabs}
\usepackage{tabularx}
\usepackage{diagbox} % case d'en-tête diagonale des tableaux à double entrée
% Colonnes extensibles centrée / gauche / droite (C, L, R), souvent utilisées par les modèles
\newcolumntype{C}{>{\centering\arraybackslash}X}
\newcolumntype{L}{>{\raggedright\arraybackslash}X}
\newcolumntype{R}{>{\raggedleft\arraybackslash}X}
\usepackage{array}
\usepackage{multicol}
\usepackage{siunitx}
% parse-numbers=false : accepte aussi \SI{2,5 \times 10^{-3}}{...}
\sisetup{locale=UK,output-decimal-marker={.},group-minimum-digits=4,parse-numbers=false}
\DeclareSIUnit\ohm{\ensuremath{\symup{\Omega}}} % U+2126 absent de la police
\DeclareSIUnit\division{div} % réglages d'oscilloscope (ms/div, V/div)
\DeclareSIUnit\year{yr}\DeclareSIUnit\annum{yr}\DeclareSIUnit\Ma{Ma}\DeclareSIUnit\tr{tr} % tours (vitesse de rotation, ex. \SI{1500}{\tr\per\minute})
\usepackage{qrcode}
% \degree : commande fréquemment utilisée par les modèles pour un angle ou une
% température en dehors de \SI{}{} (non fournie par les paquets chargés).
\providecommand{\degree}{^{\circ}}
% \qed (fin de démonstration), utilisé par les modèles sans amsthm
\providecommand{\qed}{\hfill$\square$}
% \barre : double barre des valeurs interdites dans un tableau de variations (tkz-tab)
\providecommand{\barre}{\ensuremath{\|}}
% environnement proof (amsthm non chargé) utilisé par les modèles
\ifdefined\proof\else\newenvironment{proof}[1][Proof]{\par\noindent\textit{#1.}\ }{\qed\par}\fi
\usepackage{lastpage}
\usepackage{hyperref}
\hypersetup{hidelinks}

% ============================================================
% Palette « Pop » OnBuch+ — mêmes hex que la classe P (plus_ui.dart)
% ============================================================
\definecolor{poporange}{HTML}{F59321}
\definecolor{popdark}{HTML}{E07A0C}
\definecolor{popcream}{HTML}{FFF6E8}
\definecolor{popink}{HTML}{1C1714}
\definecolor{poppurple}{HTML}{7A5AE0}
\definecolor{popgreen}{HTML}{2E9E6A}
\definecolor{poppink}{HTML}{E0457B}
\definecolor{popblue}{HTML}{2F7DE1}
\definecolor{popgold}{HTML}{C98A12}
\definecolor{popmuted}{HTML}{8A7F76}
\definecolor{popline}{HTML}{EADFD2}
\definecolor{popgray}{HTML}{8A7F76}
\colorlet{popgrey}{popgray}
% Alias tolérants (le modèle nomme parfois les couleurs différemment)
\colorlet{popwhite}{white}
\colorlet{popblack}{popink}
\colorlet{popred}{poppink}
\colorlet{popyellow}{popgold}
% Teintes claires pour fonds de tableaux / zones
\colorlet{popblueL}{popblue!10!white}
\colorlet{poporangeL}{poporange!14!white}
\colorlet{popgreenL}{popgreen!12!white}
\colorlet{poppurpleL}{poppurple!12!white}
\colorlet{poppinkL}{poppink!12!white}
\colorlet{popgoldL}{popgold!14!white}

\pagecolor{popcream}

% ============================================================
% Typographie
% ============================================================
\defaultfontfeatures{Ligatures=TeX}
\setmainfont{PlusJakartaSans-Medium.ttf}[Path=../../../fonts/,BoldFont=PlusJakartaSans-Bold.ttf]
% Maths en sans-serif (Fira Math) avec chiffres et lettres droites de Plus
% Jakarta, pour que les formules se fondent dans le texte.
\usepackage[math-style=ISO,bold-style=ISO]{unicode-math}
\setmathfont{FiraMath-Regular.otf}[Path=../../../fonts/]
\setmathfont{PlusJakartaSans-Medium.ttf}[Path=../../../fonts/,range={up/num,up/{Latin,latin},\mathup}]
% (icomma retiré : décimales avec point en anglais)
% Caractères mathématiques Unicode absents des polices de texte (Plus Jakarta,
% Archivo Black) : rendus par leur équivalent mathématique, en texte comme en
% formule — sinon ils s'affichent en carré vide (ex. « Arithmétique dans ℤ »).
\usepackage{newunicodechar}
\newunicodechar{ℝ}{\ensuremath{\mathbb{R}}}
\newunicodechar{ℤ}{\ensuremath{\mathbb{Z}}}
\newunicodechar{ℂ}{\ensuremath{\mathbb{C}}}
\newunicodechar{ℕ}{\ensuremath{\mathbb{N}}}
\newunicodechar{ℚ}{\ensuremath{\mathbb{Q}}}
\newunicodechar{𝔻}{\ensuremath{\mathbb{D}}}
\newunicodechar{α}{\ensuremath{\alpha}}
\newunicodechar{β}{\ensuremath{\beta}}
\newunicodechar{γ}{\ensuremath{\gamma}}
\newunicodechar{δ}{\ensuremath{\delta}}
\newunicodechar{ε}{\ensuremath{\varepsilon}}
\newunicodechar{θ}{\ensuremath{\theta}}
\newunicodechar{λ}{\ensuremath{\lambda}}
\newunicodechar{μ}{\ensuremath{\mu}}
\newunicodechar{σ}{\ensuremath{\sigma}}
\newunicodechar{τ}{\ensuremath{\tau}}
\newunicodechar{φ}{\ensuremath{\varphi}}
\newunicodechar{ψ}{\ensuremath{\psi}}
\newunicodechar{ω}{\ensuremath{\omega}}
\newunicodechar{Σ}{\ensuremath{\Sigma}}
% majuscules produites par \MakeUppercase dans les titres (α-aminés -> Α)
\newunicodechar{Α}{\ensuremath{\alpha}}
\newunicodechar{Β}{\ensuremath{\beta}}
\newunicodechar{Γ}{\ensuremath{\Gamma}}
\newunicodechar{Δ}{\ensuremath{\Delta}}
\newunicodechar{Ε}{\ensuremath{\varepsilon}}
\newunicodechar{Θ}{\ensuremath{\Theta}}
\newunicodechar{Λ}{\ensuremath{\Lambda}}
\newunicodechar{Μ}{\ensuremath{\mu}}
\newunicodechar{Τ}{\ensuremath{\tau}}
\newunicodechar{Φ}{\ensuremath{\Phi}}
\newunicodechar{Ψ}{\ensuremath{\Psi}}
\newunicodechar{Ω}{\ensuremath{\Omega}}
\newunicodechar{↦}{\ensuremath{\mapsto}}
\newunicodechar{∈}{\ensuremath{\in}}
\newunicodechar{∉}{\ensuremath{\notin}}
\newunicodechar{∧}{\ensuremath{\wedge}}
\newunicodechar{⇔}{\ensuremath{\Leftrightarrow}}
\newunicodechar{⟺}{\ensuremath{\Longleftrightarrow}}
\newunicodechar{⇒}{\ensuremath{\Rightarrow}}
\newunicodechar{⟹}{\ensuremath{\Longrightarrow}}
\newunicodechar{∘}{\ensuremath{\circ}}
\newunicodechar{∪}{\ensuremath{\cup}}
\newunicodechar{∩}{\ensuremath{\cap}}
\newunicodechar{≡}{\ensuremath{\equiv}}
\newunicodechar{⊂}{\ensuremath{\subset}}
\newunicodechar{⊥}{\ensuremath{\perp}}
\newunicodechar{∃}{\ensuremath{\exists}}
\newunicodechar{∀}{\ensuremath{\forall}}
\newunicodechar{∛}{\ensuremath{\sqrt[3]{\,}}}
\newunicodechar{∜}{\ensuremath{\sqrt[4]{\,}}}
\newunicodechar{⃗}{\ensuremath{{}^{\to}}}
\newunicodechar{ⁿ}{\ifmmode^{n}\else\textsuperscript{\ensuremath{n}}\fi}
\newunicodechar{ˣ}{\ifmmode^{x}\else\textsuperscript{\ensuremath{x}}\fi}
\newunicodechar{ᵇ}{\ifmmode^{b}\else\textsuperscript{\ensuremath{b}}\fi}
\newunicodechar{ʳ}{\ifmmode^{r}\else\textsuperscript{\ensuremath{r}}\fi}
\newunicodechar{ᵘ}{\ifmmode^{u}\else\textsuperscript{\ensuremath{u}}\fi}
\newunicodechar{ʸ}{\ifmmode^{y}\else\textsuperscript{\ensuremath{y}}\fi}
\newunicodechar{ᵃ}{\ifmmode^{a}\else\textsuperscript{\ensuremath{a}}\fi}
\newunicodechar{ᵉ}{\ifmmode^{e}\else\textsuperscript{\ensuremath{e}}\fi}
\newunicodechar{ᵏ}{\ifmmode^{k}\else\textsuperscript{\ensuremath{k}}\fi}
\newunicodechar{ᵖ}{\ifmmode^{p}\else\textsuperscript{\ensuremath{p}}\fi}
\newunicodechar{ᴺ}{\ifmmode^{N}\else\textsuperscript{\ensuremath{N}}\fi}
\newunicodechar{ᶜ}{\ifmmode^{c}\else\textsuperscript{\ensuremath{c}}\fi}
\newunicodechar{ʰ}{\ifmmode^{h}\else\textsuperscript{\ensuremath{h}}\fi}
\newunicodechar{ᵠ}{\ifmmode^{q}\else\textsuperscript{\ensuremath{q}}\fi}
\newunicodechar{ₙ}{\ifmmode_{n}\else\textsubscript{\ensuremath{n}}\fi}
\newunicodechar{ₐ}{\ifmmode_{a}\else\textsubscript{\ensuremath{a}}\fi}
\newunicodechar{ᵢ}{\ifmmode_{i}\else\textsubscript{\ensuremath{i}}\fi}
\newunicodechar{ᵣ}{\ifmmode_{r}\else\textsubscript{\ensuremath{r}}\fi}
\newunicodechar{ₖ}{\ifmmode_{k}\else\textsubscript{\ensuremath{k}}\fi}
\newunicodechar{ᵦ}{\ifmmode_{\beta}\else\textsubscript{\ensuremath{\beta}}\fi}
\newunicodechar{ₓ}{\ifmmode_{x}\else\textsubscript{\ensuremath{x}}\fi}
\newfontfamily\popdisplay{ArchivoBlack-Regular.ttf}[Path=../../../fonts/]
\newfontfamily\poplogo{SpaceGrotesk-Bold.ttf}[Path=../../../fonts/]
\newfontfamily\popsemi{PlusJakartaSans-SemiBold.ttf}[Path=../../../fonts/]
\newfontfamily\popxbold{PlusJakartaSans-ExtraBold.ttf}[Path=../../../fonts/]
\newfontfamily\popxboldls{PlusJakartaSans-ExtraBold.ttf}[Path=../../../fonts/,LetterSpace=3.0]

\setlist[enumerate]{leftmargin=1.7em,itemsep=5pt,topsep=4pt}
\setlist[itemize]{leftmargin=1.7em,itemsep=4pt,topsep=4pt,label={\color{poporange}\textbullet}}
\setlength{\parindent}{0pt}
\setlength{\parskip}{5pt}
\renewcommand{\arraystretch}{1.25}

% pgfplots : style Pop par défaut
\pgfplotsset{
  every axis/.append style={axis line style={popink,line width=1pt},tick style={popink},
    grid style={popline,line width=0.5pt},label style={font=\small},tick label style={font=\footnotesize}},
  popcurve/.style={poporange,line width=1.8pt,no markers,smooth},
}
% Même style disponible dans un \draw TikZ simple (hors pgfplots)
\tikzset{popcurve/.style={poporange,line width=1.8pt}}
\tikzset{popgrid/.style={popline,very thin}}
\colorlet{popcurve}{poporange} % le nom du style est parfois utilisé comme couleur

% ============================================================
% Pill / squiggle
% ============================================================
\newcommand{\poppill}[3]{%
  \tikz[baseline=-0.55ex]\node[draw=popink,line width=1.2pt,rounded corners=2.6pt,%
    fill=#2,inner xsep=6.5pt,inner ysep=3pt]{%
    \popxboldls\fontsize{7}{7}\selectfont\color{#3}\MakeUppercase{#1}};%
}
\newcommand{\popsquiggle}[1]{%
  \tikz[baseline=-0.4ex]{
    \draw[#1,line width=2.6pt,line cap=round]
      plot [smooth] coordinates {(0,0.09) (0.34,0.01) (0.68,0.21) (1.02,0.01) (1.36,0.10)};
  }%
}
% Mot-clé surligné (marqueur orange sous le texte)
\newcommand{\cle}[1]{{\popxbold\color{popdark}#1}}

% ============================================================
% En-tête / pied
% ============================================================
\pagestyle{fancy}
\fancyhf{}
\renewcommand{\headrulewidth}{0pt}
\renewcommand{\footrulewidth}{0pt}
\lhead{\raisebox{-0.15ex}{\poplogo\fontsize{13}{13}\selectfont\color{popink}OnBuch\color{poporange}+}}
\rhead{\poppill{\DOCMATIERE\ \textbullet\ \DOCNIVEAU\ \textbullet\ Chapter \DOCLECON}{white}{popink}}
\lfoot{\popxbold\fontsize{7.6}{7.6}\selectfont\color{popmuted}\MakeUppercase{Signed course OnBuch+}\ \textbullet\ {\popsemi\DOCTITRE}}
\rfoot{\tikz[baseline=-0.6ex]\node[rounded corners=8.5pt,fill=popink,minimum height=17pt,minimum width=17pt,inner xsep=5pt,inner sep=0]{\popxbold\fontsize{8}{8}\selectfont\color{popcream}\thepage\,/\,\pageref*{LastPage}};}

% ============================================================
% Titres
% ============================================================
\newcommand{\chaptitre}[2]{%
  \begin{tcolorbox}[enhanced,colback=poporange,colframe=popink,boxrule=2.6pt,arc=11pt,
      drop shadow=popink,left=15pt,right=15pt,top=12pt,bottom=11pt,before skip=3pt,after skip=18pt]
    \poppill{#2}{popink}{popcream}\par\vspace{9pt}
    {\raggedright\hyphenpenalty=10000\exhyphenpenalty=10000\popdisplay\fontsize{22}{24}\selectfont\color{popcream}\MakeUppercase{#1}\par}
  \end{tcolorbox}
}
% Section numérotée : pastille orange avec le numéro + titre Archivo
\newcounter{popsec}
\newcommand{\coursec}[1]{%
  \stepcounter{popsec}\setcounter{popsub}{0}%
  \par\vspace{16pt}\noindent
  \begin{minipage}{\linewidth}
  \tikz[baseline=-0.7ex]\node[draw=popink,line width=1.6pt,fill=poporange,rounded corners=4pt,minimum size=22pt,inner sep=0]{\popdisplay\fontsize{12}{12}\selectfont\color{popcream}\thepopsec};%
  \hspace{8pt}{\popdisplay\fontsize{15}{17}\selectfont\color{popink}\MakeUppercase{#1}}\par
  \vspace{4pt}\noindent{\color{poporange}\rule{36pt}{3.6pt}}
  \end{minipage}\par\nopagebreak\vspace{8pt}
}
\newcounter{popsub}
\newcommand{\courssub}[1]{%
  \stepcounter{popsub}%
  \par\vspace{9pt}\noindent{\popxbold\fontsize{11.5}{13}\selectfont\color{popink}{\color{poporange}\thepopsec.\thepopsub}\hspace{6pt}#1}\par\nopagebreak\vspace{4pt}
}

% ============================================================
% Boîtes pédagogiques — même recette : fond blanc, bordure épaisse colorée,
% ombre dure encre, badge plein. Titre optionnel [#1].
% ============================================================
\tcbset{popbase/.style={breakable,enhanced,colback=white,boxrule=2.3pt,arc=9pt,
  shadow={3pt}{-3pt}{0pt}{popink},left=14pt,right=14pt,top=11pt,bottom=11pt,before skip=11pt,after skip=15pt,
  fontupper=\popsemi\fontsize{10.6}{14.5}\selectfont\color{popink},
  fontlower=\popsemi\fontsize{10.6}{14.5}\selectfont\color{popink}}}
% Coche dessinée (le glyphe ✓ n'existe pas dans les polices du document)
\newcommand{\popcheck}[1]{\tikz[baseline=-0.5ex]\draw[#1,line width=1.6pt,line cap=round,line join=round](-0.13,0.0)--(-0.04,-0.09)--(0.13,0.1);}
\providecommand{\coche}{\popcheck{popgreen}}
\providecommand{\resultat}[1]{\par\smallskip\noindent\fcolorbox{popgreen}{popgreenL}{\parbox{\dimexpr\linewidth-2\fboxsep-2\fboxrule\relax}{\textbf{Result.} #1}}\par\smallskip}
\newcommand{\popboxhead}[3]{\poppill{#1}{#2}{white}\ifx\relax#3\relax\else\hspace{6pt}{\popxbold\fontsize{10.6}{12}\selectfont\color{popink}#3}\fi\par\vspace{7pt}}

\newtcolorbox{aretenir}[1][]{popbase,colframe=popgreen,before upper={\popboxhead{Key points}{popgreen}{#1}}}
\newtcolorbox{exemplebox}[1][]{popbase,colframe=poporange,before upper={\popboxhead{Exemple}{poporange}{#1}}}
\newtcolorbox{definition}[1][]{popbase,colframe=popblue,before upper={\popboxhead{Definition}{popblue}{#1}}}
\newtcolorbox{propriete}[1][]{popbase,colframe=poppurple,before upper={\popboxhead{Property}{poppurple}{#1}}}
\newtcolorbox{methode}[1][]{popbase,colframe=popgold,before upper={\popboxhead{Method}{popgold}{#1}}}
\newtcolorbox{attention}[1][]{popbase,colframe=poppink,before upper={\popboxhead{Warning}{poppink}{#1}}}
\newtcolorbox{experience}[1][]{popbase,colframe=popblue,colback=popblueL,before upper={\popboxhead{Experiment}{popblue}{#1}}}
% « Le savais-tu ? » : carte violette pleine (contexte, culture, Cameroun)
\newtcolorbox{savaistu}[1][]{popbase,colback=poppurple,colframe=popink,
  fontupper=\popsemi\fontsize{10.4}{14.2}\selectfont\color{white},
  before upper={\poppill{Did you know?}{white}{poppurple}\ifx\relax#1\relax\else\hspace{6pt}{\popxbold\color{white}#1}\fi\par\vspace{7pt}}}
% Exercice résolu : énoncé en haut, \tcblower puis la solution sur fond crème
\newcounter{popexo}\newcounter{popexr}
\newtcolorbox{exoresolu}[1][]{popbase,colframe=poporange,colbacklower=popcream,
  segmentation style={popink,line width=1.2pt,dashed},
  before upper={\stepcounter{popexr}\popboxhead{Worked example \thepopexr}{poporange}{#1}},
  before lower={\poppill{Solution}{popink}{popcream}\par\vspace{6pt}}}
% Exercice (non résolu en place) — la correction est dans la partie Corrigés
\newtcolorbox{exercice}[1][]{popbase,colframe=popink,
  before upper={\stepcounter{popexo}\popboxhead{Exercise \thepopexo}{popink}{#1}}}
% Corrigé
\newtcolorbox{corrige}[1][]{popbase,colframe=popgreen,colback=popgreenL,
  before upper={\popboxhead{Solution}{popgreen}{#1}}}
% Objectifs / prérequis / bilan
\newtcolorbox{objectifs}{popbase,colframe=popink,colback=poporangeL,
  before upper={\popboxhead{Chapter objectives}{popink}{}}}
\newtcolorbox{prerequis}{popbase,colframe=popink,colback=popblueL,
  before upper={\popboxhead{Prerequisites}{popblue}{}}}
\newtcolorbox{bilan}{popbase,colframe=popink,colback=white,boxrule=2.8pt,
  before upper={\popboxhead{Fiche bilan}{poporange}{L'essentiel à connaître pour l'examen}}}
% Figure encadrée avec légende
\newtcolorbox{popfigure}[1][]{enhanced,breakable=false,colback=white,colframe=popline,boxrule=1.4pt,arc=8pt,
  left=10pt,right=10pt,top=10pt,bottom=8pt,before skip=10pt,after skip=12pt,halign=center,
  lower separated=false,halign lower=center,
  fontlower=\popsemi\fontsize{9}{11.5}\selectfont\color{popmuted},
  #1}
\newcommand{\legende}[1]{\par\vspace{6pt}{\popsemi\fontsize{9}{11.5}\selectfont\color{popmuted}#1}}

% ============================================================
% Signature OnBuch+
% ============================================================
\newcommand{\poplogoinline}[1]{{\poplogo\fontsize{#1}{#1}\selectfont\color{popink}OnBuch\color{poporange}+}}
\newcommand{\signatureonbuch}{%
  \begin{tcolorbox}[enhanced,colback=popink,colframe=popink,boxrule=2.2pt,arc=10pt,
      drop shadow=poporange,left=14pt,right=14pt,top=10pt,bottom=10pt,before skip=0pt,after skip=0pt]
    \tikz[baseline=-0.7ex]\node[circle,fill=popgreen,draw=popcream,line width=1.3pt,minimum size=22pt,inner sep=0]{\popcheck{white}};%
    \hspace{9pt}\begin{minipage}[c]{0.62\linewidth}
      {\popxbold\fontsize{12}{13}\selectfont\color{popcream}Signed\ \textbullet\ {\poplogo OnBuch\color{poporange}+}}\par\vspace{2pt}
      {\raggedright\popsemi\fontsize{8}{10.5}\selectfont\color{popline}\MakeUppercase{OnBuch+ teaching team}\\\MakeUppercase{Follows the official MINESEC syllabus}\par}
    \end{minipage}\hfill
    {\poplogo\fontsize{22}{22}\selectfont\color{popcream}OnBuch\color{poporange}+}
  \end{tcolorbox}%
}

% ============================================================
% Page de garde
% #1 titre · #2 matière · #3 niveau · #4 module · #5 n° leçon · #6 durée ·
% #7 description / objectifs (texte ou itemize)
% ============================================================
\newcommand{\pagegarde}[7]{%
  \thispagestyle{empty}
  \newgeometry{a4paper,margin=1.7cm,top=1.6cm,bottom=1.4cm}%
  \begin{tikzpicture}[remember picture,overlay]
    \fill[poporange!18] ([xshift=-3.2cm,yshift=-3.6cm]current page.north east) circle (4.6cm);
    \fill[poppurple!14] ([xshift=2.4cm,yshift=7.6cm]current page.south west) circle (3.8cm);
  \end{tikzpicture}%
  {\poplogo\fontsize{30}{30}\selectfont\color{popink}OnBuch\color{poporange}+}\hfill
  \raisebox{0.6ex}{\poppill{Signed course \textbullet\ Premium}{popink}{popcream}}\par
  \vspace{1pt}\popsquiggle{poppurple}\par
  \vspace{8pt}
  \begin{tcolorbox}[enhanced,colback=poporange,colframe=popink,boxrule=2.8pt,arc=13pt,
      drop shadow=popink,left=18pt,right=18pt,top=12pt,bottom=13pt,after skip=10pt]
    \poppill{#2 \textbullet\ #3}{popink}{popcream}\hspace{5pt}\poppill{#4}{white}{popink}\par\vspace{8pt}
    {\popdisplay\fontsize{14}{14}\selectfont\color{popink}CHAPTER #5}\par\vspace{4pt}
    {\raggedright\hyphenpenalty=10000\exhyphenpenalty=10000\popdisplay\fontsize{25}{27}\selectfont\color{popcream}\MakeUppercase{#1}\par}
    \vspace{8pt}
    {\popxbold\fontsize{9.5}{9.5}\selectfont\color{popink}\MakeUppercase{Suggested time: #6}}
  \end{tcolorbox}
  \begin{tcolorbox}[enhanced,colback=white,colframe=popink,boxrule=2.2pt,arc=10pt,
      drop shadow=popink,left=16pt,right=16pt,top=10pt,bottom=10pt,after skip=8pt]
    \poppill{In this chapter}{poppurple}{white}\par\vspace{5pt}
    \popsemi\fontsize{9.3}{12.6}\selectfont\color{popink}#7
  \end{tcolorbox}
  \noindent\begin{minipage}[t]{0.487\linewidth}
    \begin{tcolorbox}[enhanced,colback=popgreen,colframe=popink,boxrule=2.2pt,arc=10pt,
        drop shadow=popink,left=11pt,right=11pt,top=10pt,bottom=10pt,height=3.1cm,valign=center]
      \tcbox[colback=white,colframe=popink,boxrule=1.3pt,arc=5pt,boxsep=0pt,nobeforeafter,
        left=3pt,right=3pt,top=3pt,bottom=3pt,valign=center]{\qrcode[height=1.9cm]{https://play.google.com/store/apps/details?id=com.luvvix.onbuch&hl=en}}%
      \hspace{8pt}\begin{minipage}[c]{0.5\linewidth}
        {\popdisplay\fontsize{11}{12.5}\selectfont\color{white}\MakeUppercase{Download OnBuch}}\par\vspace{4pt}
        {\raggedright\popsemi\fontsize{8.4}{11}\selectfont\color{white}Free \textbullet\ Google Play\\AI tutor, past papers, results\par}
      \end{minipage}
    \end{tcolorbox}
  \end{minipage}\hfill
  \begin{minipage}[t]{0.487\linewidth}
    \begin{tcolorbox}[enhanced,colback=poppurple,colframe=popink,boxrule=2.2pt,arc=10pt,
        drop shadow=popink,left=11pt,right=11pt,top=10pt,bottom=10pt,height=3.1cm,valign=center]
      \tikz[baseline=-0.7ex]\node[circle,fill=white,draw=popink,line width=1.3pt,minimum size=1.6cm,inner sep=0]{\popdisplay\fontsize{24}{24}\selectfont\color{poppurple}+};%
      \hspace{8pt}\begin{minipage}[c]{0.6\linewidth}
        {\popdisplay\fontsize{11}{12.5}\selectfont\color{white}\MakeUppercase{Go OnBuch+}}\par\vspace{4pt}
        {\raggedright\popsemi\fontsize{8.4}{11}\selectfont\color{white}All signed courses, unlimited AI tutor, PDF sheets — OnBuch+ tab in the app\par}
      \end{minipage}
    \end{tcolorbox}
  \end{minipage}
  \vspace{8pt}
  \popcontenu
  \vfill
  \signatureonbuch
  \restoregeometry
  \newpage
}

% Rangée « Ce cours contient » (page de garde)
\newcommand{\poptile}[3]{% #1 couleur #2 gros texte #3 légende
  \begin{tcolorbox}[enhanced,colback=#1,colframe=popink,boxrule=2pt,arc=9pt,shadow={2.5pt}{-2.5pt}{0pt}{popink},
      left=6pt,right=6pt,top=9pt,bottom=9pt,halign=center,valign=center,height=1.75cm,width=\linewidth,nobeforeafter]
    {\popdisplay\fontsize{12.5}{13}\selectfont\color{white}\MakeUppercase{#2}}\par\vspace{3pt}
    {\centering\popsemi\fontsize{7.8}{9.5}\selectfont\color{white}#3\par}
  \end{tcolorbox}}
\newcommand{\popcontenu}{%
  \par\noindent{\popxboldls\fontsize{7.5}{7.5}\selectfont\color{popmuted}THIS SIGNED COURSE CONTAINS}\par\vspace{7pt}
  \noindent\begin{minipage}[t]{0.235\linewidth}\poptile{popblue}{Course}{complete, illustrated, step by step}\end{minipage}\hfill
  \begin{minipage}[t]{0.235\linewidth}\poptile{poporange}{Methods}{and worked examples}\end{minipage}\hfill
  \begin{minipage}[t]{0.235\linewidth}\poptile{poppink}{Exercises}{exam style + solutions}\end{minipage}\hfill
  \begin{minipage}[t]{0.235\linewidth}\poptile{popgreen}{Summary sheet}{the essentials to remember}\end{minipage}\par}

% Sommaire
\newtcolorbox{sommaire}{enhanced,colback=white,colframe=popink,boxrule=2.2pt,arc=10pt,
  drop shadow=popink,left=18pt,right=18pt,top=13pt,bottom=13pt,before skip=6pt,after skip=20pt,
  before upper={\poppill{Contents}{poppurple}{white}\par\vspace{9pt}\popsemi\fontsize{10.6}{14.5}\selectfont\color{popink}}}

% Signature de fin de cours — poussée tout en bas de la dernière page
\newcommand{\findecours}{%
  \par\vfill\nopagebreak
  \begin{center}\popsquiggle{poporange}\end{center}\vspace{4pt}
  \signatureonbuch
}
\AtBeginDocument{\def\llbracket{\mathopen{[\mkern-3.5mu[}}\def\rrbracket{\mathclose{]\mkern-3.5mu]}}} % intervalles d'entiers (glyphe ⟦ absent de la police)
% Glyphes absents de Fira Math : redéfinis à partir de glyphes disponibles
\AtBeginDocument{%
  \def\setminus{\mathbin{\backslash}}\let\smallsetminus\setminus
  \def\vdots{\mathord{\vbox{\baselineskip3.2pt\lineskiplimit0pt\kern4pt\hbox{.}\hbox{.}\hbox{.}}}}%
  \def\ddots{\mathinner{\mkern1mu\raise6.5pt\hbox{.}\mkern2mu\raise3.6pt\hbox{.}\mkern2mu\raise0.7pt\hbox{.}\mkern1mu}}%
  \def\triangle{\mathord{\scalebox{0.9}{$\symup{\Delta}$}}}%
  \def\nabla{\mathord{\raisebox{\depth}{\scalebox{1}[-1]{$\symup{\Delta}$}}}}%
  \def\checkmark{\popcheck{popdark}}%
  \def\star{\mathbin{\ast}}\def\bigstar{\mathord{\ast}}%
  \def\diamond{\mathbin{\scalebox{0.6}{$\Diamond$}}}%
  \def\bigcup{\mathop{\mathchoice{\raisebox{-0.3em}{\scalebox{1.7}{$\cup$}}}{\scalebox{1.25}{$\cup$}}{\cup}{\cup}}\displaylimits}%
  \def\bigcap{\mathop{\mathchoice{\raisebox{-0.3em}{\scalebox{1.7}{$\cap$}}}{\scalebox{1.25}{$\cap$}}{\cap}{\cap}}\displaylimits}%
  \def\lesseqqgtr{\mathrel{\lessgtr}}\def\bigoplus{\mathop{\oplus}\displaylimits}%
}
\providecommand{\ssi}{if and only if} % abréviation fréquente
\AtBeginDocument{\providecommand{\wideparen}{\overparen}} % arc de cercle

% Fonctions usuelles que les modèles utilisent sans que amsmath les fournisse
\providecommand{\arsinh}{\operatorname{arsinh}}\providecommand{\arcosh}{\operatorname{arcosh}}
\providecommand{\artanh}{\operatorname{artanh}}\providecommand{\arsech}{\operatorname{arsech}}
\providecommand{\arcsch}{\operatorname{arcsch}}\providecommand{\arcoth}{\operatorname{arcoth}}
\providecommand{\arcsinh}{\operatorname{arsinh}}\providecommand{\arccosh}{\operatorname{arcosh}}
\providecommand{\arctanh}{\operatorname{artanh}}\providecommand{\sech}{\operatorname{sech}}
\providecommand{\csch}{\operatorname{csch}}\providecommand{\arcsec}{\operatorname{arcsec}}
\providecommand{\arccsc}{\operatorname{arccsc}}\providecommand{\arccot}{\operatorname{arccot}}
\providecommand{\sgn}{\operatorname{sgn}}\providecommand{\lcm}{\operatorname{lcm}}
\providecommand{\Var}{\operatorname{Var}}\providecommand{\Cov}{\operatorname{Cov}}
\providecommand{\permil}{\ensuremath{{}^{0}\!/_{00}}}\providecommand{\textperthousand}{\ensuremath{{}^{0}\!/_{00}}}

%% FIGFIX-BEGIN v5 — correctifs globaux des figures (docs/onbuch-generation/tools/figfix)
\usepackage{adjustbox}\usepackage{etoolbox}\tcbuselibrary{raster}
% toute figure TikZ/pgfplots plus large que la ligne est réduite à la largeur disponible
\BeforeBeginEnvironment{tikzpicture}{\begin{adjustbox}{max width=\linewidth}}
\AfterEndEnvironment{tikzpicture}{\end{adjustbox}}
% légendes pgfplots lisibles par-dessus les courbes
\pgfplotsset{every axis/.append style={legend style={fill=white,fill opacity=0.92,text opacity=1,draw=black!15,inner sep=3pt}}}
% étiquettes d'arêtes (syntaxe edge["..."]) avec fond blanc
\tikzset{every edge quotes/.append style={fill=white,inner sep=1.2pt}}
%% FIGFIX-END
\begin{document}

\pagegarde{Identifying different Transmission Types}{ICT}{Upper Sixth}{Upper Sixth — ICT}{8}{3 periods}{This chapter explores how data moves from one device to another: the difference between broadband and narrowband transmission, the types of data transmission (serial/parallel, simplex/half-duplex/full-duplex, synchronous/asynchronous), and the guided and unguided media that carry signals. It equips learners to analyse and choose appropriate transmission solutions for real networks, a core GCE Advanced Level ICT competence.
\vspace{4pt}
\begin{multicols}{2}\small
\begin{itemize}
\item By the end of this chapter I can define data transmission and distinguish between baseband, narrowband and broadband transmission.
\item By the end of this chapter I can compare narrowband and broadband transmission in terms of bandwidth, speed, cost and typical applications.
\item By the end of this chapter I can differentiate serial from parallel transmission and state where each is used.
\item By the end of this chapter I can distinguish simplex, half-duplex and full-duplex transmission modes with everyday examples.
\item By the end of this chapter I can explain synchronous and asynchronous transmission and identify their advantages and limitations.
\item By the end of this chapter I can classify transmission media into guided (wired) and unguided (wireless) categories.
\end{itemize}\end{multicols}}

\begin{sommaire}
\begin{multicols}{2}
\begin{enumerate}
\item Fundamentals of Data Transmission and Bandwidth
\item Broadband and Narrowband Transmission (Lesson 28)
\item Types of Data Transmission (Lesson 29)
\item Transmission Media: Guided and Unguided (Lesson 30)
\item Integration activity
\item Methods and worked examples
\item Exercises
\item Solutions to the exercises
\item Summary sheet
\end{enumerate}
\end{multicols}
\end{sommaire}

\begin{objectifs}
\begin{itemize}
\item By the end of this chapter I can define data transmission and distinguish between baseband, narrowband and broadband transmission.
\item By the end of this chapter I can compare narrowband and broadband transmission in terms of bandwidth, speed, cost and typical applications.
\item By the end of this chapter I can differentiate serial from parallel transmission and state where each is used.
\item By the end of this chapter I can distinguish simplex, half-duplex and full-duplex transmission modes with everyday examples.
\item By the end of this chapter I can explain synchronous and asynchronous transmission and identify their advantages and limitations.
\item By the end of this chapter I can classify transmission media into guided (wired) and unguided (wireless) categories.
\item By the end of this chapter I can describe the structure, characteristics and uses of twisted pair, coaxial and fibre optic cables.
\item By the end of this chapter I can describe radio waves, microwaves, infrared and satellite transmission and their applications in Cameroon.
\item By the end of this chapter I can select and justify an appropriate transmission type and medium for a given networking scenario.
\end{itemize}
\end{objectifs}

\begin{prerequis}
\begin{itemize}
\item Definition of a computer network and its components (Lower Sixth: introduction to networks)
\item Concept of data, signals (analogue vs digital) and bits/bytes
\item Basic network topologies (bus, star, ring) seen in Lower Sixth
\item Units of measurement: hertz, bits per second, and their multiples (kHz, MHz, GHz, kbps, Mbps)
\item The role of communication devices such as modems, hubs and switches
\end{itemize}
\end{prerequis}

\newpage

\coursec{Fundamentals of Data Transmission and Bandwidth}

At 7 a.m. in Douala, a student streams a video lesson on her phone while her father complains that the office internet ``crawls'' whenever everyone connects at once. Meanwhile, the radio in the kitchen plays clearly, and the neighbour's satellite dish brings in Canal+ without a single cable. Why do some connections carry huge amounts of data while others struggle, and why do some signals travel through wires while others cross the air? The answers lie in how data is \cle{transmitted}: the \cle{bandwidth} of the channel, the \cle{mode of transmission}, and the \cle{medium} that carries the signal. This chapter unpacks these three pillars of data communication. In this first section, we lay the foundations: what data transmission is, what a transmission system is made of, the difference between analogue and digital signals, and the crucial notion of bandwidth, which prepares us for the study of broadband and narrowband transmission in the next section.

\courssub{What is Data Transmission?}

\textbf{Intuition.} When you send a WhatsApp voice note to a friend in Bamenda, your phone does not send ``your voice'' itself. It converts the sound into numbers, converts those numbers into electrical or radio \emph{signals}, and pushes those signals across a network until they reach your friend's phone, which rebuilds the sound. The whole journey of the information, from your phone to your friend's phone, is a \emph{data transmission}.

\begin{definition}[Data transmission]
\cle{Data transmission} is the process of sending data, in the form of \cle{signals}, from a \cle{sender} (transmitter) to a \cle{receiver} over a \cle{communication channel} (also called a transmission medium), according to agreed rules called \cle{protocols}.
\end{definition}

Data transmission is the foundation of every network: a school computer room, a mobile money transfer with MTN MoMo, a CRTV radio broadcast, or a bank transaction between two branches in Yaoundé all rely on the same basic scheme.

\textbf{Elements of a transmission system.} Every data transmission, however simple or complex, involves five essential elements:

\begin{itemize}
\item \textbf{The sender (transmitter):} the device that originates the data, for example a computer, a smartphone or a radio studio transmitter.
\item \textbf{The receiver:} the device that accepts the data, for example another computer, a phone, or a radio set.
\item \textbf{The signal:} the physical form taken by the data while travelling (electrical voltage on a wire, light pulses in a fibre, radio waves in the air).
\item \textbf{The channel or medium:} the path the signal follows, such as a copper cable, an optical fibre, or the atmosphere.
\item \textbf{The protocol:} the set of rules both parties follow so that the receiver can correctly interpret the signal (how fast bits are sent, how a message starts and ends, how errors are detected).
\end{itemize}

\begin{popfigure}
\begin{center}
\begin{tikzpicture}[font=\small, >=stealth]
\node[draw=popblue, fill=popblueL, rounded corners, minimum width=2.4cm, minimum height=1.1cm, align=center] (s) at (0,0) {\textbf{Sender}\\ (computer)};
\node[draw=popgreen, fill=popgreenL, rounded corners, minimum width=2.4cm, minimum height=1.1cm, align=center] (r) at (10,0) {\textbf{Receiver}\\ (computer)};
\draw[line width=2.2pt, popmuted] (2.2,0) -- (7.8,0);
\node[below=2pt, popmuted] at (5,0) {Medium (cable, fibre, air)};
\draw[->, line width=1.4pt, poporange] (2.4,0.55) -- (7.6,0.55) node[midway, above, popdark]{Signal (data)};
\draw[->, line width=1.1pt, poppurple, dashed] (7.6,-0.55) -- (2.4,-0.55) node[midway, below, popdark]{Acknowledgements / replies};
\node[draw=poppurple, dashed, rounded corners, align=center, text=poppurple] at (5,1.5) {Protocol: rules agreed by both sides};
\end{tikzpicture}
\end{center}
\legende{A basic data transmission system: sender, signal, medium, receiver and protocol.}
\end{popfigure}

\begin{exemplebox}[A mobile money transfer as a transmission system]
When a trader in Buea sends money by Orange Money: the \emph{sender} is her phone, the \emph{signal} is a radio wave carrying coded data, the \emph{medium} is the air between the phone and the Orange antenna, the \emph{receiver} is the operator's server, and the \emph{protocols} are the technical rules (GSM/4G standards, encryption rules) that make the transaction secure and understandable on both sides.
\end{exemplebox}

\courssub{Analogue and Digital Signals}

\textbf{Observation.} Speak into a microphone: the voltage it produces rises and falls smoothly, copying the smooth variation of the sound. Now look at the data travelling inside a network cable: the voltage jumps abruptly between two levels, ``low'' and ``high'', representing 0 and 1. These are the two great families of signals.

\begin{definition}[Analogue and digital signals]
An \cle{analogue signal} is a signal that varies \emph{continuously} over time and can take any value within a range (example: a smooth sine wave carrying a voice on a telephone line or a radio broadcast). A \cle{digital signal} is a signal that takes only a finite number of \emph{discrete} values, usually two levels representing the binary digits 0 and 1.
\end{definition}

\begin{popfigure}
\begin{center}
\begin{tikzpicture}
\begin{axis}[
width=0.9\linewidth, height=5.2cm,
xlabel={Time}, ylabel={Signal level},
xmin=0, xmax=6.3, ymin=-1.6, ymax=1.6,
xtick=\empty, ytick={-1,0,1},
grid=both, grid style={popline!40},
legend style={at={(0.5,-0.28)}, anchor=north, legend columns=2, draw=popline},
axis line style={popmuted}]
\addplot[popblue, very thick, domain=0:6.3, samples=200] {sin(deg(2*x))};
\addlegendentry{Analogue (continuous sine wave)}
\addplot[poporange, very thick, const plot mark mid] coordinates {(0,1) (0.7,1) (0.7,-1) (1.6,-1) (1.6,1) (2.4,1) (2.4,-1) (3.1,-1) (3.1,-1) (3.9,1) (3.9,1) (4.8,-1) (4.8,-1) (5.6,1) (5.6,1) (6.3,1)};
\addlegendentry{Digital (discrete pulses: 0 / 1)}
\end{axis}
\end{tikzpicture}
\end{center}
\legende{An analogue signal varies smoothly; a digital signal jumps between discrete levels.}
\end{popfigure}

\textbf{Why do computers use digital signals?} Computers are built on binary logic: every piece of data (text, images, sound, video) is stored as bits, 0 and 1. Digital signals are therefore natural for computers, and they have decisive advantages:

\begin{itemize}
\item \textbf{Resistance to noise:} a receiver only has to decide ``is this a 0 or a 1?''. Even if a pulse is weakened or slightly distorted, it can be recognised and regenerated perfectly. An analogue signal, once distorted, is degraded forever.
\item \textbf{Easy regeneration:} devices called \emph{repeaters} can rebuild a clean digital signal, allowing transmission over very long distances (this is how CAMTEL's fibre links carry data across the country).
\item \textbf{Easy storage and processing:} digital data can be copied, compressed and encrypted without loss of quality.
\end{itemize}

\begin{attention}[Analogue is not ``old and useless'']
The physical world is analogue: sound, light and radio waves are continuous. Radio broadcasting (CRTV), older telephone lines and satellite TV still use analogue or analogue-like carrier waves. Modern systems therefore constantly \emph{convert} between the two: a modem (\textbf{mo}dulator--\textbf{dem}odulator) converts digital computer data into analogue-like signals for a telephone or radio channel, and back.
\end{attention}

\courssub{Bandwidth: The Width of the Road}

\textbf{Intuition.} Think of the road between Douala and Yaoundé. A narrow road with one lane lets only a few vehicles pass per hour; a wide highway with several lanes lets many more vehicles through. A communication channel behaves the same way: the ``width'' of the channel, measured in frequencies, limits how much data can flow per second.

Every signal can be described by its \cle{frequency}, the number of oscillations per second, measured in \cle{hertz} (Hz). A channel never carries all frequencies equally well: it transmits correctly only a certain \emph{range} of frequencies.

\begin{definition}[Bandwidth]
The \cle{bandwidth} of a transmission channel is the \emph{range of frequencies} that the channel can carry, measured in \cle{hertz} (Hz). If a channel passes frequencies from $f_{\min}$ to $f_{\max}$, its bandwidth is
\[ B = f_{\max} - f_{\min}. \]
For example, a telephone line carrying frequencies from $300\ \mathrm{Hz}$ to $3400\ \mathrm{Hz}$ has a bandwidth of $B = 3400 - 300 = 3100\ \mathrm{Hz}$.
\end{definition}

Bandwidth is directly linked to the \cle{data transfer rate} (also called channel capacity or bit rate), which is the quantity of data transmitted per unit of time, measured in \cle{bits per second} (bps) and its multiples: kbps, Mbps, Gbps.

\begin{propriete}[Bandwidth and data rate]
The \emph{greater the bandwidth} of a channel, the \emph{higher the potential data transfer rate} it can support. A wide channel can carry more signal changes per second, hence more bits per second. (The precise quantitative relationships, such as the Nyquist and Shannon formulas, are beyond the scope of this course; only the qualitative link is examinable.)
\end{propriete}

This explains the father's complaint in our opening story: when many employees share one internet connection, the available capacity is divided among them, so each user gets a smaller share and the connection ``crawls''.

\begin{attention}[Do not confuse bandwidth and data rate]
A very common exam mistake is to treat \emph{bandwidth} and \emph{data rate} as the same thing. They are \textbf{related but different quantities}:
\begin{itemize}
\item \textbf{Bandwidth} is a range of \emph{frequencies}, measured in \textbf{hertz (Hz)}. It is a property of the channel.
\item \textbf{Data rate} (bit rate) is a quantity of \emph{bits per second}, measured in \textbf{bps}. It describes how fast data actually flows.
\end{itemize}
A wide bandwidth \emph{allows} a high data rate, but the actual data rate also depends on the equipment, the encoding and the network congestion. Writing ``bandwidth = 10 Mbps'' in an exam costs marks: say ``bandwidth = 10 MHz'' or ``data rate = 10 Mbps''.
\end{attention}

\begin{methode}[Calculating the bandwidth of a channel]
\begin{enumerate}
\item Identify the lowest frequency $f_{\min}$ and the highest frequency $f_{\max}$ that the channel can transmit (usually given in the problem statement or in a specification sheet).
\item Ensure both frequencies are expressed in the \emph{same unit} (Hz, kHz, MHz, GHz).
\item Apply the formula: $B = f_{\max} - f_{\min}$.
\item State the result with the correct unit (Hz, kHz, MHz, GHz).
\end{enumerate}
\end{methode}

\begin{exoresolu}[Computing a bandwidth and comparing two channels]
\textbf{Statement.} Channel A carries frequencies from $1\ \mathrm{MHz}$ to $5\ \mathrm{MHz}$. Channel B carries frequencies from $88\ \mathrm{MHz}$ to $108\ \mathrm{MHz}$ (the FM radio band used by CRTV and private stations). (a) Compute the bandwidth of each channel. (b) Which channel can potentially carry data faster? Justify.
\tcblower
\textbf{Solution.}
(a) Bandwidth is the difference between the highest and the lowest frequency:
\begin{align*}
B_A &= f_{\max} - f_{\min} = 5\ \mathrm{MHz} - 1\ \mathrm{MHz} = 4\ \mathrm{MHz},\\
B_B &= 108\ \mathrm{MHz} - 88\ \mathrm{MHz} = 20\ \mathrm{MHz}.
\end{align*}
(b) Channel B has the larger bandwidth ($20\ \mathrm{MHz} > 4\ \mathrm{MHz}$). Since a greater bandwidth allows a higher potential data rate, channel B can potentially carry data faster than channel A. Note that the \emph{position} of the band on the frequency axis does not matter for this comparison; only its \emph{width} does.
\end{exoresolu}

\courssub{Baseband and Broadband Transmission}

Knowing that a channel has a certain bandwidth, a natural question arises: should we give the \emph{whole} bandwidth to \emph{one} signal, or \emph{divide} it so that \emph{several} signals travel at the same time? Both strategies exist, and they have names.

\begin{definition}[Baseband transmission]
In \cle{baseband transmission}, a \emph{single digital signal} uses the \emph{entire capacity} of the medium at any instant. The signal is sent directly, without being shifted to a carrier frequency. Baseband is simple and cheap, but only one communication can use the medium at a time, so devices must take turns. It is typical of \emph{Ethernet} local area networks (LANs), such as a school computer room network.
\end{definition}

\begin{definition}[Broadband transmission]
In \cle{broadband transmission}, the bandwidth of the medium is \emph{divided into several frequency channels}, so that \emph{multiple signals share the same medium simultaneously}. Each signal is carried on its own frequency band; this technique is called \emph{frequency division multiplexing (FDM)}. Examples: a cable network carrying television channels \emph{and} internet on the same coaxial cable, or the many radio stations (CRTV, private FM stations) sharing the airwaves at once.
\end{definition}

\begin{popfigure}
\begin{center}
\begin{tikzpicture}[font=\small]
% Baseband axis
\draw[->, thick, popmuted] (0,2.6) -- (10.4,2.6) node[right]{frequency};
\fill[popblue!55] (0.4,2.15) rectangle (9.6,3.05);
\node[popdark] at (5,2.6) {\textbf{One signal uses the whole band}};
\node[left, popdark, align=right] at (0,2.6) {\textbf{Baseband}};
% Broadband axis
\draw[->, thick, popmuted] (0,0.6) -- (10.4,0.6) node[right]{frequency};
\fill[poporange!60] (0.4,0.15) rectangle (3.2,1.05);
\fill[popgreen!55] (3.5,0.15) rectangle (6.3,1.05);
\fill[poppurple!50] (6.6,0.15) rectangle (9.6,1.05);
\node[popdark] at (1.8,0.6) {Channel 1};
\node[popdark] at (4.9,0.6) {Channel 2};
\node[popdark] at (8.1,0.6) {Channel 3};
\node[left, popdark, align=right] at (0,0.6) {\textbf{Broadband}};
\end{tikzpicture}
\end{center}
\legende{Baseband: one signal occupies the whole bandwidth. Broadband: the bandwidth is split into several frequency channels used simultaneously.}
\end{popfigure}

\begin{exemplebox}[Broadband in everyday Cameroonian life]
\begin{itemize}
\item \textbf{Cable TV + internet:} one coaxial cable entering a home can carry dozens of TV channels and an internet connection at the same time, each on its own frequency band.
\item \textbf{FM radio:} in Yaoundé you can tune to CRTV or to several private stations; all the signals cross the same air simultaneously, separated by frequency. Your radio simply selects one channel.
\item \textbf{Mobile networks:} MTN and Orange divide their licensed frequency bands into many channels so that thousands of subscribers can call and browse at the same time.
\end{itemize}
\end{exemplebox}

\begin{savaistu}[Why ``broadband'' means ``fast internet'']
In everyday language, ``broadband'' has come to mean any fast, always-on internet connection (CAMTEL fibre, 4G). This popular meaning comes from the technical one: broadband techniques exploit a \emph{wide} bandwidth, and a wide bandwidth allows a \emph{high data rate}. So the marketing word and the physics agree: broadband is fast because it is wide!
\end{savaistu}

\begin{propriete}[Narrowband vs. broadband (classification by bandwidth)]
In the context of \emph{transmission capacity} (Lesson 28), channels are also classified by the \emph{size} of their bandwidth:
\begin{itemize}
\item A \cle{narrowband} channel has a \textbf{small bandwidth} (typically a few kHz), sufficient for voice or low-speed data. Example: a traditional telephone line ($B \approx 3.1\ \mathrm{kHz}$).
\item A \cle{broadband} channel has a \textbf{large bandwidth} (typically MHz or GHz), allowing high-speed data, video, and multiple simultaneous services. Examples: coaxial cable, optical fibre, 4G/5G radio channels.
\end{itemize}
This classification by \emph{width} is distinct from the \emph{transmission method} (baseband vs. broadband) defined above. A baseband Ethernet cable can be broadband in capacity (e.g., 100 MHz), while a broadband FM radio channel is narrowband in capacity (each station $\approx 200\ \mathrm{kHz}$).
\end{propriete}

\begin{aretenir}[Key points of this section]
\begin{itemize}
\item \textbf{Data transmission} = sending data as signals from a sender to a receiver over a channel, following a protocol.
\item The five elements of a transmission system: \textbf{sender, receiver, signal, medium, protocol}.
\item \textbf{Analogue} signals vary continuously; \textbf{digital} signals take discrete values (0 and 1). Computers use digital signals because they resist noise and can be regenerated perfectly.
\item \textbf{Bandwidth} = range of frequencies a channel can carry, $B = f_{\max} - f_{\min}$, measured in \textbf{hertz}.
\item Greater bandwidth $\Rightarrow$ higher \emph{potential} \textbf{data rate} (in bps) — but the two quantities are not the same.
\item \textbf{Baseband}: one digital signal uses the whole medium (Ethernet LANs). \textbf{Broadband (method)}: several signals share the medium on different frequency channels simultaneously (FDM).
\item \textbf{Narrowband (capacity)}: small bandwidth, low data rate. \textbf{Broadband (capacity)}: large bandwidth, high data rate.
\end{itemize}
\end{aretenir}

\begin{exercice}[Checking your understanding]
\begin{enumerate}
\item Name the five elements of a data transmission system and identify each of them in the case of a student sending an email from a cybercafé in Bafoussam.
\item A channel carries frequencies from $2\ \mathrm{MHz}$ to $14\ \mathrm{MHz}$. Calculate its bandwidth.
\item State two advantages of digital signals over analogue signals.
\item A classmate writes: ``The bandwidth of my internet connection is 20 Mbps.'' Correct the statement and explain the error.
\item Classify each situation as baseband or broadband (transmission method): (a) two computers exchanging a file on an Ethernet cable in a school LAN; (b) a household receiving TV channels and internet through the same coaxial cable.
\item A traditional telephone line has a bandwidth of $3.1\ \mathrm{kHz}$. A fibre optic link has a bandwidth of $10\ \mathrm{GHz}$. Which is narrowband and which is broadband in terms of capacity? Justify.
\end{enumerate}
\end{exercice}

\begin{corrige}[Exercise 1]
\begin{enumerate}
\item Sender: the student's computer; receiver: the mail server (then the recipient's device); signal: electrical/radio pulses encoding the email; medium: the cybercafé's cable and the ISP's network; protocol: the rules for sending mail and transmitting data on the network.
\item $B = f_{\max} - f_{\min} = 14 - 2 = 12\ \mathrm{MHz}$.
\item Digital signals resist noise (a 0 or 1 can still be recognised despite distortion) and can be regenerated perfectly by repeaters, allowing long-distance transmission without loss of quality.
\item The unit is wrong: bandwidth is measured in hertz, not in bits per second. Correct statement: ``The \emph{data rate} of my connection is 20 Mbps'' (or give the bandwidth in MHz if that is what is meant).
\item (a) Baseband: one signal uses the whole cable at a time. (b) Broadband: several signals (TV channels, internet) share the cable on different frequency channels simultaneously.
\item Telephone line: narrowband (small bandwidth $\approx 3.1\ \mathrm{kHz}$). Fibre optic: broadband (very large bandwidth $\approx 10\ \mathrm{GHz}$). The fibre can carry vastly more data per second.
\end{enumerate}
\end{corrige}

Now that the notions of signal, bandwidth, baseband and broadband are clear, we are ready to study in detail the two great families of transmission classified by bandwidth — \emph{broadband and narrowband transmission} — in the next section.

\coursec{Broadband and Narrowband Transmission (Lesson 28)}

In the previous section we saw that every communication channel is characterised by its \cle{bandwidth} --- the range of frequencies it can carry --- and that bandwidth is the main factor limiting the \cle{data rate} of a transmission. We now use this idea to classify transmission systems into two great families: \cle{narrowband} and \cle{broadband}. This classification is a classic GCE Advanced Level question, and it is also the key to understanding why your mobile phone in Yaoundé can stream a video while a village radio set can only carry one voice at a time.

\courssub{Intuition: a road with one lane or many lanes}

Imagine the road between Douala and Bafoussam. If the road has a single narrow lane, vehicles must pass one at a time: the traffic (the data) moves slowly, and two lorries cannot travel side by side. If the road is a wide motorway with many lanes, several vehicles travel simultaneously and the total traffic per hour is enormous.

A transmission channel behaves in exactly the same way. The ``width of the road'' is the \textbf{bandwidth} (measured in hertz), and the ``vehicles'' are the \textbf{signals}. A channel with a narrow band of frequencies can carry only one signal at a time, at a modest data rate: this is \emph{narrowband}. A channel with a wide band of frequencies can be divided into many sub-channels, each carrying its own signal at the same time: this is \emph{broadband}.

\courssub{Narrowband Transmission}

\begin{definition}[Narrowband transmission]
\cle{Narrowband transmission} is data communication that uses a \textbf{narrow range of frequencies} (a small bandwidth). The channel carries \textbf{one signal at a time}, and the achievable data rate is relatively \textbf{low}.
\end{definition}

Because only one signal occupies the channel at any instant, narrowband systems are simple: there is no need for complex equipment to separate many simultaneous signals. This explains their main characteristics.

\begin{aretenir}[Characteristics of narrowband transmission]
\begin{itemize}
\item \textbf{Limited bandwidth}: the channel occupies a small slice of the frequency spectrum.
\item \textbf{Lower speed}: data rates are modest (from a few bits per second up to about $56\ \mathrm{kbps}$ for a dial-up modem).
\item \textbf{One signal at a time}: no simultaneous voice, video and data on the same channel.
\item \textbf{Cheaper equipment}: transmitters and receivers are simple and inexpensive.
\item \textbf{Longer possible distances}: for some technologies (e.g.\ low-frequency radio), narrowband signals travel far and penetrate obstacles well, which is why they are still used for rural radio communication.
\end{itemize}
\end{aretenir}

\begin{exemplebox}[Everyday examples of narrowband]
\begin{itemize}
\item \textbf{Traditional dial-up telephone modem}: uses the voice band of an ordinary telephone line and reaches at most $56\ \mathrm{kbps}$. While connected, the telephone cannot be used for a voice call --- one signal at a time!
\item \textbf{Some radio communication}: walkie-talkies used by security guards, or a simple FM voice channel, occupy a narrow band and carry a single conversation.
\item \textbf{Early mobile networks (2G)}: designed essentially for voice calls and SMS text messages at low data rates.
\end{itemize}
\end{exemplebox}

\courssub{Broadband Transmission}

\begin{definition}[Broadband transmission]
\cle{Broadband transmission} is \textbf{high-bandwidth} data communication in which a \textbf{wide range of frequencies} is used, so that \textbf{multiple signals can be carried simultaneously} over the same medium, at \textbf{high data rates}.
\end{definition}

How can several signals share one cable without mixing? The wide frequency band is divided into many \textbf{sub-channels}, each centred on a different carrier frequency. Voice, television and Internet data are each placed on different sub-channels, exactly like different radio stations broadcasting at the same time on different frequencies. This technique has a name you should know.

\begin{propriete}[Frequency division multiplexing]
Broadband transmission allows \textbf{several signals to share one physical medium at the same time} by assigning each signal a different frequency sub-channel. This technique is called \cle{frequency division multiplexing} (FDM). A device at the receiving end filters the sub-channels to recover each original signal. (The detailed working of FDM is not examinable; knowing its purpose is sufficient.)
\end{propriete}

\begin{aretenir}[Characteristics of broadband transmission]
\begin{itemize}
\item \textbf{Wide frequency range}: a large bandwidth is available.
\item \textbf{High speed}: from a few megabits per second to gigabits per second.
\item \textbf{Many simultaneous signals}: voice, video and data travel together on one line.
\item \textbf{Always-on connection}: no need to ``dial''; the connection is permanently available.
\item \textbf{Supports multimedia}: video calls, streaming, online learning and large downloads are possible.
\end{itemize}
\end{aretenir}

\begin{exemplebox}[Broadband technologies]
\begin{itemize}
\item \textbf{ADSL}: broadband over the ordinary copper telephone line; the voice call and the Internet data use different frequency bands, so both work at the same time. Typical speeds: $2$ to $20\ \mathrm{Mbps}$.
\item \textbf{Fibre optic (FTTH, Fibre To The Home)}: light signals in a glass fibre; typical speeds from $100\ \mathrm{Mbps}$ up to $1\ \mathrm{Gbps}$ and beyond.
\item \textbf{Cable Internet}: broadband delivered over the coaxial cable of a television network.
\item \textbf{Mobile broadband (3G, 4G, 5G)}: high-speed Internet through the mobile phone network; 4G typically offers $10$ to $100\ \mathrm{Mbps}$.
\end{itemize}
\end{exemplebox}

\courssub{Comparing the Two Families}

The following diagram shows the essential difference: a narrowband channel is a single thin frequency band carrying one signal, whereas a broadband channel is a wide band divided into many sub-channels used in parallel.

\begin{popfigure}
\begin{center}
\begin{tikzpicture}[scale=1.0]
% Narrowband
\draw[->,thick,popink] (0,0) -- (10.5,0) node[below left,font=\small]{frequency};
\draw[fill=popblueL,draw=popblue,thick] (4.2,0) rectangle (5.4,1.1);
\node[font=\small,align=center] at (4.8,0.55) {one\\signal};
\draw[decorate,decoration={brace,amplitude=4pt},popblue] (4.2,1.25) -- (5.4,1.25) node[midway,above=5pt,font=\small,text=popblue]{narrow band};
\node[font=\small\bfseries,text=popink] at (5.2,-0.7) {Narrowband: one signal at a time};
% Broadband
\begin{scope}[yshift=-3.4cm]
\draw[->,thick,popink] (0,0) -- (10.5,0) node[below left,font=\small]{frequency};
\foreach \x/\c/\t in {1.2/popblueL/voice, 2.9/popgreenL/data, 4.6/popblueL/video, 6.3/popgreenL/data, 8.0/popblueL/TV}{
\draw[fill=\c,draw=popblue,thick] (\x,0) rectangle (\x+1.3,1.1);
\node[font=\small] at (\x+0.65,0.55) {\t};}
\draw[decorate,decoration={brace,amplitude=4pt},poporange] (1.2,1.25) -- (9.3,1.25) node[midway,above=5pt,font=\small,text=poporange]{wide band: many sub-channels (FDM)};
\node[font=\small\bfseries,text=popink] at (5.2,-0.7) {Broadband: many signals simultaneously};
\end{scope}
\end{tikzpicture}
\end{center}
\legende{Narrowband channel (single thin band, one signal) versus broadband channel (wide band split into sub-channels by frequency division multiplexing).}
\end{popfigure}

The speed gap between the technologies is so large that a logarithmic scale is needed to display it on one chart:

\begin{popfigure}
\begin{center}
\begin{tikzpicture}
\begin{axis}[
    ybar, ymode=log, log origin=infty,
    width=12cm, height=6.5cm,
    bar width=0.9cm,
    ymin=0.01, ymax=3000,
    ylabel={typical speed (Mbps, log scale)},
    symbolic x coords={Dial-up, ADSL, 4G, Fibre},
    xtick=data,
    nodes near coords, every node near coord/.append style={font=\small},
    ymajorgrids=true, grid style={popline},
    tick label style={font=\small},
    label style={font=\small},
    fill=popblueL, draw=popblue
]
\addplot coordinates {(Dial-up,0.056) (ADSL,10) (4G,50) (Fibre,500)};
\end{axis}
\end{tikzpicture}
\end{center}
\legende{Typical speeds on a logarithmic scale: dial-up $56\ \mathrm{kbps}$, ADSL $2$--$20\ \mathrm{Mbps}$, 4G $10$--$100\ \mathrm{Mbps}$, fibre $100\ \mathrm{Mbps}$--$1\ \mathrm{Gbps}$ (bar heights show representative values).}
\end{popfigure}

\begin{aretenir}[Comparison table: narrowband vs broadband]
\begin{center}
\begin{tabularx}{\linewidth}{|l|X|X|}
\hline
\rowcolor{popblueL} \textbf{Criterion} & \textbf{Narrowband} & \textbf{Broadband} \\
\hline
Bandwidth & Narrow (small frequency range) & Wide (large frequency range) \\
\hline
Speed & Low (up to about $56\ \mathrm{kbps}$ for dial-up) & High (Mbps to Gbps) \\
\hline
Simultaneous signals & One signal at a time & Many signals at once (FDM) \\
\hline
Connection & Often dial-up (on demand) & Always on \\
\hline
Equipment cost & Cheap, simple & More expensive, more complex \\
\hline
Typical uses & Voice calls, SMS, simple radio, telemetry & Internet, video streaming, video calls, multimedia \\
\hline
Examples & Dial-up modem, 2G voice, walkie-talkie & ADSL, fibre (FTTH), cable, 3G/4G/5G \\
\hline
\end{tabularx}
\end{center}
\end{aretenir}

\courssub{Choosing Between Narrowband and Broadband}

\begin{methode}[Choosing the right transmission type for an application]
\begin{enumerate}
\item \textbf{Identify the needed speed.} Estimate the data rate the application requires (an SMS needs almost nothing; a video call needs several $\mathrm{Mbps}$).
\item \textbf{Count the simultaneous services.} Will voice, video and data be used together? If yes, broadband is required.
\item \textbf{Check availability.} Is broadband infrastructure present at the location (fibre, 4G coverage)? In some rural areas only narrowband radio or 2G is available.
\item \textbf{Consider the budget.} Narrowband equipment and subscriptions are cheaper; broadband costs more but delivers far more.
\item \textbf{Decide and justify.} Choose the cheapest solution that satisfies steps 1--3, and state the reason clearly in an examination answer.
\end{enumerate}
\end{methode}

\begin{exoresolu}[Choosing a transmission type]
A health post in a remote village of the North-West Region must (a) send short daily text reports of a few kilobytes to the district hospital, and (b) allow a doctor in Bamenda to hold a monthly video consultation. The village has 2G coverage only; the nearest 4G coverage is $30\ \mathrm{km}$ away. Propose a transmission type for each need, justifying your answer.
\tcblower
\textbf{Solution.}
\begin{itemize}
\item \textbf{(a) Daily text reports.} A few kilobytes per day is a tiny amount of data: even $56\ \mathrm{kbps}$ transfers $5\ \mathrm{KB}$ in less than one second. \textbf{Narrowband} (2G SMS or a low-rate data connection) is sufficient, cheap, and it is the only technology available on site. Need: low speed --- satisfied; budget: minimal; availability: yes.
\item \textbf{(b) Monthly video consultation.} A video call requires a continuous flow of voice \emph{and} video, i.e.\ several megabits per second and two simultaneous signals: this is impossible on narrowband. \textbf{Broadband} is required. Since the village has no 4G or fibre, a practical solution is to travel to the 4G-covered town for the consultation, or to install a broadband link (e.g.\ a satellite or fixed-wireless broadband connection) if the budget allows.
\end{itemize}
\textbf{Conclusion:} narrowband for the text reports (sufficient, available, cheap); broadband for the video consultation (the only option that supports simultaneous voice and video at high speed).
\end{exoresolu}

\begin{attention}[Common mistake: ``broadband means wireless'']
Many candidates write that broadband is Internet ``without cables''. This is \textbf{wrong}. \textbf{Broadband refers to the bandwidth (wide frequency range, high speed), not to the transmission medium.} ADSL and fibre optic are broadband technologies that use \emph{cables}; conversely, a 2G mobile link is \emph{wireless} but \emph{narrowband}. In the examination, always define broadband by its wide bandwidth and multiple simultaneous signals --- never by the absence of wires.
\end{attention}

\begin{savaistu}[Broadband in Cameroon (enrichment, not examinable in detail)]
Cameroon has invested in a national \textbf{fibre-optic backbone} operated mainly by \textbf{CAMTEL}, linking major towns and landing stations of submarine cables on the Atlantic coast at Douala and Kribi. In \textbf{Yaoundé and Douala}, 4G mobile broadband from operators such as MTN and Orange is widely available, and fibre-to-the-home offers are growing. However, \textbf{rural connectivity remains a challenge}: long distances, the cost of extending fibre, and low population density mean that many villages still depend on 2G narrowband or satellite links. Projects supported by the state and by \textbf{ANTIC} aim to extend broadband access to schools and public services across the country.
\end{savaistu}

\begin{exercice}[Narrowband or broadband?]
For each situation, state whether narrowband or broadband is appropriate and justify in one sentence:
\begin{enumerate}
\item A weather station on Mount Cameroon sends temperature readings of a few bytes every hour by radio.
\item A cybercafé in Buea offers twenty customers simultaneous web browsing and video streaming.
\item A family in Douala wants one telephone line used for both voice calls and Internet at the same time.
\end{enumerate}
\end{exercice}

\begin{corrige}[Exercise 1]
\begin{enumerate}
\item \textbf{Narrowband}: the data volume is tiny and infrequent, and simple radio equipment is cheap and works over long distances.
\item \textbf{Broadband}: many users and multimedia traffic require high speed and many simultaneous signals.
\item \textbf{Broadband (ADSL)}: frequency division multiplexing places the voice call and the Internet data on different sub-channels of the same line, so both work simultaneously.
\end{enumerate}
\end{corrige}

\begin{aretenir}[Key points of Lesson 28]
\begin{itemize}
\item \textbf{Narrowband}: narrow frequency range, one signal at a time, low speed, cheap equipment (dial-up $56\ \mathrm{kbps}$, 2G, simple radio).
\item \textbf{Broadband}: wide frequency range, many simultaneous signals via \textbf{frequency division multiplexing}, high speed, always on (ADSL, fibre, cable, 3G/4G/5G).
\item Broadband describes the \textbf{bandwidth}, not the medium: broadband can be wired or wireless.
\item Choose by analysing: needed speed, number of simultaneous services, availability, budget.
\end{itemize}
\end{aretenir}

\coursec{Types of Data Transmission (Lesson 29)}

In the previous section we studied \emph{how much} data a channel can carry: we distinguished \cle{narrowband} channels, which offer a limited bandwidth, from \cle{broadband} channels, which carry large volumes of data at high speed, often by dividing the medium into several frequency channels. But bandwidth alone does not tell the whole story. Two links with the same bandwidth can behave very differently depending on \emph{how} the bits are organised and \emph{in which direction} they are allowed to flow. Does the sender push bits one after another on a single wire, or eight at a time on eight wires? Can both devices talk at once, or must they take turns? Are the bits framed by start and stop markers, or clocked in large blocks? Answering these questions is the object of this lesson: the \cle{types of data transmission}.

\courssub{Serial and Parallel Transmission}

\textbf{An everyday observation.}\par
Imagine 800 students who must leave the Lyc\'ee de Buea through the school gate. Two strategies are possible: open \emph{one} gate and let the students pass one after another, or open \emph{eight} gates side by side and let eight students pass simultaneously. The first strategy is slower per gate but needs only one gate; the second is faster in principle but requires eight gates, eight doorkeepers, and the students must stay aligned --- if one line moves faster than the others, the group arrives disorganised. Data transmission faces exactly the same choice.

\begin{definition}[Serial and parallel transmission]
In \cle{serial transmission}, the bits of a message are sent \textbf{one after another} over a \textbf{single} communication channel (one wire or one pair of wires for differential signalling).

In \cle{parallel transmission}, several bits (typically 8, 16, 32 or 64) are sent \textbf{simultaneously}, each on its \textbf{own wire}, so that a whole byte or word crosses the link in a single clock step.
\end{definition}

At first sight, parallel transmission seems obviously superior: with 8 wires it should be 8 times faster. In practice, three problems appear as speed and distance increase:

\begin{itemize}
\item \cle{Skew}: the wires of a parallel cable never have exactly the same length or electrical properties, so bits sent at the same instant arrive at slightly different times. At high frequency the receiver can no longer tell which bits belong together.
\item \cle{Crosstalk}: many wires side by side interfere electromagnetically with one another, corrupting the signals.
\item \cle{Cost and bulk}: a parallel cable needs many conductors, connectors and pins --- expensive and fragile over long distances.
\end{itemize}

Serial transmission avoids skew (there is only one line) and is cheap and reliable over long distances. This is why \textbf{modern systems prefer high-speed serial links}: engineers learned to make a single serial line switch so fast (billions of bits per second) that it outperforms any affordable parallel bus. The evidence is all around us: \cle{USB} (Universal \emph{Serial} Bus) replaced the old parallel printer port, and \cle{SATA} (Serial ATA) replaced the wide 40-wire \cle{PATA} ribbon cables that connected hard disks inside computers. Network cables (Ethernet) are also serial.

\begin{popfigure}
\begin{tikzpicture}[scale=0.9, every node/.style={font=\small}]
% Serial
\node[font=\small\bfseries, popblue] at (1.2,2.6) {Serial (one line)};
\draw[thick, popink] (0,2.0) -- (7.5,2.0);
\foreach \x/\b in {0.2/1, 1.0/0, 1.8/1, 2.6/1, 3.4/0, 4.2/0, 5.0/1, 5.8/0}{
  \draw[popblue, thick] (\x,1.8) rectangle (\x+0.8,2.2);
  \node at (\x+0.4,2.0) {\b};}
\node[font=\footnotesize, popmuted] at (3.9,1.5) {bits follow one another in time};
% Parallel
\node[font=\small\bfseries, poporange] at (1.2,0.6) {Parallel (8 lines)};
\foreach \y/\b in {0.0/1, -0.5/0, -1.0/1, -1.5/1, -2.0/0, -2.5/0, -3.0/1, -3.5/0}{
  \draw[thick, popink] (0,\y) -- (7.5,\y);
  \draw[poporange, thick] (0.2,\y-0.18) rectangle (1.0,\y+0.18);
  \node at (0.6,\y) {\b};}
\draw[popdark, dashed, thick] (0.1,0.35) -- (0.1,-3.85);
\draw[popdark, dashed, thick] (1.1,0.35) -- (1.1,-3.85);
\node[font=\footnotesize, popmuted] at (3.9,-4.2) {all 8 bits cross during the same time slot};
\end{tikzpicture}
\legende{Timing comparison: in serial transmission the 8 bits of a byte travel one after another on a single line; in parallel transmission all 8 bits travel at once on 8 separate lines.}
\end{popfigure}

\begin{center}
\begin{tabularx}{\linewidth}{|l|X|X|}
\hline
\rowcolor{popblueL} \textbf{Criterion} & \textbf{Serial} & \textbf{Parallel} \\
\hline
Number of wires & 1 (or 1 differential pair) & Several (8, 16, 32, 64) \\
\hline
Speed per wire & Lower, but very high clock rates achievable (GHz) & Higher in theory, limited by skew and crosstalk \\
\hline
Reliable distance & Long (metres to kilometres) & Short (centimetres, inside a device) \\
\hline
Cost & Low (cable, connectors, PCB traces) & High (cable, connectors, PCB area) \\
\hline
Typical uses & USB, SATA, PCIe, Ethernet, telephone lines, fibre & Internal CPU buses (DDR RAM), old printer ports (LPT), PATA \\
\hline
\end{tabularx}
\end{center}

\begin{attention}[Parallel is not obsolete everywhere]
Parallel transmission survives \emph{inside} devices where distances are centimetres and clock rates are controlled: the memory bus (DDR4/DDR5) uses 64 data lines in parallel, and GPU memory buses use even more. The rule is: \textbf{parallel for short, high-bandwidth internal paths; serial for everything that leaves the chip or the box}.
\end{attention}

\courssub{Transmission Modes: Simplex, Half-Duplex and Full-Duplex}

Serial or parallel describes \emph{how many bits travel together}. A second, independent question is: \emph{in which direction(s)} may data flow between two devices $A$ and $B$?

\begin{definition}[Direction of transmission modes]
\begin{itemize}
\item \cle{Simplex}: data flows in \textbf{one direction only}; one device is always the sender, the other always the receiver. Examples: radio and television broadcasting (CRTV), a keyboard sending characters to the computer, a temperature sensor transmitting to a logger.
\item \cle{Half-duplex}: data flows in \textbf{both directions, but one at a time}; the two devices must take turns. Examples: walkie-talkies used by security guards (press to talk, release to listen), traditional Ethernet hubs (CSMA/CD), USB 2.0.
\item \cle{Full-duplex}: data flows in \textbf{both directions simultaneously}. This requires two independent channels (two wire pairs in copper, two fibres, or two frequency bands). Examples: an ordinary telephone call (both people can speak and hear at the same time), modern switched Ethernet networks (1000BASE-T and above), USB 3.0/3.1/4, fibre optic links.
\end{itemize}
\end{definition}

\begin{popfigure}
\begin{tikzpicture}[scale=1.0, every node/.style={font=\small}]
% Simplex
\node[draw, rounded corners, fill=popblueL, minimum width=1.1cm, minimum height=0.7cm] (A1) at (0,0) {$A$};
\node[draw, rounded corners, fill=popblueL, minimum width=1.1cm, minimum height=0.7cm] (B1) at (4,0) {$B$};
\draw[-{Stealth}, very thick, popblue] (A1) -- (B1);
\node[font=\small\bfseries] at (2,-0.8) {Simplex};
% Half-duplex
\node[draw, rounded corners, fill=popgreenL, minimum width=1.1cm, minimum height=0.7cm] (A2) at (7,0) {$A$};
\node[draw, rounded corners, fill=popgreenL, minimum width=1.1cm, minimum height=0.7cm] (B2) at (11,0) {$B$};
\draw[-{Stealth}, very thick, popgreen] (A2.north east) to[bend left=20] (B2.north west);
\draw[-{Stealth}, very thick, popgreen, dashed] (B2.south west) to[bend left=20] (A2.south east);
\node[font=\small\bfseries] at (9,-0.8) {Half-duplex (turn by turn)};
% Full-duplex
\node[draw, rounded corners, fill=popblueL, minimum width=1.1cm, minimum height=0.7cm] (A3) at (2,-2.6) {$A$};
\node[draw, rounded corners, fill=popblueL, minimum width=1.1cm, minimum height=0.7cm] (B3) at (6,-2.6) {$B$};
\draw[-{Stealth}, very thick, poporange] (A3.north east) to[bend left=20] (B3.north west);
\draw[-{Stealth}, very thick, poporange] (B3.south west) to[bend left=20] (A3.south east);
\node[font=\small\bfseries] at (4,-3.5) {Full-duplex (simultaneous)};
\end{tikzpicture}
\legende{The three transmission modes. In half-duplex the two directions exist but are never used at the same instant (dashed arrow = waiting); in full-duplex both directions are active simultaneously.}
\end{popfigure}

\begin{methode}[Identifying the transmission mode in a scenario]
When a GCE question describes a situation and asks for the transmission mode, proceed in two steps:
\begin{enumerate}
\item \textbf{Direction}: can data travel in both directions at all? If only one device can ever send, the mode is \emph{simplex}.
\item \textbf{Simultaneity}: if both can send, can they send \emph{at the same time}? If they must take turns, the mode is \emph{half-duplex}; if they can talk and listen simultaneously, it is \emph{full-duplex}.
\end{enumerate}
Always justify your answer by quoting the clue in the scenario (``the guard must release the button to hear the reply'' $\Rightarrow$ half-duplex).
\end{methode}

\begin{attention}[Half-duplex is NOT simplex]
The most frequent mistake at the GCE is to classify a walkie-talkie as simplex ``because only one person speaks at a time''. Wrong: in simplex, the second device can \emph{never} send (a TV set cannot talk back to the CRTV transmitter). In half-duplex, \textbf{both} devices can send --- just not simultaneously. Two-way but not simultaneous = half-duplex.
\end{attention}

\courssub{Asynchronous and Synchronous Transmission}

A third classification concerns \emph{timing}: how does the receiver know when each bit (and each character) begins and ends?

\textbf{Intuition.}\par
Suppose a friend taps messages to you through a wall. If the taps can come at any moment, you need a clear signal for ``a letter starts now'' and ``the letter is finished'' --- otherwise you cannot separate the letters. That is the asynchronous solution. Alternatively, you could agree in advance: ``I will tap exactly one tap per second, continuously, in long regular sequences'' --- then no per-letter markers are needed, but both of you must keep perfect time. That is the synchronous solution.

\begin{definition}[Asynchronous transmission]
In \cle{asynchronous transmission}, data is sent \textbf{character by character} (typically 7 or 8 bits), and each character is framed by a \cle{start bit} (which warns the receiver that a character is beginning) and one or two \cle{stop bits} (which mark its end and return the line to the idle state). The line is idle (mark state, logic 1) between characters, and no common clock is shared between sender and receiver. The receiver re-synchronises on every start bit.
\end{definition}

\begin{popfigure}
\begin{tikzpicture}[scale=0.85, every node/.style={font=\small}]
% idle
\draw[very thick, popmuted] (0,1) -- (1,1);
\node[font=\footnotesize, popmuted] at (0.5,1.35) {idle (mark)};
% start bit
\draw[very thick, popink] (1,1) -- (1,0) -- (2,0);
\node[font=\footnotesize] at (1.5,-0.4) {start (0)};
% data bits (7 bits example)
\draw[very thick, popink] (2,0) -- (2,1) -- (3,1) -- (3,0) -- (4,0) -- (4,1) -- (5,1) -- (5,0) -- (6,0) -- (6,1) -- (7,1) -- (7,0) -- (8,0) -- (8,1) -- (9,1);
\foreach \x/\b in {2.5/1, 3.5/0, 4.5/1, 5.5/0, 6.5/0, 7.5/1, 8.5/1}{
  \node[font=\footnotesize, popblue] at (\x,0.55) {\b};}
\draw[popblue, dashed] (2,1.5) -- (9,1.5);
\node[font=\footnotesize, popblue] at (5.5,1.8) {7 data bits (LSB first)};
% stop bit(s)
\draw[very thick, popink] (9,1) -- (9,0) -- (10,0) -- (10,1) -- (11,1);
\node[font=\footnotesize] at (9.5,-0.4) {stop (1)};
\node[font=\footnotesize] at (10.5,-0.4) {stop (1)};
\draw[very thick, popmuted] (11,1) -- (12,1);
\node[font=\footnotesize, popmuted] at (11.5,1.35) {idle};
\end{tikzpicture}
\legende{Asynchronous frame (example: 7 data bits, 1 start bit, 2 stop bits). The line is held high (mark) when idle. The start bit (space, logic 0) triggers the receiver's clock. Stop bits return the line to mark.}
\end{popfigure}

The advantage is \textbf{simplicity}: cheap hardware, no shared clock, and characters can arrive at irregular intervals (ideal for a keyboard, where the user types unpredictably). The price is \cle{overhead}: for a 7-bit character with 1 start and 2 stop bits, 3 bits out of 10 --- $30\%$ of the line time --- carry no useful data. For 8 data bits, 1 start, 1 stop: 2 bits out of 10 = $20\%$ overhead.

\begin{definition}[Synchronous transmission]
In \cle{synchronous transmission}, data is sent in \textbf{blocks (frames)} of many bytes (hundreds to thousands), preceded by special \cle{synchronisation characters} (flags or preambles) and followed by control information (address, control, error-checking code such as CRC). The sender's and receiver's \textbf{clocks are synchronised} (via a separate clock line, or a clock embedded in the signal using line coding like Manchester or NRZI), so no start/stop bits are needed per character. The receiver samples the line at precise clock intervals.
\end{definition}

Synchronous transmission is \textbf{faster and more efficient for large volumes of data} because the overhead (a few control bytes per block of hundreds or thousands of bytes) is proportionally tiny (often $<5\%$). Its drawbacks are \textbf{complexity and cost}: precise clock synchronisation requires more sophisticated equipment (PLLs, clock recovery circuits). It is used in high-speed network links (Ethernet, HDLC, PPP, Fibre Channel), while asynchronous transmission survives in simple, low-speed serial interfaces (RS-232, USB 1.x/2.0 low/full speed control endpoints).

\begin{center}
\begin{tabularx}{\linewidth}{|l|X|X|}
\hline
\rowcolor{popblueL} \textbf{Criterion} & \textbf{Asynchronous} & \textbf{Synchronous} \\
\hline
Unit sent & One character at a time (5--9 bits) & Blocks / frames of many bytes (hundreds--thousands) \\
\hline
Synchronisation & Start and stop bits per character & Shared clock + sync flags / preamble per frame \\
\hline
Overhead & High (20--30\% per character) & Low for large blocks ($<5\%$) \\
\hline
Speed & Low to moderate (typically $<115.2$ kbps for RS-232) & High (Mbps to Gbps) \\
\hline
Cost / complexity & Simple, cheap UART hardware & Complex, costly (clock recovery, buffers) \\
\hline
Best for & Irregular, small data (keyboard, mouse, GPS, old modems) & Large volumes, continuous streams (LAN, WAN, storage) \\
\hline
\end{tabularx}
\end{center}

\begin{exemplebox}[Overhead calculation]
A device sends ASCII characters (7 bits) asynchronously with 1 start bit and 2 stop bits. What percentage of the transmitted bits are useful data?
\textbf{Solution:} Total bits per character = 1 (start) + 7 (data) + 2 (stop) = 10 bits. Useful bits = 7. Overhead = 3 bits. Percentage useful = $\frac{7}{10}\times100\% = 70\%$. Overhead = $30\%$.
If the same data were sent in synchronous frames of 1000 bytes (8000 bits) with 48 bits of overhead (flag, address, control, CRC), overhead = $\frac{48}{8048}\times100\% \approx 0.6\%$.
\end{exemplebox}

\courssub{Choosing a Transmission Type: Selection Criteria}

In an examination or a real network design, the choice among these types is guided by four criteria:

\begin{itemize}
\item \cle{Speed}: full-duplex synchronous serial links achieve the highest throughputs; asynchronous links are the slowest.
\item \cle{Cost}: serial and asynchronous solutions need fewer wires and simpler hardware, hence are cheaper.
\item \cle{Distance}: serial transmission resists skew over long distances; parallel is confined to short links inside equipment.
\item \cle{Volume of data}: large, regular data flows justify synchronous framing; small, irregular flows favour asynchronous transmission.
\end{itemize}

\begin{methode}[Selecting the appropriate transmission type for a given scenario]
\begin{enumerate}
\item \textbf{Analyse the data pattern}: Is it a continuous high-volume stream (video, bulk file transfer) $\rightarrow$ synchronous? Or sporadic small messages (keystrokes, sensor readings) $\rightarrow$ asynchronous?
\item \textbf{Determine directionality}: One-way only $\rightarrow$ simplex. Two-way but alternating $\rightarrow$ half-duplex. Two-way simultaneous $\rightarrow$ full-duplex.
\item \textbf{Assess distance and speed}: Long distance or high speed $\rightarrow$ serial. Very short, ultra-high bandwidth inside a box $\rightarrow$ parallel.
\item \textbf{Consider cost constraints}: Low budget $\rightarrow$ asynchronous serial. Performance critical $\rightarrow$ synchronous serial (or parallel if distance permits).
\end{enumerate}
\end{methode}

\begin{exoresolu}[Identifying transmission types in a Cameroonian scenario]
A cybercaf\'e in Bamenda uses: (a) a radio receiver tuned to CRTV FM; (b) walkie-talkies for its two security guards; (c) a telephone line for voice calls; (d) a USB 3.0 cable between a computer and an external SSD; (e) a keyboard connected to a PC, which sends each typed character framed by a start and a stop bit. For each case, give the transmission type(s) involved and justify.

\tcblower
\textbf{Solution.}
\begin{itemize}
\item (a) \textbf{Simplex}: the radio station transmits, the receiver can never send back. One direction only. Also \textbf{asynchronous} in the sense that the receiver does not share a clock with the transmitter (but broadcast uses its own framing).
\item (b) \textbf{Half-duplex}: both guards can talk, but each must press the button and release it to hear the other --- two directions, one at a time. Typically \textbf{serial} (single frequency).
\item (c) \textbf{Full-duplex}: both speakers can talk and hear simultaneously during the call. The telephone network uses \textbf{synchronous} transmission internally (TDM/PCM), but the analogue subscriber line appears as a full-duplex analogue channel.
\item (d) \textbf{Serial} transmission (Universal \emph{Serial} Bus): bits travel one after another over differential pairs. USB 3.0 uses \textbf{full-duplex} (separate TX and RX pairs) and \textbf{synchronous} transmission (clock embedded in 8b/10b or 128b/132b encoding).
\item (e) \textbf{Asynchronous} transmission: data is sent character by character with start and stop bits and no shared clock. The link keyboard $\to$ PC is also \textbf{simplex} (keyboard only sends) and \textbf{serial}.
\end{itemize}
Method used: for direction, ask ``can both devices send?'' then ``can they send at the same time?''; for timing, look for start/stop bits (asynchronous) or clocked frames (synchronous); for wiring, count the data paths.
\end{exoresolu}

\begin{attention}[Do not mix the three classifications]
Serial/parallel, simplex/half-duplex/full-duplex and asynchronous/synchronous are \textbf{three independent classifications}. A single link combines one choice from each family: for example, a USB 3.0 link is \emph{serial}, \emph{full-duplex}, and \emph{synchronous}. An old RS-232 modem link is \emph{serial}, \emph{full-duplex} (on two wires), and \emph{asynchronous}. Writing ``parallel full-duplex'' as if it were a single type is a classic error.
\end{attention}

\begin{savaistu}[Exam tip: memorise two examples per mode]
GCE Advanced Level questions very often ask: ``Give \emph{one} real-life example of each transmission mode.'' Prepare at least two per mode: \textbf{simplex} --- TV/radio broadcasting (CRTV), keyboard to computer, temperature sensor to data logger; \textbf{half-duplex} --- walkie-talkies, traditional Ethernet hub (CSMA/CD), USB 2.0; \textbf{full-duplex} --- telephone conversation (landline or mobile), modern switched Ethernet (1 Gbps+), USB 3.0+, fibre optic link. A precise, correctly justified example scores full marks; a vague one scores nothing.
\end{savaistu}

\begin{savaistu}[Cameroon context: Mobile Money and Network Links]
When you send money via MTN MoMo or Orange Money, your phone communicates with the operator's server over the cellular network (GSM/3G/4G). The air interface uses \textbf{serial} transmission (one frequency channel at a time per user), \textbf{full-duplex} (FDD: separate uplink/downlink frequencies, or TDD: alternating time slots), and \textbf{synchronous} framing (TDMA/OFDMA frames with precise timing). The core network uses synchronous optical fibre (SDH/OTN) and Ethernet. Understanding these types explains why you can talk and browse simultaneously (full-duplex) and why the network can handle millions of users (synchronous time/frequency division).
\end{savaistu}

\begin{exercice}[Practice]
\begin{enumerate}
\item A bank in Douala transfers its daily records (several gigabytes) every night to its head office over a dedicated fibre line. Would you recommend asynchronous or synchronous transmission? Justify with two criteria.
\item Explain why the wide PATA ribbon cable inside old computers was abandoned in favour of SATA, even though PATA sent 16 bits in parallel.
\item Classify each link as simplex, half-duplex or full-duplex: (i) a baby monitor with a ``talk-back'' button; (ii) a sensor sending temperature readings to a display; (iii) two computers on a modern Gigabit switch exchanging files; (iv) a TV remote control.
\item A device transmits 8-bit characters asynchronously with 1 start bit and 1 stop bit at 9600 baud. Calculate the effective data rate in useful bits per second.
\item Why does synchronous transmission require larger buffers at the receiver than asynchronous transmission?
\end{enumerate}
\end{exercice}

\begin{corrige}[Exercise n°1]
\textbf{1.} Synchronous transmission: the \emph{volume} is very large and the transfer is \emph{regular/scheduled}, so the low per-frame overhead gives much better efficiency; the higher equipment cost is justified by the speed gained. Also, synchronous allows continuous clock recovery, essential for high-speed fibre.

\textbf{2.} At high clock rates, the 16 parallel wires suffered from \emph{skew} (bits arriving at different times due to slight length/impedance differences) and \emph{crosstalk} (electromagnetic coupling between adjacent wires), and the cable was bulky, fragile and costly. A single very high-speed serial line (SATA) with embedded clock is faster in practice, cheaper, more reliable and allows thinner, longer cables.

\textbf{3.} (i) Baby monitor with talk-back: \emph{half-duplex} (both directions, one at a time, push-to-talk). (ii) Temperature sensor: \emph{simplex} (sensor only sends, display only receives). (iii) Computers on a Gigabit switch: \emph{full-duplex} (simultaneous exchange in both directions on separate wire pairs). (iv) TV remote control: \emph{simplex} (remote only sends IR pulses, TV only receives).

\textbf{4.} Total bits per character = 1 start + 8 data + 1 stop = 10 bits. Baud rate = 9600 symbols/s = 9600 bits/s (since 1 bit per symbol in baseband). Useful bits per second = $\frac{8}{10} \times 9600 = 7680\ \mathrm{bps}$.

\textbf{5.} Synchronous transmission sends large frames (hundreds to thousands of bytes) continuously. The receiver must store the entire frame (or a large portion) to perform error checking (CRC) before delivering data to the upper layer. Asynchronous transmission delivers one character at a time (after start/stop detection), so only a single-character buffer is needed.
\end{corrige}

\begin{aretenir}[Key points of Lesson 29]
\begin{itemize}
\item \textbf{Serial}: bits one after another on one channel --- cheap, reliable over long distances (USB, SATA, PCIe, Ethernet, fibre). \textbf{Parallel}: several bits at once on several wires --- fast only over very short distances, suffers from skew, crosstalk and cost (internal buses, DDR RAM).
\item \textbf{Simplex}: one direction only (radio, keyboard, sensor). \textbf{Half-duplex}: both directions, one at a time (walkie-talkie, USB 2.0, old Ethernet hub). \textbf{Full-duplex}: both directions simultaneously (telephone, switched Ethernet, USB 3.0+, fibre).
\item \textbf{Asynchronous}: character by character with start/stop bits --- simple, cheap, high overhead (20--30\%), good for irregular low-speed data (RS-232, keyboard, GPS). \textbf{Synchronous}: clocked blocks/frames with sync flags --- efficient for large volumes (overhead $<5\%$), complex, costly, used in LAN/WAN/storage (Ethernet, HDLC, PPP, Fibre Channel).
\item Selection criteria: \cle{speed}, \cle{cost}, \cle{distance}, \cle{volume of data}, \cle{directionality}, \cle{regularity of traffic}.
\item A real link combines one from each family: e.g., USB 3.0 = serial + full-duplex + synchronous; RS-232 modem = serial + full-duplex + asynchronous; DDR5 RAM = parallel + simplex (per direction) + synchronous.
\end{itemize}
\end{aretenir}

\coursec{Transmission Media: Guided and Unguided (Lesson 30)}

\courssub{Introduction: The Physical Layer in Action}

Imagine you are the network administrator of a new secondary school campus in Bafoussam. You have to connect 50 computers in the main laboratory to a server \SI{50}{m} away, link the administration block \SI{300}{m} across the playground, and provide internet access to a partner school in a remote village \SI{40}{km} away with no existing cable infrastructure. You cannot use the same "wire" for all three links. The physical path --- the \cle{transmission medium} --- dictates the speed, distance, cost and reliability of every connection. In Lessons 28 and 29 we studied \emph{how} data is encoded (bandwidth, modes, serial/parallel, synchronous/asynchronous). Now we study \emph{where} the signals travel: the media themselves.

\courssub{Definition and Classification of Transmission Media}

\begin{definition}[Transmission medium]
A \cle{transmission medium} is the physical path or channel through which data signals (electrical, optical or electromagnetic) propagate from a transmitter to a receiver in a communication system.
\end{definition}

\begin{definition}[Guided and unguided media]
Transmission media are divided into two great families:
\begin{itemize}
\item A \cle{guided} (or \emph{wired}, \emph{bounded}) medium provides a physical conductor that constrains and guides the signal along a fixed path. The three main types are \textbf{twisted pair cable}, \textbf{coaxial cable} and \textbf{fibre optic cable}.
\item An \cle{unguided} (or \emph{wireless}, \emph{unbounded}) medium transports the signal through the atmosphere or free space as electromagnetic waves, without a physical conductor. The main types are \textbf{radio waves}, \textbf{microwaves}, \textbf{satellite transmission} and \textbf{infrared}.
\end{itemize}
\end{definition}

\begin{aretenir}[Guided vs. Unguided --- the fundamental difference]
In guided media, the signal energy is confined inside a solid material (copper or glass); in unguided media, the energy spreads in space. This single fact explains almost all differences in bandwidth, distance, interference, security and installation cost.
\end{aretenir}

\courssub{Guided (Wired) Transmission Media}

\textbf{1. Twisted Pair Cable}\par
A twisted pair cable consists of two insulated copper wires twisted around each other in a helix. The twisting ensures that both wires are equally exposed to external electromagnetic fields, so the induced noise (common-mode interference) is nearly identical on both wires and cancels out at the differential receiver. This also reduces \cle{crosstalk} between adjacent pairs in the same sheath.

Two shielding variants exist:

\begin{itemize}
\item \cle{UTP} (Unshielded Twisted Pair): no metallic foil or braid around the pairs. It is inexpensive, thin, flexible and easy to terminate with RJ-45 connectors. It is the universal medium for \textbf{Local Area Networks (LANs)} (Ethernet) and traditional telephone lines.
\item \cle{STP} (Shielded Twisted Pair): each pair (or the whole bundle) is wrapped in a metallic foil or braid, often with a drain wire. It offers superior protection against external electromagnetic interference (EMI) but is thicker, heavier, more expensive, harder to install and requires careful grounding. It is used in electrically noisy environments (factories, near heavy machinery).
\end{itemize}

UTP cables are standardised into \emph{categories} (Cat) by the TIA/EIA-568 standard. For GCE purposes, know the most common ones:

\begin{center}
\begin{tabularx}{\linewidth}{|l|X|X|X|}
\hline
\rowcolor{popblueL} \textbf{Category} & \textbf{Max Data Rate (typical)} & \textbf{Bandwidth (MHz)} & \textbf{Common Application} \\
\hline
Cat 5e & \SI{1}{Gbit/s} & 100 & Standard office LAN, Gigabit Ethernet \\
\hline
Cat 6 & \SI{10}{Gbit/s} (up to \SI{55}{m}) & 250 & High-performance LAN, PoE+ \\
\hline
Cat 6a & \SI{10}{Gbit/s} (full \SI{100}{m}) & 500 & Data centres, 10GBASE-T \\
\hline
\end{tabularx}
\end{center}

\begin{attention}[The \SI{100}{m} Rule]
In Ethernet networks (10BASE-T, 100BASE-TX, 1000BASE-T, 10GBASE-T), the maximum length of a single copper segment (horizontal link) is \textbf{\SI{100}{metres}} (\SI{90}{m} solid cable + \SI{10}{m} patch cords). Beyond this, signal attenuation and timing violations require a \textbf{switch} or \textbf{repeater}. Writing "unlimited distance" for twisted pair in the exam is a fatal error.
\end{attention}

\textbf{2. Coaxial Cable}\par
A coaxial cable ("coax") is built in concentric layers sharing a common geometric axis:
\begin{enumerate}
\item A solid or stranded \textbf{central copper core} (conductor).
\item A \textbf{dielectric insulator} (solid polyethylene or foam) maintaining constant spacing.
\item A \textbf{metallic shield} (braided copper, aluminium foil, or both) acting as the return conductor and Faraday cage against EMI.
\item A protective \textbf{outer jacket} (PVC or plenum-rated).
\end{enumerate}

Because the shield confines the electromagnetic field between the core and the shield, coaxial cable supports much higher frequencies (hundreds of MHz to several GHz) and longer distances than twisted pair. It was the backbone of early Ethernet (10BASE2 "Thinnet", 10BASE5 "Thicknet" in bus topology). Today it remains the standard for:
\begin{itemize}
\item \textbf{Cable Television (CATV)} distribution (RG-6, 75 \SI{}{\ohm}).
\item Connecting \textbf{satellite dishes} (LNB) to decoders (RG-6).
\item \textbf{Cable Internet} (DOCSIS) last-mile access.
\end{itemize}

\textbf{3. Fibre Optic Cable}\par
Fibre optic cable guides \emph{light pulses} (typically \SI{850}{nm}, \SI{1310}{nm} or \SI{1550}{nm}) through a high-purity glass or plastic \textbf{core}. The core is surrounded by a \textbf{cladding} with a slightly lower refractive index. When light travelling in the core strikes the core-cladding interface at an angle greater than the \cle{critical angle}, it is reflected back into the core rather than transmitted into the cladding. This \cle{total internal reflection} allows light to propagate over kilometres with minimal loss.

Two propagation modes define the two main fibre types:

\begin{definition}[Single-mode fibre (SMF)]
Core diameter $\approx$ \SI{9}{\micro\meter}. Only \emph{one} propagation path (mode) exists. Pulse spreading (modal dispersion) is virtually eliminated. Used with \textbf{laser} sources (\SI{1310}{nm}/\SI{1550}{nm}). Distances: \textbf{tens to hundreds of kilometres} without regeneration. Standard for \textbf{inter-city backbones}, submarine cables and long-haul WANs.
\end{definition}

\begin{definition}[Multi-mode fibre (MMF)]
Core diameter \SI{50}{\micro\meter} (OM2/OM3/OM4/OM5) or \SI{62.5}{\micro\meter} (OM1). \emph{Multiple} light paths exist; rays entering at different angles arrive at different times, causing \cle{modal dispersion} that limits bandwidth-distance product. Used with cheaper \textbf{LED} or \textbf{VCSEL} sources (\SI{850}{nm}). Distances: \textbf{a few hundred metres to \SI{2}{km}}. Common inside buildings (campus LANs, data centres).
\end{definition}

\begin{propriete}[Advantages and drawbacks of fibre optics]
\begin{itemize}
\item \textbf{Bandwidth:} Highest of all guided media (Terabits/s on a single strand with WDM).
\item \textbf{Distance:} Tens of km (SMF) without repeaters.
\item \textbf{Immunity:} Zero electromagnetic interference (EMI/RFI); safe in high-voltage environments.
\item \textbf{Security:} Extremely difficult to tap without detection (light leak causes immediate loss).
\item \textbf{Drawbacks:} High material and installation cost (specialised tools, fusion splicers, trained technicians); fragility (glass core breaks if bent beyond minimum bend radius); connectors require polishing.
\end{itemize}
\end{propriete}

\begin{popfigure}
\begin{center}
\begin{tikzpicture}[scale=0.95]
% Twisted pair (UTP)
\draw[thick, popblue, fill=popblueL] (0,0) circle (1.2);
\foreach \a in {45,135,225,315}{
  \draw[thick, poporange, fill=poporange!25] (\a:0.6) circle (0.3);
  \draw[thick, popdark, fill=popgold] (\a:0.6) circle (0.13);}
\node[align=center, font=\small] at (0,-1.7) {Twisted Pair (UTP)};
\node[align=center, font=\scriptsize, text=popmuted] at (0,-2.15) {4 twisted pairs, PVC jacket, no shield};
% Coaxial
\begin{scope}[xshift=4.5cm]
\draw[thick, popblue, fill=popblueL] (0,0) circle (1.2);
\draw[thick, popdark, fill=white] (0,0) circle (0.9);
\draw[thick, poppurple, fill=poppurple!20] (0,0) circle (0.68);
\draw[thick, white, fill=white] (0,0) circle (0.38);
\draw[thick, popdark, fill=popgold] (0,0) circle (0.18);
\node[align=center, font=\small] at (0,-1.7) {Coaxial Cable};
\node[align=center, font=\scriptsize, text=popmuted] at (0,-2.15) {Jacket, Braid, Foil, Dielectric, Core};
\end{scope}
% Fibre Optic
\begin{scope}[xshift=9.0cm]
\draw[thick, popblue, fill=popblueL] (0,0) circle (1.2);
\draw[thick, popgreen, fill=popgreenL] (0,0) circle (0.78);
\draw[thick, poporange, fill=poporange!30] (0,0) circle (0.32);
\node[align=center, font=\small] at (0,-1.7) {Fibre Optic Cable};
\node[align=center, font=\scriptsize, text=popmuted] at (0,-2.15) {Jacket, Strength member, Cladding, Core};
\end{scope}
\end{tikzpicture}
\end{center}
\legende{Cross-sections of the three guided media. From left: UTP (4 pairs), Coaxial (concentric layers), Fibre Optic (core/cladding).}
\end{popfigure}

\begin{popfigure}
\begin{center}
\begin{tikzpicture}[scale=1.1]
% Cladding boundaries
\draw[thick, popgreen] (0,0.8) -- (10,0.8);
\draw[thick, popgreen] (0,-0.8) -- (10,-0.8);
\fill[popgreenL, opacity=0.3] (0,0.8) rectangle (10,1.2);
\fill[popgreenL, opacity=0.3] (0,-0.8) rectangle (10,-1.2);
\node[font=\scriptsize, text=popdark] at (9.2,1.0) {cladding ($n_2$)};
\node[font=\scriptsize, text=popdark] at (9.2,-1.0) {cladding ($n_2$)};
\node[font=\scriptsize, text=poporange] at (0.8,0.4) {core ($n_1 > n_2$)};
% Light ray zigzag (total internal reflection)
\draw[very thick, poporange, ->] (-0.5,0) -- (1.2,0.8);
\draw[very thick, poporange] (1.2,0.8) -- (3.6,-0.8);
\draw[very thick, poporange] (3.6,-0.8) -- (6.0,0.8);
\draw[very thick, poporange] (6.0,0.8) -- (8.4,-0.8);
\draw[very thick, poporange, ->] (8.4,-0.8) -- (10.5,0.0);
% Critical angle annotation
\draw[dashed, thin, popmuted] (3.6,-0.8) -- (3.6,0.2);
\draw[<->, thin, popink] (3.65,-0.75) arc (-90:-30:0.35);
\node[font=\tiny, text=popink] at (4.1,-0.4) {$\theta_c$};
\node[align=center, font=\small, text=popink] at (5.0,-1.6) {Total internal reflection at core--cladding boundary};
\end{tikzpicture}
\end{center}
\legende{Light propagation in a step-index optical fibre by total internal reflection. $n_1 > n_2$ ensures confinement.}
\end{popfigure}

\courssub{Unguided (Wireless) Transmission Media}

Unguided media use an \textbf{antenna} to radiate electromagnetic waves into space. The key parameters are \textbf{frequency} (determines bandwidth, propagation behaviour, antenna size) and \textbf{directionality} (omnidirectional vs. directional).

\textbf{1. Radio Waves (LF/MF/HF/VHF/UHF --- \SIrange{30}{3000}{MHz})}\par
Radio waves are typically generated by an antenna fed with alternating current. At lower frequencies they follow the Earth's curvature (ground wave) or reflect off the ionosphere (sky wave). At VHF/UHF (\SI{30}{MHz}--\SI{3}{GHz}) they propagate essentially in \textbf{line-of-sight (LOS)} but diffract around obstacles and penetrate buildings reasonably well. They are naturally \textbf{omnidirectional} (radiate in all horizontal directions) unless shaped by antenna arrays.

\begin{itemize}
\item \textbf{Applications:} FM/AM broadcasting, \textbf{Mobile cellular networks} (2G/3G/4G/5G) --- in Cameroon operated by \textbf{MTN} and \textbf{Orange}, \textbf{Wi-Fi} (IEEE 802.11, \SI{2.4}{GHz} and \SI{5}{GHz} bands), \textbf{Bluetooth} (\SI{2.4}{GHz} ISM band), LoRaWAN, Zigbee.
\item \textbf{Characteristics:} Moderate bandwidth (shared medium), susceptible to interference and multipath fading, easily intercepted $\rightarrow$ \textbf{encryption mandatory} (WPA3, AES).
\end{itemize}

\textbf{2. Microwaves (\SIrange{1}{300}{GHz})}\par
Microwaves are high-frequency radio waves. Because their wavelength is centimetric to millimetric, they can be focused into a \textbf{narrow, directional beam} using parabolic dish antennas (high gain). This allows point-to-point links with high bandwidth and frequency reuse.

\begin{propriete}[Line-of-Sight (LOS) Requirement]
\cle{Microwaves} (and infrared) require an unobstructed \textbf{line of sight} between the transmitting and receiving antennas. The Fresnel zone (an ellipsoidal volume around the visual LOS) must be at least 60\% clear to avoid diffraction losses. Because the Earth curves ($\approx$ \SI{8}{cm} drop per \SI{1}{km}), terrestrial microwave relays are mounted on towers or hilltops every \SI{30}{--}\SI{50}{km}.
\end{propriete}

\begin{itemize}
\item \textbf{Applications:} \textbf{Backhaul links} connecting mobile base stations to the core network, inter-city trunk links, private corporate WANs between buildings (building-to-building), satellite feeder links.
\item \textbf{Characteristics:} High bandwidth (hundreds of Mbps to Gbps), low interference (narrow beam), requires precise alignment, affected by heavy rain (rain fade above \SI{10}{GHz}).
\end{itemize}

\textbf{3. Satellite Transmission}\par
A communication satellite is a microwave relay station in orbit. The most common for fixed services is the \cle{geostationary (GEO) satellite}:
\begin{itemize}
\item Orbit: $\approx$ \SI{35\,786}{km} above the equator.
\item Period: \SI{24}{h} (matches Earth rotation) $\rightarrow$ appears fixed in the sky.
\item Footprint: Covers $\approx$ \SI{1}{3} of Earth's surface.
\end{itemize}
Communication flow: \textbf{Uplink} (Earth $\rightarrow$ Satellite, e.g. \SI{14}{GHz} Ku-band) $\rightarrow$ Transponder (amplifies, frequency shifts) $\rightarrow$ \textbf{Downlink} (Satellite $\rightarrow$ Earth, e.g. \SI{12}{GHz}).

\begin{itemize}
\item \textbf{VSAT (Very Small Aperture Terminal):} Small dishes (\SI{0.75}{--}\SI{2.4}{m}) for two-way data (internet, banking, SCADA) in remote areas. In Cameroon, VSATs connect rural schools, health centres and bank agencies where CAMTEL fibre has not reached.
\item \textbf{Satellite TV (DTH):} Receive-only (downlink) for bouquets like Canal+.
\item \textbf{Drawback --- Latency:} Round-trip distance $\approx$ \SI{72\,000}{km}. Propagation delay $\approx$ \SI{240}{ms} one-way, \SI{480}{ms} round-trip. This degrades real-time applications (VoIP, video conferencing, gaming, TCP throughput).
\end{itemize}

\begin{savaistu}[New orbits: LEO constellations]
Low Earth Orbit (LEO) constellations (Starlink, OneWeb, Kuiper) operate at \SI{550}{--}\SI{1200}{km} altitude. Latency drops to \SI{20}{--}\SI{40}{ms}, rivalling fibre. They require user terminals with phased-array antennas that track moving satellites. This technology is rapidly reaching Cameroon and changing the "remote connectivity" landscape.
\end{savaistu}

\textbf{4. Infrared (IR) (\SIrange{300}{400}{THz}, \SIrange{750}{1000}{nm})}\par
Infrared uses light just below visible red. It is \textbf{non-penetrating}: blocked by walls, doors, even paper. This confines communication to a single room (excellent security, no licence needed) but limits range to a few metres.

\begin{itemize}
\item \textbf{Applications:} TV/decoder/AC remote controls (consumer IR, \SI{38}{kHz} modulation), legacy IrDA ports on old laptops/phones (up to \SI{4}{Mbps}), free-space optics (FSO) for building-to-building links (high bandwidth, requires strict LOS, affected by fog).
\end{itemize}

\begin{popfigure}
\begin{center}
\begin{tikzpicture}[scale=0.85]
% Ground line
\draw[thick, popdark] (0,0) -- (14,0);
% Hills for microwave relays
\draw[thick, popgreen, fill=popgreenL] (4.0,0) .. controls (4.8,1.8) and (5.8,1.8) .. (6.6,0);
\draw[thick, popgreen, fill=popgreenL] (9.5,0) .. controls (10.3,1.5) and (11.2,1.5) .. (12.0,0);
% FM Radio Tower (Omnidirectional)
\draw[thick, popink] (1.5,0) -- (1.5,2.5);
\draw[thick, popink] (1.2,0) -- (1.5,2.5) -- (1.8,0);
\foreach \r in {0.6, 1.1, 1.6} \draw[popblue, thick] (1.5,2.5) circle (\r);
\node[align=center, font=\scriptsize] at (1.5,-0.6) {FM/TV Tower\\(Omnidirectional)};
% Microwave dishes on hills
\draw[thick, popink] (5.3,1.75) -- (5.3,2.5);
\draw[thick, poporange, fill=poporange!20] (5.3,2.5) arc (200:340:0.4);
\draw[thick, popink] (10.75,1.45) -- (10.75,2.2);
\draw[thick, poporange, fill=poporange!20] (10.75,2.2) arc (200:340:0.4);
\draw[very thick, poppurple, dashed, ->] (5.7,2.65) -- (10.35,2.45);
\node[align=center, font=\scriptsize, text=poppurple] at (8,3.1) {Microwave Link (LOS)};
% Geostationary Satellite
\draw[thick, popdark, fill=popgold] (12.5,5.5) rectangle (13.3,6.0);
\draw[thick, popblue, fill=popblueL] (11.7,5.6) rectangle (12.5,5.9);
\draw[thick, popblue, fill=popblueL] (13.3,5.6) rectangle (14.1,5.9);
\node[align=center, font=\scriptsize] at (12.9,6.4) {GEO Satellite};
% VSAT Dish on ground
\draw[thick, popink] (13.0,0) -- (13.0,0.6);
\draw[thick, poporange, fill=poporange!20] (13.0,0.6) arc (250:110:0.5);
\draw[very thick, poppink, dashed, ->] (13.15,1.1) -- (12.85,5.4);
\node[align=center, font=\scriptsize] at (13.0,-0.6) {VSAT Terminal};
% Mobile phone / Wi-Fi symbols near FM tower
\node[font=\scriptsize, text=popblue, align=center] at (3.0,1.2){4G/5G \\ Wi-Fi \\ Bluetooth};
\end{tikzpicture}
\end{center}
\legende{Unguided media landscape: Omnidirectional FM/Cellular tower, Point-to-point Microwave relay chain (LOS), GEO Satellite with VSAT ground station.}
\end{popfigure}

\courssub{Comparison of Transmission Media}

The following table summarises the key parameters for examination answers. Values are typical orders of magnitude for standard deployments.

\begin{center}
\begin{tabularx}{\linewidth}{|l|X|X|X|X|}
\hline
\rowcolor{popblueL} \textbf{Medium} & \textbf{Typical Bandwidth} & \textbf{Max Distance (no repeater)} & \textbf{Interference / Security} & \textbf{Typical Use Cases} \\
\hline
UTP (Cat5e/6) & \SI{1}{--}\SI{10}{Gbit/s} & \SI{100}{m} (Ethernet) & Poor (EMI, crosstalk, easy tap) & LANs, Telephony, PoE \\
\hline
STP & \SI{1}{--}\SI{10}{Gbit/s} & \SI{100}{m} & Good (shielded, needs grounding) & Industrial LANs, High EMI areas \\
\hline
Coaxial (RG-6/11) & \SI{100}{Mbit/s}{--}\SI{10}{Gbit/s} (DOCSIS 3.1) & \SI{200}{--}\SI{500}{m} & Good (shielded) & Cable TV, Satellite LNB, Broadband \\
\hline
Fibre Optic (SMF) & \SI{10}{Gbit/s}{--}\SI{100+}{Tbit/s} (WDM) & \SI{40}{--}\SI{100+}{km} & Immune (EMI), Very hard to tap & Backbones, Metro, FTTH, Submarine \\
\hline
Fibre Optic (MMF) & \SI{1}{--}\SI{100}{Gbit/s} & \SI{300}{m}{--}\SI{2}{km} & Immune (EMI), Very hard to tap & Campus LAN, Data Centres \\
\hline
Radio (Wi-Fi/4G) & \SI{10}{Mbit/s}{--}\SI{10}{Gbit/s} (shared) & \SI{10}{m}{--}\SI{30}{km} (cell) & High interference, Easy intercept & Mobile, WLAN, IoT, Broadcast \\
\hline
Microwave (PTP) & \SI{100}{Mbit/s}{--}\SI{10}{Gbit/s} & \SI{30}{--}\SI{50}{km}/hop & Low (narrow beam), Needs LOS & Backhaul, Building-to-Building \\
\hline
Satellite (GEO) & \SI{10}{Mbit/s}{--}\SI{1}{Gbit/s} (shared) & Continental (Footprint) & Interceptable, High Latency (\SI{480}{ms} RTT) & TV, Remote Internet, Maritime \\
\hline
Infrared (IR) & \SI{1}{--}\SI{4}{Mbit/s} (IrDA) / \SI{10}{Gbit/s} (FSO) & \SI{1}{--}\SI{10}{m} (room) / \SI{1}{--}\SI{2}{km} (FSO) & Blocked by obstacles, Secure (room) & Remotes, Legacy IrDA, FSO links \\
\hline
\end{tabularx}
\end{center}

\courssub{Choosing a Transmission Medium: A Structured Method}

GCE questions often ask: "A company wants to connect Site A to Site B. Recommend a suitable transmission medium and justify your choice." Use this method:

\begin{methode}[Selecting a transmission medium --- Step by step]
\begin{enumerate}
\item \textbf{Analyse the Distance:} Same room (\SI{<10}{m})? Same building (\SI{<100}{m})? Campus (\SI{<2}{km})? Inter-city (\SI{>10}{km})? Remote/Isolated?
\item \textbf{Determine Required Bandwidth:} Simple data entry (\SI{<10}{Mbit/s})? Standard office (\SI{1}{Gbit/s})? Heavy video/backup/Data Centre (\SI{10}{Gbit/s}+)?
\item \textbf{Assess Budget Constraints:} Cable cost (fibre \textgreater\ coax \textgreater\ UTP), Installation cost (trenching, fusion splicing \textgreater\ pulling UTP), Equipment cost (switches, media converters, VSAT terminals).
\item \textbf{Evaluate the Environment:} High EMI (factories, near power lines $\rightarrow$ Fibre/STP)? Obstacles/Hills (block LOS $\rightarrow$ avoid Microwave/IR)? No ducts/poles (Wireless)? Corrosive/Outdoor (Armoured cable)?
\item \textbf{Consider Security \& Reliability:} Military/Banking $\rightarrow$ Fibre (hard to tap). High availability $\rightarrow$ Redundant paths (Fibre + Microwave).
\item \textbf{State the Choice Clearly:} Name the specific medium (e.g., "Single-mode Fibre Optic", "UTP Cat6a", "VSAT Ku-band", "Microwave PTP link").
\item \textbf{Justify against the criteria:} "Chosen because distance is \SI{15}{km} (exceeds copper), bandwidth requirement is \SI{10}{Gbit/s} (exceeds microwave), environment has high EMI (fibre immune), budget allows fibre installation."
\end{enumerate}
\end{methode}

\begin{exoresolu}[Choosing media for a school campus and remote link]
A secondary school in Bamenda wants to: (1) Connect 30 PCs in a computer laboratory to a server in the same building, \SI{60}{m} away. Budget is tight. (2) Connect the school campus to the internet backbone in the town centre, \SI{15}{km} away. Propose a medium for each link and justify.
\tcblower
\textbf{Solution.}

\textbf{Link 1: Laboratory to Server (\SI{60}{m}, tight budget).}
\begin{itemize}
\item \textbf{Choice:} \textbf{UTP Category 6 (or Cat5e) Twisted Pair} with Gigabit Ethernet switches.
\item \textbf{Justification:} Distance (\SI{60}{m}) is well within the \SI{100}{m} Ethernet limit. Bandwidth (\SI{1}{Gbit/s}) is sufficient for a school lab. UTP is the cheapest guided medium to purchase and install (RJ-45 crimping requires minimal skill/tools). No high EMI reported in a school.
\end{itemize}

\textbf{Link 2: Campus to Town Backbone (\SI{15}{km}).}
\begin{itemize}
\item \textbf{Primary Choice:} \textbf{Single-Mode Fibre Optic (SMF)} (e.g., OS2, \SI{1310}{nm} lasers).
\item \textbf{Justification:} Distance (\SI{15}{km}) far exceeds copper limits. SMF offers virtually unlimited bandwidth for future growth, total immunity to lightning/EMI (critical for outdoor runs), and high security. CAMTEL's national backbone in Cameroon is fibre-based.
\item \textbf{Alternative (if fibre impossible):} \textbf{Licensed Microwave Point-to-Point Link} (\SI{18}{GHz} or \SI{23}{GHz} band).
\item \textbf{Justification for Alternative:} Can span \SI{15}{km} in a single hop if LOS exists (mount hilltop masts). High bandwidth (\SI{1}{Gbit/s}+). Lower civil works cost (no trenching), but requires spectrum licence (ANTIC), precise alignment, and suffers rain fade.
\end{itemize}
\end{exoresolu}

\begin{attention}[Wireless $\neq$ "No Equipment" / "No Cost"]
A classic GCE trap: "Wireless is cheaper because no cables." \textbf{False.} Unguided media require: Towers/Masts, High-gain Antennas (Dishes, Sector antennas), Radio Transceivers (ODU/IDU), Power systems (Solar/Battery/Generator for remote sites), Spectrum Licence fees (ANTIC in Cameroon). A \SI{15}{km} microwave link often costs \emph{more} than a fibre trench if ducts exist. Always compare \textbf{Total Cost of Ownership (TCO)}.
\end{attention}

\begin{attention}[UTP vs STP --- The Shield Confusion]
\begin{itemize}
\item \textbf{UTP} = \textbf{U}nshielded Twisted Pair. \textbf{No} metallic foil/braid. Relies only on twisting for noise rejection.
\item \textbf{STP} = \textbf{S}hielded Twisted Pair. \textbf{Has} metallic shield (foil/braid) \textbf{per pair} (U/FTP) or \textbf{overall} (F/UTP, S/FTP).
\item \textbf{Sc/FTP} (Screened Foiled Twisted Pair) = Overall braid + Foil per pair (Best shielding).
\end{itemize}
Writing "UTP has a shield" loses marks. Remember: \textbf{S} in STP stands for \textbf{Shielded}.
\end{attention}

\begin{attention}[Fibre: Single-mode vs Multi-mode Connectors]
Never connect a Single-mode transceiver (laser) to a Multi-mode patch cord, or vice-versa, without a mode-conditioning patch cord (for legacy 1Gbps LX on MMF). Core diameter mismatch (\SI{9}{\micro\meter} vs \SI{50}{\micro\meter}) causes massive insertion loss (\SI{>20}{dB}). In the exam, state: "SMF uses laser, yellow jacket (usually); MMF uses LED/VCSEL, orange/aqua/violet jacket (usually)."
\end{attention}

\begin{savaistu}[Cameroon's Connectivity Backbone]
Cameroon lands international submarine fibre cables (WACS, SAT-3, ACE, SAIL) at \textbf{Douala}. \textbf{CAMTEL} operates the national fibre backbone (thousands of km) linking Douala, Yaoundé, Bafoussam, Bamenda, Garoua, Maroua, etc. \textbf{MTN} and \textbf{Orange} lease capacity on this backbone and deploy their own fibre/microwave for last-mile and backhaul. Rural Universal Service Fund projects heavily use \textbf{VSAT} and \textbf{Solar-powered Microwave} relays on hilltops (e.g., Adamawa, Far North regions) where fibre deployment is economically unviable. This mix of media is a textbook case of "Right medium for the geography".
\end{savaistu}

\begin{exercice}[Media Selection for a Bank]
A commercial bank headquartered in Yaoundé has the following requirements:
\begin{enumerate}
\item Link Headquarters to a new Branch \SI{8}{km} away. Requirement: \SI{10}{Gbit/s}, maximum security (financial data), high reliability.
\item Provide wireless internet access for customers in the HQ waiting room (approx. \SI{200}{m^2}).
\item Connect a new micro-finance agency in a remote village in the East Region. No fibre, no microwave LOS, no mobile coverage (3G/4G). Requirement: Basic banking transactions (low bandwidth, high availability).
\end{enumerate}
For each case (a), (b), (c), choose the \textbf{most appropriate transmission medium} and \textbf{justify} your choice using the criteria: Distance, Bandwidth, Environment, Security, Cost.
\end{exercice}

\begin{corrige}[Exercise 1]
\begin{enumerate}
\item \textbf{(a) Headquarters $\leftrightarrow$ Branch (\SI{8}{km}): Single-Mode Fibre Optic (SMF).}
\textbf{Justification:} Distance (\SI{8}{km}) exceeds copper limits. Bandwidth (\SI{10}{Gbit/s}) requires fibre (or licensed microwave, but fibre is more secure/reliable). Security: Fibre is immune to EMI and extremely difficult to tap without detection (critical for banking). Reliability: No rain fade, no LOS alignment drift. Cost: High initial CAPEX, but lowest OPEX and future-proof.
\item \textbf{(b) HQ Waiting Room Wi-Fi: Radio Waves (Wi-Fi 6 / 802.11ax, \SI{2.4}{GHz} + \SI{5}{GHz}).}
\textbf{Justification:} Distance: Short range (single room). Bandwidth: Shared \SI{1}{--}\SI{5}{Gbit/s} sufficient for guest browsing. Environment: Indoor, omnidirectional coverage needed (penetrates walls). Cost: Low (Access Points + PoE switch). Security: WPA3-Enterprise / Captive Portal isolates guest traffic from bank LAN.
\item \textbf{(c) Remote Village Agency: VSAT Satellite (Ku-band or Ka-band GEO, or LEO e.g. Starlink).}
\textbf{Justification:} Distance: Continental coverage, no terrestrial infrastructure needed. Environment: No fibre, no LOS for microwave, no cellular. Bandwidth: Low (\SI{1}{--}\SI{50}{Mbit/s}) sufficient for transactional traffic (ATM, POS, core banking sync). Latency (\SI{500}{ms} GEO / \SI{40}{ms} LEO) acceptable for non-real-time banking. Cost: Moderate CAPEX (Dish, Modem, Router), recurring OPEX (Bandwidth subscription). Reliability: High (satellite uptime \SI{99.5}{\%}+), independent of local cuts.
\end{enumerate}
\end{corrige}

\begin{aretenir}[Key Points for Revision]
\begin{itemize}
\item \textbf{Guided Media:} Twisted Pair (UTP/STP, \SI{100}{m} limit, Categories 5e/6/6a), Coaxial (Concentric, Shielded, CATV/Satellite), Fibre Optic (Light, Total Internal Reflection, SMF \SI{9}{\micro\meter} core \SI{>40}{km} Laser, MMF \SI{50}{\micro\meter} core \SI{<2}{km} LED/VCSEL, Highest BW, Immune to EMI, Costly/Fragile).
\item \textbf{Unguided Media:} Radio (Omnidirectional, Penetrates walls, FM/Wi-Fi/Bluetooth/4G/5G, Interference/Interception), Microwave (Directional, LOS required, \SI{30}{--}\SI{50}{km}/hop, Backhaul), Satellite (GEO \SI{36\,000}{km}, VSAT, High Latency \SI{480}{ms} RTT, Remote areas), Infrared (Short range, Blocked by obstacles, Remote controls/FSO).
\item \textbf{Selection Criteria:} Distance $\rightarrow$ Bandwidth $\rightarrow$ Budget $\rightarrow$ Environment (EMI, LOS, Ducts) $\rightarrow$ Security $\rightarrow$ Justify.
\item \textbf{Traps:} Wireless needs expensive equipment/licences; UTP is NOT shielded; SMF $\neq$ MMF connectors; \SI{100}{m} copper limit; Satellite Latency.
\end{itemize}
\end{aretenir}

\newpage

\coursec{Integration activity}

\coursec{Identifying different Transmission Types}

\courssub{Introduction}

Data transmission is the backbone of every computer network. Whether you are sending a WhatsApp message from Douala to Yaoundé, streaming a lesson from the Ministry of Education's platform, or submitting GCE registration data to the Board's server, the information travels across a \cle{transmission medium} using a specific \cle{transmission mode} and occupies a certain \cle{bandwidth}. This chapter equips you with the knowledge to analyse a real-world connectivity problem, compare the available technologies, and justify your technical choices — exactly what a network technician or ICT consultant does every day in Cameroon.

\begin{aretenir}[Key concepts of this chapter]
\begin{itemize}
\item \textbf{Bandwidth} is the range of frequencies a medium can carry, measured in hertz (Hz); it determines the maximum theoretical bit rate.
\item \textbf{Attenuation} is the loss of signal strength as it travels; it limits the maximum distance of a link without repeaters.
\item \textbf{Electromagnetic interference (EMI)} from generators, power lines, or lightning can corrupt signals on copper cables; optical fibre and properly shielded twisted pair are immune or resistant.
\item \textbf{Guided media} (twisted pair, coaxial, optical fibre) confine the signal to a physical path; \textbf{unguided media} (radio, microwave, satellite) propagate through free space.
\item \textbf{Transmission modes}: simplex (one way), half-duplex (both ways but not simultaneously), full-duplex (both ways simultaneously).
\item \textbf{Serial vs parallel}: serial sends bits sequentially on one channel; parallel sends multiple bits simultaneously on multiple channels. Over distance, serial wins because of skew and cost.
\item \textbf{Broadband} divides a wide band into multiple channels (or uses a wide channel at high bit rate) to carry many simultaneous communications; \textbf{narrowband} uses a single narrow channel for one low-rate communication.
\end{itemize}
\end{aretenir}

\courssub{Lesson 28: Broadband and Narrowband Transmission}

\begin{definition}[Bandwidth and Channel Capacity]
The \cle{bandwidth} of a transmission medium is the difference between the highest and lowest frequencies it can effectively carry, expressed in hertz (Hz). According to the Nyquist and Shannon-Hartley theorems, the maximum data rate (capacity) is proportional to bandwidth. A \cle{narrowband} channel occupies a small frequency range (typically < 64 kbit/s in traditional telephony) and carries a single low-bit-rate stream. A \cle{broadband} system uses a wide frequency range (MHz to GHz) to multiplex many channels or to achieve very high bit rates (Mbit/s to Gbit/s), allowing multiple users and services (voice, video, data) to share the medium simultaneously.
\end{definition}

\begin{propriete}[Broadband vs Narrowband in Practice]
\begin{itemize}
\item \textbf{Narrowband examples:} Traditional PSTN landline (3.1 kHz voice channel), SMS over GSM control channels, old dial-up modems (56 kbit/s).
\item \textbf{Broadband examples:} ADSL, fibre-to-the-home (FTTH), 4G/5G mobile data, Wi-Fi (802.11n/ac/ax), cable TV (DOCSIS), CAMTEL fibre corporate links.
\item \textbf{Examinable distinction:} Broadband enables \emph{simultaneous} multiple communications (frequency-division multiplexing, time-division multiplexing, or high-speed packet switching); narrowband does not.
\end{itemize}
\end{propriete}

\begin{exemplebox}[School internet subscription]
A school with 30 simultaneous users browsing, downloading past papers, and uploading registration files needs a broadband connection. A narrowband link (e.g. a single 64 kbit/s ISDN channel) would allow only one very slow session at a time. A broadband fibre subscription of 100 Mbit/s shared among 30 users gives ~3 Mbit/s each — sufficient for educational use.
\end{exemplebox}

\courssub{Lesson 29: Types of Data Transmission}

\begin{definition}[Transmission Modes]
\begin{itemize}
\item \textbf{Simplex}: Data flows in one direction only (e.g. radio broadcast, keyboard to CPU, PA bell system).
\item \textbf{Half-duplex}: Data flows in both directions but only one direction at a time (e.g. walkie-talkie, legacy Ethernet hub, USB 1.x control transfers).
\item \textbf{Full-duplex}: Data flows in both directions simultaneously (e.g. telephone conversation, modern switched Ethernet, USB 3.x, fibre links with separate Tx/Rx fibres).
\end{itemize}
\end{definition}

\begin{definition}[Serial and Parallel Transmission]
\begin{itemize}
\item \textbf{Parallel transmission}: Multiple bits (usually 8, 16, 32 or 64) are sent simultaneously over multiple parallel wires. Used internally in computers (PCI, old IDE, old printer port) and for very short distances.
\item \textbf{Serial transmission}: Bits are sent one after another over a single channel (or a differential pair). Used for all long-distance communication (Ethernet, USB, SATA, HDMI, PCIe, fibre).
\end{itemize}
\end{definition}

\begin{propriete}[Why Serial Dominates Over Distance]
\begin{enumerate}
\item \textbf{Cost and bulk}: Parallel needs $n$ signal wires + ground + control lines $\rightarrow$ thick, expensive, fragile cables. Serial needs 1 or 2 pairs.
\item \textbf{Skew and timing}: Over metres, signals on different wires arrive at slightly different times (\cle{clock skew}). At high speeds, the receiver cannot reliably reassemble the parallel word. Serial has no inter-wire skew.
\item \textbf{Crosstalk and EMI}: Many parallel wires in close proximity couple noise into each other. A twisted pair or fibre rejects common-mode noise.
\item \textbf{Result}: Modern high-speed links (USB 3, SATA, PCIe, 10G Ethernet) are all serial, running at multi-Gbit/s rates over metres.
\end{enumerate}
\end{propriete}

\begin{methode}[Choosing the Transmission Mode for a Link]
\begin{enumerate}
\item Identify whether the communicating devices ever need to send and receive \emph{at the same time}.
\item If yes $\rightarrow$ full-duplex required (modern LAN switches, routers, fibre).
\item If communication is strictly alternating (like a walkie-talkie) $\rightarrow$ half-duplex acceptable.
\item If data only ever flows one way (sensor to logger, broadcast) $\rightarrow$ simplex sufficient.
\item For the physical layer: always use serial transmission for any link longer than a few centimetres.
\end{enumerate}
\end{methode}

\courssub{Lesson 30: Transmission Media (Guided and Unguided)}

\begin{definition}[Guided Transmission Media]
Signals are confined to a solid physical conductor.
\begin{itemize}
\item \textbf{Twisted Pair (UTP/STP)}: Two insulated copper wires twisted together. Categories: Cat 5e (100 MHz, 1 Gbit/s up to 100 m), Cat 6 (250 MHz, 1 Gbit/s up to 100 m, 10 Gbit/s up to 55 m), Cat 6a (500 MHz, 10 Gbit/s up to 100 m). Cheap, easy to terminate (RJ45), but susceptible to EMI and limited to 100 m for Ethernet.
\item \textbf{Coaxial Cable}: Central conductor, dielectric, braided shield, outer jacket. Higher bandwidth than twisted pair, better EMI immunity, used for cable TV (DOCSIS) and legacy Ethernet (10BASE2/5). Distance up to 500 m (thicknet) but bulky and harder to install.
\item \textbf{Optical Fibre}: Glass or plastic core guides light. \textbf{Multimode} (OM3/OM4, 50/125 µm) for LAN up to 550 m at 10 Gbit/s; \textbf{Singlemode} (OS2, 9/125 µm) for km to 100+ km at 100 Gbit/s+. Immune to EMI, no lightning conduction, highest bandwidth, but higher cost, specialised termination (fusion splicing), and fragile if not protected.
\end{itemize}
\end{definition}

\begin{definition}[Unguided Transmission Media]
Signals propagate as electromagnetic waves through air/vacuum.
\begin{itemize}
\item \textbf{Radio / Wi-Fi (IEEE 802.11)}: 2.4 GHz, 5 GHz, 6 GHz bands. Shared medium, half-duplex (CSMA/CA), susceptible to interference, distance 30–100 m indoor, up to several km outdoor with directional antennas.
\item \textbf{Microwave (Point-to-Point)}: Licensed bands (6–42 GHz), highly directional dishes, line-of-sight required, distances up to 50–100 km, used for backbone links (MTN, Orange, CAMTEL backhaul).
\item \textbf{Satellite}: Geostationary (high latency ~600 ms) or LEO (Starlink, lower latency). Used for remote areas.
\item \textbf{Bluetooth / Zigbee / LoRa}: Short-range, low-power IoT.
\end{itemize}
\end{definition}

\begin{center}
\begin{tabularx}{\linewidth}{|l|X|X|X|X|}
\hline
\rowcolor{popblueL} \textbf{Criterion} & \textbf{UTP Cat 6} & \textbf{STP / FTP Cat 6} & \textbf{Multimode Fibre (OM4)} & \textbf{Singlemode Fibre (OS2)} \\
\hline
Max distance (Ethernet) & 100 m (1 Gbit/s) & 100 m (1 Gbit/s) & 400 m (10 Gbit/s), 550 m (1 Gbit/s) & 10 km – 100+ km \\
Bandwidth & 250 MHz & 250 MHz & 4700 MHz·km & $\approx$ 100 THz·km \\
Cost per metre & Low & Medium & Medium–High & High \\
EMI immunity & Low & Good (shielded) & Total (dielectric) & Total (dielectric) \\
Lightning / surge risk & High & Medium (shield to ground) & None & None \\
Installation ease & Very easy (RJ45 crimp) & Easy (shield grounding needed) & Moderate (fusion splicer or pre-terminated) & Moderate–Hard \\
Typical use in Cameroon & Indoor LAN, short outdoor in conduit & Outdoor between buildings < 100 m, near generators & Building backbone, campus links < 550 m & ISP backhaul, long campus links, FTTH \\
\hline
\end{tabularx}
\end{center}

\begin{attention}[Common Misconceptions]
\begin{itemize}
\item ``Optical fibre is always more expensive than copper.'' — For a single 300 m link, the \emph{total} cost (cable + 2 media converters + labour) may be comparable to copper + repeaters + surge protectors, and fibre is future-proof.
\item ``Wi-Fi is always cheaper than cabling.'' — Outdoor point-to-point Wi-Fi needs two radios, two directional antennas, masts, lightning arrestors, and clear line-of-sight; for 300 m across a field with a generator, fibre is often more reliable and not much dearer.
\item ``Cat 6 cable works at 10 Gbit/s over 100 m.'' — Only Cat 6a guarantees 10 Gbit/s at 100 m; Cat 6 is limited to 55 m at 10 Gbit/s.
\item ``Full-duplex doubles the bandwidth.'' — Full-duplex gives 1 Gbit/s \emph{each way} simultaneously; the aggregate throughput is 2 Gbit/s, but a single flow still maxes at 1 Gbit/s.
\end{itemize}
\end{attention}

\courssub{Integration Activity: Connecting Government Bilingual High School Bafoussam}

\begin{aretenir}[Aims of the activity]
By the end of this integration activity, you should be able to:
\begin{itemize}
\item analyse a real connectivity problem and extract the technical constraints (distance, bandwidth, budget, interference);
\item choose an appropriate \cle{transmission medium} (guided or unguided) for each link of a network and justify the choice using comparison criteria (cost, bandwidth, attenuation, immunity to interference, ease of installation);
\item select the correct \cle{transmission mode} (simplex, half-duplex, full-duplex) for a given communication need;
\item distinguish \cle{broadband} from \cle{narrowband} and argue why an internet subscription for a school must be broadband;
\item explain why \cle{serial transmission} is used inside network cables rather than parallel transmission;
\item present and defend a technical solution orally, using a labelled diagram, and evaluate competing solutions critically.
\end{itemize}
\end{aretenir}

\textbf{Situation: Connecting Government Bilingual High School Bafoussam.}\par

The principal of Government Bilingual High School (GBHS) Bafoussam, in the West Region of Cameroon, has obtained funding from the Parent-Teachers' Association to modernise the school's computer facilities. The school campus is laid out as follows:

\begin{itemize}
\item \textbf{Block A} contains Computer Lab 1 with 20 desktop computers.
\item \textbf{Block B}, situated \SI{80}{m} from Block A, contains Computer Lab 2 with another 20 desktop computers.
\item The \textbf{Administration block} (principal's office, bursar's office, secretary's office: 6 computers and 1 shared printer) is \SI{300}{m} away from Block A, on the other side of the school football field.
\item A CAMTEL optical fibre line passes along the road at the school gate, about \SI{20}{m} from Block A. The school can subscribe to an internet connection there.
\end{itemize}

The principal's requirements are:

\begin{enumerate}
\item All 40 computers in the two labs must be connected into a single \cle{local area network} (LAN) so that students can share files, a printer and educational software.
\item The Administration block must be linked to the LAN so that report cards and administrative documents can be exchanged electronically.
\item The whole school must have internet access. During examination periods, at least 30 students may be online simultaneously, and the administration must be able to submit candidates' files to the GCE Board online platform without interruption.
\item The total budget is limited: the school cannot afford to lay optical fibre everywhere, and any outdoor cable must survive the rainy season and the dusty dry season in Bafoussam.
\item The path between Block A and the Administration block crosses the football field, near the school's generator house, where a large diesel generator produces strong electromagnetic interference when running.
\end{enumerate}

The principal has contacted your class, the Upper Sixth ICT students, to act as consultants. Working in groups, you must design the school's network and defend your choices.

\courssub{Tasks}

\begin{enumerate}
\item \textbf{Identify the links.} List all the separate communication links that must be created (for example: computers to switch inside Lab 1, Lab 1 to Lab 2, LAN to Administration, LAN to the internet). For each link, state its approximate length and the environment it crosses (indoor, outdoor, near a source of interference).

\item \textbf{Choose a transmission medium for each link.} For every link identified in question 1, propose a guided medium (twisted pair, coaxial cable, optical fibre) or an unguided medium (radio/Wi-Fi, microwave). Justify each choice using at least three comparison criteria studied in this chapter: maximum distance before attenuation becomes a problem, bandwidth, cost per metre, immunity to electromagnetic interference, and ease of installation.

\item \textbf{Choose the transmission mode.} Between the two labs, and between the LAN and the Administration block, users must send and receive data at the same time (a teacher downloads a file while students upload their work). State whether the links must operate in \cle{simplex}, \cle{half-duplex} or \cle{full-duplex} mode, and justify your answer. Give one example of a device in the school that would only need simplex transmission.

\item \textbf{Broadband or narrowband?} The internet provider offers two subscriptions: a cheap narrowband connection and a more expensive broadband connection. Using the definitions of \cle{broadband} and \cle{narrowband} transmission, explain which subscription the school should choose, referring to the requirement that 30 students may be online simultaneously.

\item \textbf{Serial or parallel?} Inside the network cables connecting the computers, data travels one bit after another. State whether this is \cle{serial} or \cle{parallel} transmission, and explain two reasons why this choice is made for distances of several metres or more.

\item \textbf{Draw and defend.} Produce a labelled diagram of the complete network showing all blocks, all media used, the switches, the router and the internet connection. Present your solution to the class and defend it against the constraints of cost, distance and interference. Then, as a class, compare the different group solutions and agree on the best compromise.
\end{enumerate}

\courssub{Model Answer}

\textbf{Question 1 — Identifying the links.}\par

Five distinct links are needed:

\begin{center}
\begin{tabularx}{\linewidth}{|l|X|l|l|}
\hline
\rowcolor{popblueL} \textbf{Link} & \textbf{Description} & \textbf{Length} & \textbf{Environment} \\
\hline
L1 & Computers to switch, inside each lab & $\leq \SI{10}{m}$ & Indoor \\
\hline
L2 & Lab 1 (Block A) to Lab 2 (Block B) & \SI{80}{m} & Outdoor, between buildings \\
\hline
L3 & LAN (Block A) to Administration & \SI{300}{m} & Outdoor, crosses field, near generator \\
\hline
L4 & Block A to the fibre access point at the gate & \SI{20}{m} & Outdoor, short \\
\hline
L5 & Router to internet provider (CAMTEL) & --- & Provider's fibre \\
\hline
\end{tabularx}
\end{center}

\textbf{Question 2 — Choice of transmission media.}\par

\begin{itemize}
\item \textbf{L1 (inside the labs): UTP twisted pair (Cat 5e or Cat 6).} The distances are very short (a few metres), well below the \SI{100}{m} limit of Ethernet over twisted pair. UTP cable is cheap, easy to install with RJ45 connectors, and offers bandwidths of 100 Mbit/s to 1 Gbit/s, more than enough for student workstations. Interference is low indoors, so unshielded cable is acceptable.

\item \textbf{L2 (Lab 1 to Lab 2, \SI{80}{m}): UTP/STP twisted pair or optical fibre.} At \SI{80}{m}, the link is just within the \SI{100}{m} Ethernet limit, so a good-quality shielded twisted pair (STP), run in a protective PVC conduit overhead or buried, is the cheapest acceptable solution. A group with a slightly larger budget may propose optical fibre for greater reliability and future upgrades; both answers are defensible if the justification is sound.

\item \textbf{L3 (LAN to Administration, \SI{300}{m}, near the generator): optical fibre.} This is the critical link. Twisted pair is impossible because \SI{300}{m} far exceeds the \SI{100}{m} limit (signal attenuation would make the link unusable without repeaters). Moreover, the diesel generator produces strong electromagnetic interference, and copper cable crossing an open field is also exposed to lightning-induced surges during the rainy season. Optical fibre is immune to electromagnetic interference, covers \SI{300}{m} easily (multimode fibre reaches at least \SI{550}{m} at gigabit speed) and does not conduct lightning surges. Its higher cost is justified on this single link only. A point-to-point Wi-Fi link with directional antennas is an acceptable alternative if the budget is extremely tight, but it is less reliable in heavy rain and offers lower bandwidth.

\item \textbf{L4 (Block A to the gate, \SI{20}{m}): optical fibre or UTP.} Since the provider's equipment is fibre-based, a short fibre patch into the school's router is natural; a UTP link from the provider's indoor modem is also possible.

\item \textbf{L5 (internet access): CAMTEL optical fibre subscription} at the roadside, terminated by the provider's modem/router in Block A.
\end{itemize}

\textbf{Question 3 — Transmission mode.}\par

The links between the labs and between the LAN and the Administration must operate in \cle{full-duplex} mode: data must flow in both directions simultaneously, because while the bursar downloads examination files, the secretary may be uploading enrolment data on the same link. Modern switched Ethernet is full-duplex by nature (each device has dedicated send and receive circuits and there are no collisions). Half-duplex (one direction at a time, like a walkie-talkie or old Ethernet hubs) would halve the effective throughput and is unacceptable here. An example of a simplex device in the school is the public-address bell system or a radio receiver in the staff room: information flows in one direction only.

\textbf{Question 4 — Broadband or narrowband?}\par

The school must choose the \cle{broadband} subscription. Narrowband transmission uses a single narrow channel and offers a low bit rate: it cannot carry many simultaneous communications. Broadband transmission divides a wide bandwidth into many channels (or uses a wide channel at high bit rate), so multiple users and multiple services (web browsing, e-mail, video lessons, GCE online registration) can share the connection at the same time. With 30 students online simultaneously, a narrowband link would be saturated and the administration's GCE submissions would be interrupted, violating requirement 3. Broadband is therefore technically necessary, even though it costs more.

\textbf{Question 5 — Serial or parallel?}\par

Transmission inside network cables is \cle{serial}: bits are sent one after another on a single pair of wires. Two reasons explain this choice:

\begin{enumerate}
\item \textbf{Cost and simplicity over distance:} parallel transmission needs one wire per bit (at least 8 or 16 wires plus control lines), which makes cables thick, expensive and fragile over tens or hundreds of metres; serial needs only one or two pairs.
\item \textbf{Skew and interference:} over long distances, the bits of a parallel word travel at slightly different speeds on different wires (clock skew) and arrive misaligned, causing errors at high speed; a serial stream has no such problem and can therefore run much faster over long distances. This is why even inside computers, modern technologies (USB, SATA) have abandoned parallel buses in favour of very fast serial links.
\end{enumerate}

\textbf{Question 6 — Labelled diagram and defence.}\par

\begin{popfigure}
\begin{center}
\begin{tikzpicture}[scale=0.95, every node/.style={font=\small}]
% Blocks
\node[draw=popblue, thick, rounded corners, fill=popblueL, minimum width=2.6cm, minimum height=1.2cm, align=center] (lab1) at (0,0) {Block A\\ Lab 1 (20 PCs)};
\node[draw=popblue, thick, rounded corners, fill=popblueL, minimum width=2.6cm, minimum height=1.2cm, align=center] (lab2) at (6.5,0) {Block B\\ Lab 2 (20 PCs)};
\node[draw=poppurple, thick, rounded corners, fill=popblueL, minimum width=2.8cm, minimum height=1.2cm, align=center] (admin) at (0,-4) {Administration\\ (6 PCs + printer)};
\node[draw=popgreen, thick, rounded corners, fill=popgreenL, minimum width=2.6cm, minimum height=1.2cm, align=center] (gate) at (6.5,-4) {School gate\\ CAMTEL fibre};
% Switches / router
\node[draw=popdark, fill=white, rounded corners, inner sep=2pt] (sw1) at (0,1.6) {Switch 1};
\node[draw=popdark, fill=white, rounded corners, inner sep=2pt] (sw2) at (6.5,1.6) {Switch 2};
\node[draw=poporange, thick, fill=white, rounded corners, inner sep=2pt] (rt) at (3.2,-2.6) {Router};
% Links
\draw[thick, popblue] (lab1) -- (sw1);
\draw[thick, popblue] (lab2) -- (sw2);
\draw[thick, popblue] (sw1) -- node[above]{UTP \SI{80}{m}} (sw2);
\draw[thick, poporange] (sw1) -- node[left, align=center]{Fibre \SI{300}{m}\\ full-duplex} (admin);
\draw[thick, popgreen] (rt) -- node[below]{Fibre \SI{20}{m}} (gate);
\draw[thick, popdark] (sw1) -- (rt);
\node[align=center, text=popmuted] at (3.2,2.4) {All LAN links: serial, full-duplex};
\end{tikzpicture}
\end{center}
\legende{Proposed network for GBHS Bafoussam: UTP inside labs and between labs, optical fibre to the Administration block and to the CAMTEL access point.}
\end{popfigure}

In the class debate, the best compromise agreed upon is: copper (UTP) everywhere distances are under \SI{100}{m} and interference is low, because it is cheap and easy to maintain; optical fibre only on the long, interference-exposed \SI{300}{m} link to the Administration and for the internet access, because only fibre satisfies the distance, immunity and bandwidth requirements; full-duplex switched Ethernet on all LAN links; and a broadband fibre subscription, sized for at least 30 simultaneous users. This solution respects the limited budget while meeting every requirement of the principal.

\begin{attention}[Common mistakes to avoid in your defence]
\begin{itemize}
\item Proposing UTP for the \SI{300}{m} link: beyond \SI{100}{m} the Ethernet signal is too attenuated.
\item Forgetting the generator's electromagnetic interference and the lightning risk when choosing copper across the field.
\item Confusing half-duplex (one direction at a time) with full-duplex (both directions simultaneously).
\item Claiming parallel transmission is used in network cables because ``it is faster'': over long distances it is slower and unreliable due to skew.
\item Choosing the cheapest (narrowband) subscription without checking the simultaneous-user requirement.
\end{itemize}
\end{attention}

\begin{savaistu}[Cameroon Context: National Backbone and School Connectivity]
Cameroon's national fibre backbone, managed by CAMTEL, connects major cities (Douala, Yaoundé, Bafoussam, Bamenda, Garoua, Maroua, etc.). The \emph{Projet d'Appui à la Connectivité des Établissements Scolaires} (PACES) and partnerships with ANTIC aim to bring broadband to schools. Many schools still rely on 3G/4G modems (MTN, Orange) as backup or primary link where fibre is not yet at the gate. Understanding the trade-offs between fibre, microwave, and mobile broadband is a key skill for ICT professionals in Cameroon.
\end{savaistu}

\begin{exercice}[Practice Questions]
\begin{enumerate}
\item A company in Douala wants to connect two buildings 1.2 km apart across a busy highway. They have a limited budget but need 1 Gbit/s full-duplex. Compare the options: (a) singlemode fibre, (b) licensed microwave link, (c) Cat 6a with repeaters every 100 m. Which would you recommend and why?
\item Define the following terms: bandwidth, attenuation, electromagnetic interference, simplex, half-duplex, full-duplex, serial transmission, parallel transmission, broadband, narrowband.
\item A school uses a 4G router with a 20 Mbit/s subscription for 50 computers. During exams, the connection becomes unusable. Explain, using the concepts of broadband and contention ratio, why this happens and propose two technical improvements.
\item Draw a labelled diagram of a home network with: fibre ONT, router, 8-port switch, 3 wired PCs, a Wi-Fi access point, a smart TV, and a network printer. Indicate the media, transmission modes, and whether each link is serial or parallel.
\end{enumerate}
\end{exercice}

\begin{corrige}[Exercise 1]
\textbf{(a) Singlemode fibre:} Best technical choice. Distance 1.2 km is trivial for singlemode (supports 10 km+ at 10 Gbit/s). Immune to EMI from highway traffic and lightning. Full-duplex native. Higher upfront cost (cable, two media converters/switches with SFP+ modules, fusion splicing or pre-terminated), but zero recurring cost and future-proof.
\textbf{(b) Licensed microwave:} Requires frequency licence from ANTIC, clear line-of-sight (may need masts on both roofs), professional installation. Latency very low, bandwidth high. Recurring licence fee. Good if trenching fibre is impossible (e.g. highway crossing permission denied).
\textbf{(c) Cat 6a with repeaters:} Needs 11 repeaters (switches) powered and weatherproofed along 1.2 km. Each adds cost, power requirement, failure point, and latency. Not recommended for outdoor highway crossing.
\textbf{Recommendation:} Singlemode fibre if trenching or aerial pole route is feasible; otherwise licensed microwave.
\end{corrige}

\begin{corrige}[Exercise 2]
\begin{itemize}
\item \textbf{Bandwidth:} Range of frequencies a medium can carry (Hz), determines max data rate.
\item \textbf{Attenuation:} Signal power loss over distance (dB/km or dB/100m).
\item \textbf{EMI:} Unwanted electromagnetic signals coupling into a cable, corrupting data.
\item \textbf{Simplex:} One-way only (e.g. keyboard $\rightarrow$ PC).
\item \textbf{Half-duplex:} Two-way but not simultaneous (e.g. walkie-talkie, Wi-Fi).
\item \textbf{Full-duplex:} Two-way simultaneous (e.g. switched Ethernet, fibre).
\item \textbf{Serial:} Bits sent sequentially on one channel (e.g. Ethernet, USB).
\item \textbf{Parallel:} Multiple bits sent simultaneously on multiple wires (e.g. old PCI, IDE).
\item \textbf{Broadband:} Wide bandwidth carrying multiple channels / high bit rate, many simultaneous users.
\item \textbf{Narrowband:} Single narrow channel, low bit rate, one user/service at a time.
\end{itemize}
\end{corrige}

\begin{corrige}[Exercise 3]
The 4G connection is broadband (wide channel, high bit rate) but the \emph{contention ratio} (number of users sharing the cell sector) and the \emph{shared medium} (TDMA/OFDMA scheduling) mean the 20 Mbit/s is the \emph{peak} sector capacity, not per user. With 50 computers, even 10 active users get ~2 Mbit/s each; during exams, many more are active, causing congestion, high latency, packet loss. Improvements: (1) Subscribe to a dedicated fibre or fixed wireless (microwave) link with committed information rate (CIR). (2) Implement QoS on the router to prioritise exam traffic (GCE portal) over recreational traffic. (3) Cache exam resources locally (proxy server) to reduce WAN demand.
\end{corrige}

\begin{corrige}[Exercise 4]
Diagram description: ONT (fibre) $\rightarrow$ Router (WAN port) $\rightarrow$ 8-port Switch (LAN ports). Three PCs $\rightarrow$ Switch via UTP (serial, full-duplex). Wi-Fi AP $\rightarrow$ Switch via UTP (serial, full-duplex); Smart TV and Printer connect to AP via Wi-Fi (serial, half-duplex CSMA/CA). All wired links: UTP Cat 6, serial, full-duplex. Wireless links: radio, serial, half-duplex.
\end{corrige}

\newpage

\coursec{Methods and worked examples}

\coursec{Identifying different Transmission Types}

\courssub{Introduction and Syllabus Overview}

\begin{aretenir}[Chapter ICT08 -- What you must master]
This chapter covers three lessons from the official Upper Sixth syllabus:
\begin{itemize}
\item \textbf{Lesson 28:} Broadband and Narrowband Transmission
\item \textbf{Lesson 29:} Types of data transmission (simplex, half-duplex, full-duplex; serial vs parallel; synchronous vs asynchronous)
\item \textbf{Lesson 30:} Transmission media -- guided (twisted pair, coaxial, optical fibre) and unguided (infrared, radio, microwave, satellite)
\end{itemize}
By the end of this chapter you will be able to:
\begin{enumerate}
\item Convert and compare bit rates and bandwidths, and compute transfer times.
\item Classify a link as broadband or narrowband with justification.
\item Identify the transmission mode (simplex, half-duplex, full-duplex) from a description.
\item Distinguish serial from parallel transmission and explain why serial dominates long-distance links.
\item Distinguish synchronous from asynchronous transmission and calculate framing overhead.
\item Select the appropriate guided medium (UTP, STP, coaxial, fibre) for a given scenario.
\item Select the appropriate unguided medium (infrared, radio, microwave, satellite) for a given scenario.
\item Synthesise all concepts in a multi-part GCE-style question.
\end{enumerate}
\end{aretenir}

\begin{savaistu}[Bandwidth in Cameroon]
Cameroon's international bandwidth comes mainly through the \textbf{SAT-3/WASC} and \textbf{WACS} submarine fibre-optic cables landing at Kribi and Limbe. CAMTEL manages the national backbone. Mobile operators (MTN, Orange, Nexttel) provide 3G/4G broadband; VSAT satellite links serve remote areas. Understanding transmission types helps you evaluate real offers like ``MTN Home Broadband'' or ``Orange Fibre''.
\end{savaistu}

\courssub{Skill 1: Converting between bit rate units and choosing the right bandwidth vocabulary}

\begin{methode}[Handling bit rates and bandwidth figures]
When a GCE question mentions a bandwidth or a bit rate, proceed methodically.
\begin{enumerate}
\item \textbf{Identify the quantity.} Bandwidth (in hertz, Hz) is the width of the frequency range a channel can carry; bit rate (in bits per second, bit/s) is the number of bits transmitted per second. Do not confuse them.
\item \textbf{Master the prefixes.} $1\ \mathrm{kbit/s} = 10^3\ \mathrm{bit/s}$, $1\ \mathrm{Mbit/s} = 10^6\ \mathrm{bit/s}$, $1\ \mathrm{Gbit/s} = 10^9\ \mathrm{bit/s}$. The same prefixes apply to Hz.
\item \textbf{Convert before comparing.} Always express all values in the same unit before ranking or subtracting them.
\item \textbf{Link to the transmission type.} A \cle{narrowband} channel carries a small bandwidth (typically voice-grade, a few kHz) and low bit rates; a \cle{broadband} channel carries a wide bandwidth and high bit rates, often shared between several simultaneous signals.
\item \textbf{Compute transfer time when asked.} Use $t = \dfrac{\text{size of data (bits)}}{\text{bit rate (bit/s)}}$, remembering that $1$ byte $= 8$ bits.
\end{enumerate}
\end{methode}

\begin{attention}[Common traps]
\begin{itemize}
\item Confusing \textbf{MB} (megabyte, $10^6$ bytes) with \textbf{Mbit} (megabit, $10^6$ bits). Always convert file size to bits: multiply by $8$.
\item Thinking ``M'' is always bigger than ``k'': $0.5\ \mathrm{Mbit/s} = 500\ \mathrm{kbit/s}$, which is \emph{slower} than $512\ \mathrm{kbit/s}$.
\item Forgetting that practical throughput is lower than the nominal bit rate due to protocol overhead (headers, acknowledgements, retransmissions).
\end{itemize}
\end{attention}

\begin{exoresolu}[Converting and comparing bit rates]
A cybercaf\'e in Bamenda offers three Internet packages: package A at $512\ \mathrm{kbit/s}$, package B at $2\ \mathrm{Mbit/s}$ and package C at $0.5\ \mathrm{Mbit/s}$.
\begin{enumerate}
\item Express all three bit rates in bit/s and rank the packages from slowest to fastest.
\item A student downloads a file of $4\ \mathrm{MB}$ (megabytes) using package B. Calculate the theoretical download time.
\item State, with a reason, whether package B is a narrowband or broadband service.
\end{enumerate}
\tcblower
\textbf{1. Conversion and ranking.} Convert each value to bit/s:
\begin{align*}
\text{A: } & 512\ \mathrm{kbit/s} = 512 \times 10^3 = 512\,000\ \mathrm{bit/s}\\
\text{B: } & 2\ \mathrm{Mbit/s} = 2 \times 10^6 = 2\,000\,000\ \mathrm{bit/s}\\
\text{C: } & 0.5\ \mathrm{Mbit/s} = 0.5 \times 10^6 = 500\,000\ \mathrm{bit/s}
\end{align*}
Ranking from slowest to fastest: C ($500\,000$) $<$ A ($512\,000$) $<$ B ($2\,000\,000$). Note the trap: $0.5\ \mathrm{Mbit/s}$ is \emph{slower} than $512\ \mathrm{kbit/s}$ even though ``M'' looks bigger than ``k''; always convert first.

\textbf{2. Download time.} Convert the file size to bits:
\[ 4\ \mathrm{MB} = 4 \times 10^6 \times 8 = 3.2 \times 10^7\ \mathrm{bits}. \]
Then apply the formula:
\[ t = \dfrac{3.2 \times 10^7}{2 \times 10^6} = 16\ \mathrm{s}. \]
The theoretical download time is $16$ seconds (in practice slightly longer because of protocol overheads and line quality).

\textbf{3. Classification.} Package B is a \cle{broadband} service: it offers a high bit rate (megabits per second), well above the voice-grade narrowband range, and broadband links can carry several services (voice, video, data) simultaneously over the same medium.
\end{exoresolu}

\begin{popfigure}
\begin{tikzpicture}[scale=0.9, every node/.style={font=\small}]
% Bit rate scale
\draw[->, thick] (0,0) -- (12,0) node[right] {Bit rate (bit/s)};
\foreach \x/\lbl in {0/0, 2/500k, 4/1M, 6/2M, 8/10M, 10/100M} {
  \draw (\x,-0.1) -- (\x,0.1) node[below] {\lbl};
}
% Mark packages
\draw[poporange, thick] (2.048,0.2) -- (2.048,-0.2) node[below, poporange] {C: 0.5 Mbit/s};
\draw[popblue, thick] (2.096,0.2) -- (2.096,-0.2) node[below, popblue] {A: 512 kbit/s};
\draw[popgreen, thick] (6,0.2) -- (6,-0.2) node[below, popgreen] {B: 2 Mbit/s};
% Bands
\draw[poppink, thick, decorate, decoration={brace, amplitude=5pt}] (0,-0.5) -- (3,-0.5) node[midway, below=8pt] {Narrowband (voice-grade)};
\draw[poppurple, thick, decorate, decoration={brace, amplitude=5pt}] (3,-0.5) -- (12,-0.5) node[midway, below=8pt] {Broadband};
\end{tikzpicture}
\legende{Bit rate scale showing narrowband vs broadband regions and the three packages.}
\end{popfigure}

\courssub{Skill 2: Distinguishing broadband from narrowband transmission}

\begin{definition}[Narrowband vs Broadband]
\begin{itemize}
\item \cle{Narrowband transmission}: Uses a narrow frequency band (typically $\le 4\ \mathrm{kHz}$ for voice-grade lines). Only one signal occupies the channel at a time. Low bit rate (up to $\approx 56\ \mathrm{kbit/s}$ for dial-up). Example: traditional PSTN telephone line with a dial-up modem.
\item \cle{Broadband transmission}: Uses a wide frequency band divided into multiple channels (via Frequency Division Multiplexing or other techniques). Many signals travel simultaneously. High bit rate (Mbit/s to Gbit/s). Examples: ADSL, fibre, 4G/5G, cable TV.
\end{itemize}
\end{definition}

\begin{methode}[Classifying a transmission as broadband or narrowband]
\begin{enumerate}
\item \textbf{Recall the definitions.} Narrowband: narrow band, one signal at a time, low bit rate. Broadband: wide band, many simultaneous signals, high bit rate.
\item \textbf{Look for key indicators in the question:} bit rate (low/high), number of simultaneous signals (one/many), typical technology (dial-up vs ADSL, fibre, 4G).
\item \textbf{Use a comparison table} to structure your answer in an exam: criterion, narrowband, broadband.
\item \textbf{Justify with the bandwidth}: the wider the frequency band available, the more channels and the higher the total bit rate.
\end{enumerate}
\end{methode}

\begin{exoresolu}[Broadband vs narrowband in context]
In a village near Bafoussam, an old dial-up modem connects to the Internet at $56\ \mathrm{kbit/s}$ through the telephone line, and while connected, no telephone call can be made. In Yaound\'e, a household uses an ADSL connection at $8\ \mathrm{Mbit/s}$: the family browses the Internet, watches television and makes phone calls on the same line at the same time.
\begin{enumerate}
\item Identify which connection is narrowband and which is broadband.
\item Explain why the ADSL line can carry several services simultaneously.
\item Give two other examples of broadband technologies available in Cameroon.
\end{enumerate}
\tcblower
\textbf{1. Classification.} The dial-up modem at $56\ \mathrm{kbit/s}$ is \cle{narrowband}: it uses a single narrow channel (the voice band of the telephone line) and carries only one signal at a time, which is why the telephone cannot be used during the connection. The ADSL connection at $8\ \mathrm{Mbit/s}$ is \cle{broadband}: high bit rate and several simultaneous services.

\textbf{2. Explanation.} ADSL divides the total bandwidth of the copper line into several frequency channels: a low-frequency channel for the telephone voice, and higher-frequency channels for downstream and upstream data. Because each service occupies its own frequency band (a technique called \emph{frequency division multiplexing}), they travel simultaneously without interfering. Broadband means precisely this: a wide band shared between many signals.

\textbf{3. Other broadband technologies in Cameroon.} Examples: 4G/LTE mobile Internet offered by MTN and Orange; fibre-optic access deployed along the CAMTEL backbone; cable Internet; satellite broadband (VSAT). Any two of these are acceptable.
\end{exoresolu}

\begin{center}
\begin{tabularx}{\linewidth}{|l|X|X|}
\hline
\rowcolor{popblueL} \textbf{Criterion} & \textbf{Narrowband} & \textbf{Broadband} \\
\hline
Bandwidth & Narrow (few kHz) & Wide (MHz to GHz) \\
\hline
Simultaneous signals & One at a time & Many (multiplexed) \\
\hline
Typical bit rate & $\le 56\ \mathrm{kbit/s}$ & $\ge 256\ \mathrm{kbit/s}$ (often Mbit/s) \\
\hline
Example & Dial-up modem & ADSL, Fibre, 4G, Cable, Satellite \\
\hline
\end{tabularx}
\end{center}

\courssub{Skill 3: Identifying the transmission mode (simplex, half-duplex, full-duplex)}

\begin{definition}[Transmission Modes]
\begin{itemize}
\item \cle{Simplex}: Data flows in \textbf{one direction only}. The receiver cannot transmit back on the same channel. Example: TV/radio broadcasting, keyboard $\to$ computer.
\item \cle{Half-duplex}: Data flows in \textbf{both directions but alternately} (one at a time). The channel is shared in time. Example: walkie-talkie (push-to-talk), CB radio.
\item \cle{Full-duplex}: Data flows in \textbf{both directions simultaneously}. Requires two separate channels (or frequency division). Example: telephone conversation, mobile phone call.
\end{itemize}
\end{definition}

\begin{methode}[Recognising simplex, half-duplex and full-duplex]
\begin{enumerate}
\item \textbf{Ask two questions} about the link: (i) can data flow in both directions? (ii) if yes, can it flow in both directions \emph{at the same time}?
\item \textbf{Apply the decision rule:}
\begin{itemize}
\item One direction only $\Rightarrow$ \cle{simplex}
\item Both directions but alternately $\Rightarrow$ \cle{half-duplex}
\item Both directions simultaneously $\Rightarrow$ \cle{full-duplex}
\end{itemize}
\item \textbf{Give an example} each time: examiners award marks for a correct definition \emph{plus} a correct example.
\item \textbf{Watch the wording:} ``users must press a button to speak'' signals half-duplex; ``both can speak and listen at the same time'' signals full-duplex.
\end{enumerate}
\end{methode}

\begin{popfigure}
\begin{tikzpicture}[scale=0.9, every node/.style={font=\small}]
% Simplex
\node[draw, rounded corners, fill=popblueL] (S1) at (0,2) {Transmitter};
\node[draw, rounded corners, fill=popblueL] (S2) at (4,2) {Receiver};
\draw[->, thick, popblue] (S1) -- (S2) node[midway, above] {Simplex: one way};
% Half-duplex
\node[draw, rounded corners, fill=poporangeL] (H1) at (0,0) {Station A};
\node[draw, rounded corners, fill=poporangeL] (H2) at (4,0) {Station B};
\draw[<->, thick, poporange, dashed] (H1) -- (H2) node[midway, above] {Half-duplex: alternate};
% Full-duplex
\node[draw, rounded corners, fill=popgreenL] (F1) at (0,-2) {Station A};
\node[draw, rounded corners, fill=popgreenL] (F2) at (4,-2) {Station B};
\draw[->, thick, popgreen] (F1) -- (F2) node[midway, above] {Forward};
\draw[->, thick, popgreen] (F2) -- (F1) node[midway, below] {Return};
\node at (2,-2.7) {Full-duplex: simultaneous};
\end{tikzpicture}
\legende{Transmission modes: simplex (one way), half-duplex (alternate), full-duplex (simultaneous two-way).}
\end{popfigure}

\begin{exoresolu}[Classifying communication situations]
For each situation below, state the transmission mode and justify your answer.
\begin{enumerate}
\item CRTV broadcasts television programmes to viewers' sets.
\item Two gendarmes on patrol communicate with walkie-talkies; each must press a button to talk and release it to listen.
\item A customer in Douala calls his bank on his mobile phone; both can speak and hear each other at the same time.
\end{enumerate}
\tcblower
\textbf{1. Television broadcasting: simplex.} The signal travels in one direction only, from the transmitter to the receivers. Viewers cannot send anything back through the same channel. Hence it is a \cle{simplex} transmission.

\textbf{2. Walkie-talkies: half-duplex.} Communication is possible in both directions (each gendarme can talk and can listen), but not simultaneously: the push-to-talk button forces the two parties to alternate. One channel is shared in time. Hence it is a \cle{half-duplex} transmission.

\textbf{3. Mobile phone call: full-duplex.} Data (voice) flows in both directions \emph{at the same time}: each speaker can talk while hearing the other, and interruptions are possible, exactly as in a face-to-face conversation. Hence it is a \cle{full-duplex} transmission.

\textbf{Summary table for revision:}
\begin{center}
\begin{tabularx}{\linewidth}{|l|X|X|}
\hline
\rowcolor{popblueL} \textbf{Mode} & \textbf{Direction of flow} & \textbf{Example} \\
\hline
Simplex & One direction only & TV/radio broadcasting, keyboard $\to$ CPU \\
\hline
Half-duplex & Both directions, alternately & Walkie-talkie, CB radio \\
\hline
Full-duplex & Both directions, simultaneously & Telephone call, mobile call, video conference \\
\hline
\end{tabularx}
\end{center}
\end{exoresolu}

\courssub{Skill 4: Distinguishing serial from parallel transmission}

\begin{definition}[Serial vs Parallel Transmission]
\begin{itemize}
\item \cle{Serial transmission}: Bits are sent \textbf{one after another} on a \textbf{single} data line (or a pair for full-duplex). Used for long distances: USB, RS-232, Ethernet, fibre.
\item \cle{Parallel transmission}: Multiple bits (e.g., 8, 16, 32) are sent \textbf{simultaneously} on \textbf{multiple} data lines. Used for very short distances: internal buses (PCI, ISA), old printer port (Centronics/LPT).
\end{itemize}
\end{definition}

\begin{methode}[Analysing serial vs parallel transmission]
\begin{enumerate}
\item \textbf{Count the data lines.} Serial: one line; Parallel: $k$ lines (e.g., 8).
\item \textbf{Compare speed and distance.} Parallel seems faster per clock cycle, but over long distances the bits on different wires arrive at slightly different times (\emph{skew}) and the cables are costly; serial links can be clocked much faster and are cheaper, so serial dominates long-distance communication.
\item \textbf{Compute effective times if asked:} time to send $n$ bits serially at rate $R$ is $t = \dfrac{n}{R}$; for parallel with $k$ lines, effective rate $= k \times R$, so $t = \dfrac{n}{kR}$.
\item \textbf{Give examples:} serial $\rightarrow$ USB, RS-232, Ethernet, fibre; parallel $\rightarrow$ old printer port (LPT), internal data bus.
\end{enumerate}
\end{methode}

\begin{attention}[Skew and crosstalk]
In parallel cables, \emph{skew} (different propagation delays) and \emph{crosstalk} (electromagnetic coupling between adjacent wires) worsen with distance and frequency. This limits parallel transmission to short runs (typically $< 2\ \mathrm{m}$ for high speeds). Serial links avoid these problems by using one or two differential pairs.
\end{attention}

\begin{exoresolu}[Serial vs parallel: a numerical comparison]
A school in Buea must send a document of $1\,600$ bits to a device.
\begin{enumerate}
\item Using a serial link at $800\ \mathrm{bit/s}$, calculate the transmission time.
\item Using a parallel link with $8$ wires, each wire carrying $800\ \mathrm{bit/s}$, calculate the transmission time.
\item Despite this advantage, explain why modern long-distance networks use serial rather than parallel transmission.
\end{enumerate}
\tcblower
\textbf{1. Serial link.} All $1\,600$ bits pass one by one on a single line:
\[ t_{\text{serial}} = \dfrac{1\,600}{800} = 2\ \mathrm{s}. \]

\textbf{2. Parallel link.} With $8$ wires, $8$ bits are transmitted per clock step, so the effective rate is $8 \times 800 = 6\,400\ \mathrm{bit/s}$:
\[ t_{\text{parallel}} = \dfrac{1\,600}{6\,400} = 0.25\ \mathrm{s}. \]
The parallel link is $8$ times faster for the same per-wire rate.

\textbf{3. Why serial wins in practice.} Over long distances, parallel cables suffer from \emph{skew}: signals on different wires do not arrive exactly together, causing errors at high speed. Parallel cables also need many conductors, making them thick, expensive and subject to crosstalk. Serial links use one or two wires, can be clocked at very high frequencies, and tolerate long distances. This is why USB, Ethernet and all wide-area networks are serial technologies, while parallel transmission survives only over very short distances (e.g.\ inside the computer, on the motherboard bus).
\end{exoresolu}

\begin{popfigure}
\begin{tikzpicture}[scale=0.8, every node/.style={font=\small}]
% Serial
\node[draw, rounded corners, fill=popblueL] (Ss) at (0,1) {Sender};
\node[draw, rounded corners, fill=popblueL] (Sr) at (6,1) {Receiver};
\draw[->, thick] (Ss) -- (Sr) node[midway, above] {Serial: 1 line, bits sequential};
\foreach \i in {1,...,4} {
  \node[circle, draw, fill=popblue, inner sep=1pt] at (1.2+\i*1, 1) {};
}
% Parallel
\node[draw, rounded corners, fill=poporangeL] (Ps) at (0,-1) {Sender};
\node[draw, rounded corners, fill=poporangeL] (Pr) at (6,-1) {Receiver};
\foreach \y/\bit in {0.2/7, 0.1/6, 0/5, -0.1/4, -0.2/3, -0.3/2, -0.4/1, -0.5/0} {
  \draw[->, thick] (Ps) -- ++(0,\y) -- (6, -1+\y) -- (Pr);
  \node at (3, -1+\y+0.15) {\tiny bit \bit};
}
\node at (3, -1.8) {Parallel: 8 lines, bits simultaneous};
\end{tikzpicture}
\legende{Serial vs parallel transmission: serial sends bits sequentially on one line; parallel sends multiple bits at once on multiple lines.}
\end{popfigure}

\courssub{Skill 5: Distinguishing synchronous and asynchronous transmission}

\begin{definition}[Synchronous vs Asynchronous Transmission]
\begin{itemize}
\item \cle{Asynchronous transmission}: Data is sent \textbf{character by character} (typically 7 or 8 bits). Each character is framed by a \emph{start bit} (logic 0) and one or two \emph{stop bits} (logic 1). Sender and receiver clocks are \textbf{independent}; the start bit resynchronises the receiver for each character. Idle gaps (line held at logic 1) are allowed between characters. Overhead: 2-3 bits per character. Example: keyboard input, RS-232, old modem links.
\item \cle{Synchronous transmission}: Data is sent in \textbf{continuous blocks (frames)}. Sender and receiver clocks are \textbf{synchronised} (by a separate clock line, embedded clock, or synchronisation pattern/preamble). No start/stop bits per character; only a few control bits per frame. Very high efficiency. Example: HDLC, Ethernet frames, fibre channel, data transfers between computers.
\end{itemize}
\end{definition}

\begin{methode}[Identifying synchronous vs asynchronous transmission]
\begin{enumerate}
\item \textbf{Check how timing is handled.} Asynchronous: start/stop bits per character, independent clocks, idle gaps allowed. Synchronous: continuous frames, common clock or sync bits, no idle gaps within a frame.
\item \textbf{Compare efficiency.} Asynchronous adds overhead (start/stop bits) to every character; synchronous has little overhead and suits high-speed, large-volume transfers.
\item \textbf{Compute the overhead if asked:} efficiency $= \dfrac{\text{data bits}}{\text{total bits transmitted}} \times 100\ \%$.
\item \textbf{Examples:} asynchronous $\rightarrow$ keyboard, old modem; synchronous $\rightarrow$ network frames, inter-computer transfers.
\end{enumerate}
\end{methode}

\begin{exoresolu}[Efficiency of asynchronous transmission]
A device sends data asynchronously. Each character consists of $8$ data bits, framed by $1$ start bit and $1$ stop bit.
\begin{enumerate}
\item How many bits are actually transmitted per character?
\item Calculate the efficiency of this transmission.
\item The line carries $12\,000$ characters per minute. Calculate the useful data rate in bit/s.
\item Explain why synchronous transmission would be preferred for transferring a $500\ \mathrm{MB}$ database backup.
\end{enumerate}
\tcblower
\textbf{1. Bits per character.} Each character carries $1$ start bit $+ 8$ data bits $+ 1$ stop bit $= 10$ bits transmitted.

\textbf{2. Efficiency.} Only the $8$ data bits are useful:
\[ \text{efficiency} = \dfrac{8}{10} \times 100\ \% = 80\ \%. \]
So $20\ \%$ of the line capacity is consumed by framing bits.

\textbf{3. Useful data rate.} Characters per second: $\dfrac{12\,000}{60} = 200$ characters/s. Each character yields $8$ useful bits:
\[ 200 \times 8 = 1\,600\ \mathrm{bit/s}. \]
(The raw line rate is $200 \times 10 = 2\,000\ \mathrm{bit/s}$, of which only $1\,600\ \mathrm{bit/s}$ is useful data.)

\textbf{4. Why synchronous is better for a large backup.} In synchronous transmission, data is sent in large continuous frames with only a few synchronisation and control bits per frame, so the overhead is proportionally tiny and the efficiency approaches $100\ \%$. For a $500\ \mathrm{MB}$ backup, the $20\ \%$ overhead of asynchronous transmission would waste about $100\ \mathrm{MB}$ of transferred volume and considerably lengthen the operation. Synchronous transmission, with both clocks locked together, sustains high bit rates and is therefore the standard choice for bulk data transfer between computers.
\end{exoresolu}

\begin{center}
\begin{tabularx}{\linewidth}{|l|X|X|}
\hline
\rowcolor{popblueL} \textbf{Aspect} & \textbf{Asynchronous} & \textbf{Synchronous} \\
\hline
Timing & Start/stop bits per character & Common clock / sync pattern \\
\hline
Clock & Independent & Synchronised \\
\hline
Gap between characters & Allowed (idle line) & Not within a frame \\
\hline
Overhead per character & 2-3 bits (20-30\%) & Few bits per frame ($<1\%$) \\
\hline
Typical use & Keyboard, RS-232, low speed & Networks, high-speed links \\
\hline
\end{tabularx}
\end{center}

\courssub{Skill 6: Selecting an appropriate guided transmission medium}

\begin{definition}[Guided Transmission Media]
Signals propagate along a physical cable. Three main types:
\begin{itemize}
\item \cle{Twisted pair} (UTP/STP): Two insulated copper wires twisted together. UTP = Unshielded; STP = Shielded. Cheap, easy to install. Max segment $\approx 100\ \mathrm{m}$ for Ethernet. Categories: Cat5e (1 Gbit/s), Cat6 (10 Gbit/s up to 55 m), Cat6a (10 Gbit/s up to 100 m). Sensitive to EMI (STP better).
\item \cle{Coaxial cable}: Central conductor, dielectric, braided shield, outer jacket. Better shielding, higher bandwidth than twisted pair. Used for cable TV (CATV), older Ethernet (10BASE2/5). Medium cost.
\item \cle{Optical fibre}: Glass/core carries light. Single-mode (long distance, high bandwidth) and multi-mode (shorter distance). Immune to EMI, very high bit rate (Gbit/s to Tbit/s), long distance (km to 100+ km), secure (hard to tap). Expensive, fragile, requires specialised termination.
\end{itemize}
\end{definition}

\begin{methode}[Choosing between twisted pair, coaxial cable and optical fibre]
\begin{enumerate}
\item \textbf{List the requirements} in the question: distance, required bit rate, budget, electromagnetic environment, security.
\item \textbf{Recall the characteristics of each guided medium:}
\begin{itemize}
\item Twisted pair: cheap, easy, limited distance ($\approx 100\ \mathrm{m}$), moderate bit rate, sensitive to EMI (STP shielded).
\item Coaxial: better shielding, higher bandwidth, used for cable TV and older LANs, medium cost.
\item Optical fibre: very high bit rate, very long distance, immune to EMI, secure, but expensive and fragile.
\end{itemize}
\item \textbf{Match medium to situation} and justify with at least two characteristics.
\item \textbf{Never answer with the medium name alone}: the GCE mark scheme rewards the justification.
\end{enumerate}
\end{methode}

\begin{exoresolu}[Choosing cabling for a school network]
A secondary school in Ngaound\'er\'e wants to network three situations:
\begin{enumerate}
\item Connecting $20$ computers inside one computer laboratory to a switch, over distances of at most $30\ \mathrm{m}$, with a small budget.
\item Linking two buildings $8\ \mathrm{km}$ apart at very high speed, in an area with many high-voltage power lines.
\item Carrying cable television signals to homes in a neighbourhood.
\end{enumerate}
For each situation, choose the most appropriate guided medium and justify your choice.
\tcblower
\textbf{1. Computer laboratory: UTP twisted pair (Cat5e or Cat6).} The distances are short (well under the $\approx 100\ \mathrm{m}$ Ethernet limit), the budget is small, and installation must be simple. Unshielded twisted pair is the cheapest guided medium, is easy to crimp and connect to the switch with RJ45 plugs, and comfortably supports Fast/Gigabit Ethernet over $30\ \mathrm{m}$. Fibre would be an unnecessary expense here.

\textbf{2. Link between two buildings $8\ \mathrm{km}$ apart: optical fibre (single-mode).} Three requirements point to fibre: (i) the distance of $8\ \mathrm{km}$ far exceeds what twisted pair or coaxial cable can carry without repeaters; (ii) very high speed is demanded, and fibre supports gigabit rates; (iii) the high-voltage lines create strong electromagnetic interference, to which fibre is completely immune because it transmits light, not electricity. Fibre is also difficult to tap, which secures the inter-building traffic.

\textbf{3. Cable television distribution: coaxial cable.} Coaxial cable has a metallic shield giving good protection against interference, a bandwidth wide enough to carry many TV channels simultaneously (broadband transmission), and a moderate cost acceptable for mass deployment to homes. This is precisely why cable TV operators historically built their networks with coaxial cable.

\textbf{Moral:} there is no ``best'' medium in absolute terms; the correct answer always depends on distance, bit rate, interference and cost.
\end{exoresolu}

\begin{popfigure}
\begin{tikzpicture}[scale=0.8, every node/.style={font=\small}]
% Twisted pair
\draw[popblue, thick] (0,2) -- (3,2);
\draw[poporange, thick] (0,1.8) -- (3,1.8);
\draw[decorate, decoration={snake, amplitude=1pt, segment length=4pt}] (0,2) -- (3,2);
\node at (3.5,1.9) {Twisted pair (UTP)};
% Coaxial
\draw[popblue, thick] (0,1) circle (0.15);
\draw[poporange, thick] (0,1) circle (0.35);
\draw[popblue, thick] (3,1) circle (0.15);
\draw[poporange, thick] (3,1) circle (0.35);
\draw[popblue, thick] (0.15,1) -- (2.85,1);
\draw[poporange, thick] (0.35,1) -- (2.65,1);
\node at (3.5,1) {Coaxial cable};
% Fibre
\draw[popgreen, thick] (0,0) -- (3,0);
\draw[popgreen!50, thick] (0,-0.1) -- (3,-0.1);
\draw[popgreen!50, thick] (0,0.1) -- (3,0.1);
\node at (3.5,0) {Optical fibre (core/cladding)};
\end{tikzpicture}
\legende{Cross-sections of guided media: twisted pair, coaxial cable, optical fibre.}
\end{popfigure}

\courssub{Skill 7: Selecting an appropriate unguided (wireless) transmission medium}

\begin{definition}[Unguided Transmission Media]
Signals propagate through free space (air, vacuum) as electromagnetic waves.
\begin{itemize}
\item \cle{Infrared}: Wavelength $700\ \mathrm{nm}$--$1\ \mathrm{mm}$. Short range (few metres), requires \textbf{line of sight}, cannot penetrate walls. Low power, cheap. Uses: remote controls, IrDA (obsolete).
\item \cle{Radio waves}: Wavelength $1\ \mathrm{mm}$--$100\ \mathrm{km}$. Omnidirectional, penetrates walls (with attenuation). Range: metres (Bluetooth, Wi-Fi) to kilometres (cellular 2G/3G/4G/5G). Used for mobile networks, Wi-Fi, broadcasting.
\item \cle{Microwaves}: High-frequency radio ($1\ \mathrm{GHz}$--$300\ \mathrm{GHz}$). Highly \textbf{directional}, requires \textbf{line of sight} between parabolic dishes. Range: tens of km per hop. Used for terrestrial point-to-point links (backhaul between base stations).
\item \cle{Satellite}: Microwave link via geostationary satellite ($\approx 36\,000\ \mathrm{km}$ altitude). Continental/global coverage. High latency ($\approx 250\ \mathrm{ms}$ one-way). Used for remote areas, broadcasting, VSAT.
\end{itemize}
\end{definition}

\begin{methode}[Choosing between radio waves, microwaves, infrared and satellite]
\begin{enumerate}
\item \textbf{Check the range and the obstacles.} Infrared: few metres, line of sight, no walls. Radio (Wi-Fi, cellular): metres to km, penetrates walls. Microwave: tens of km, line of sight, directional. Satellite: global, independent of ground infrastructure.
\item \textbf{Check whether line of sight exists.} No line of sight $\Rightarrow$ eliminate infrared and terrestrial microwave.
\item \textbf{Check mobility.} Mobile users $\Rightarrow$ cellular radio (2G/3G/4G) or Wi-Fi, not directional microwave.
\item \textbf{Justify} with range, directionality and obstacle behaviour.
\end{enumerate}
\end{methode}

\begin{exoresolu}[Matching wireless media to Cameroonian situations]
Choose the most appropriate unguided medium for each case, with justification:
\begin{enumerate}
\item A teacher uses a remote control to change slides on a projector.
\item Students connect their laptops to the Internet inside a school campus covered by a hotspot.
\item A mobile operator links two base stations on two hills $40\ \mathrm{km}$ apart, with a clear view between them.
\item Internet access must be provided to a health centre in an isolated area of the Far North Region where no cable or cellular coverage exists.
\end{enumerate}
\tcblower
\textbf{1. Remote control: infrared.} The distance is a few metres, the remote is pointed at the projector (line of sight available), and the power required is tiny. Infrared is cheap and perfect for this; its inability to cross walls is even an advantage, as it avoids interfering with the projector in the next room.

\textbf{2. Campus hotspot: radio waves (Wi-Fi).} Wi-Fi uses radio waves broadcast in all directions, covering rooms and corridors without line of sight, since radio waves cross walls (with attenuation). It supports many mobile users simultaneously, which matches students moving with laptops.

\textbf{3. Hill-to-hill operator link: terrestrial microwave.} The two antennas are in line of sight and $40\ \mathrm{km}$ apart, within the range of a microwave hop. Microwaves are highly directional: a focused beam between the two dishes gives a high-capacity, point-to-point link with little interference. This is exactly how operators interconnect distant base stations where laying fibre would be too costly.

\textbf{4. Isolated health centre: satellite (VSAT).} No terrestrial infrastructure (cable, fibre, cellular) reaches the site, and distances to the nearest connected town are large. A satellite link provides coverage anywhere under the satellite's footprint, independently of ground infrastructure. The trade-offs are higher cost and noticeable propagation delay, but it is the only viable unguided solution here.

\textbf{Summary:}
\begin{center}
\begin{tabularx}{\linewidth}{|l|X|X|}
\hline
\rowcolor{popblueL} \textbf{Medium} & \textbf{Range / constraint} & \textbf{Typical use} \\
\hline
Infrared & Few metres, line of sight, no walls & Remote controls, IrDA \\
\hline
Radio (Wi-Fi, Bluetooth, cellular) & Metres to km, penetrates walls & Hotspots, mobile phones, IoT \\
\hline
Microwave (terrestrial) & Tens of km, line of sight, directional & Backhaul between base stations \\
\hline
Satellite (GEO) & Continental / global & Remote areas, broadcasting, VSAT \\
\hline
\end{tabularx}
\end{center}
\end{exoresolu}

\begin{popfigure}
\begin{tikzpicture}[scale=0.8, every node/.style={font=\small}]
% Infrared
\draw[poppink, thick, ->] (0,3) -- (1.5,3) node[midway, above] {Infrared: few m, LOS};
% Wi-Fi
\draw[popblue, thick] (0,2) circle (1.5);
\node at (0,2) {AP};
\foreach \angle in {0,45,...,315} {
  \draw[popblue, ->] (0,2) -- ++(\angle:1.5);
}
\node at (2.5,2) {Wi-Fi: omnidirectional, penetrates walls};
% Microwave
\draw[poporange, thick, ->] (0,1) -- (4,1) node[midway, above] {Microwave: directional, LOS, 40 km};
\draw[poporange, thick] (0,1) -- (0.5,1.3) (0,1) -- (0.5,0.7);
\draw[poporange, thick] (4,1) -- (3.5,1.3) (4,1) -- (3.5,0.7);
% Satellite
\draw[poppurple, thick] (0,0) -- (2,2) node[midway, above left, sloped] {Uplink};
\draw[poppurple, thick] (2,2) -- (4,0) node[midway, above right, sloped] {Downlink};
\node[draw, circle, fill=poppurple!20] at (2,2) {Satellite};
\node at (2,-0.5) {Satellite: global coverage, high latency};
\end{tikzpicture}
\legende{Unguided media characteristics: infrared (short, LOS), Wi-Fi (omnidirectional), microwave (directional LOS), satellite (global).}
\end{popfigure}

\courssub{Skill 8: Analysing a complete transmission scenario (synthesis)}

\begin{methode}[Answering a multi-part GCE scenario on transmission]
\begin{enumerate}
\item \textbf{Read the scenario twice} and underline every technical clue: distances, bit rates, ``at the same time'', ``one after another'', ``wireless'', ``start/stop bits'', etc.
\item \textbf{Answer each sub-question with the pattern:} name the concept $\rightarrow$ define it briefly $\rightarrow$ apply it to the scenario $\rightarrow$ justify.
\item \textbf{Reuse the vocabulary of the chapter precisely:} broadband/narrowband, simplex/half-duplex/full-duplex, serial/parallel, synchronous/asynchronous, guided/unguided.
\item \textbf{Manage your time:} definitions are quick marks; calculations (bit rate, time, efficiency) must show the formula, the substitution and the unit.
\end{enumerate}
\end{methode}

\begin{exoresolu}[GCE-style synthesis question]
The ministry in charge of secondary education connects a regional delegation to the central database in Yaound\'e through a fibre-optic link of $10\ \mathrm{Mbit/s}$. Inside the delegation, staff computers send data to a server character by character with start and stop bits; the link between the delegation and Yaound\'e carries data in continuous frames, in both directions at the same time, and simultaneously carries telephone (VoIP) conversations and video conferences on separate frequency channels. A technician also transfers a file of $20\ \mathrm{MB}$ to Yaound\'e.
\begin{enumerate}
\item State whether the fibre link is narrowband or broadband, justifying with two observations from the text.
\item Give the transmission mode (simplex, half-duplex or full-duplex) between the delegation and Yaound\'e, with justification.
\item Identify the type of transmission (synchronous or asynchronous) used: (i) inside the delegation; (ii) on the fibre link.
\item State whether the fibre link is a guided or unguided medium, and give two advantages of fibre over copper cable.
\item Calculate the theoretical time to transfer the $20\ \mathrm{MB}$ file.
\end{enumerate}
\tcblower
\textbf{1. Broadband.} The fibre link is \cle{broadband} because: (i) it offers a high bit rate ($10\ \mathrm{Mbit/s}$); (ii) it carries several signals simultaneously (data, VoIP, video) on separate frequency channels, which is the defining feature of broadband.

\textbf{2. Full-duplex.} The text states that data flows ``in both directions at the same time'' between the delegation and Yaound\'e. Simultaneous two-way communication is the definition of \cle{full-duplex} transmission.

\textbf{3. Synchronous or asynchronous?}
\begin{itemize}
\item (i) Inside the delegation: data is sent character by character with start and stop bits framing each character $\Rightarrow$ \cle{asynchronous} transmission.
\item (ii) On the fibre link: data is sent in continuous frames with synchronised clocks $\Rightarrow$ \cle{synchronous} transmission.
\end{itemize}

\textbf{4. Guided medium, advantages of fibre.} Optical fibre is a \cle{guided} medium: the signal (light) is confined and propagated along a physical cable. Two advantages over copper cable: (i) much higher bandwidth and bit rate over much longer distances without repeaters; (ii) total immunity to electromagnetic interference (and it is also difficult to tap, hence more secure).

\textbf{5. Transfer time.} Convert the file size to bits:
\[ 20\ \mathrm{MB} = 20 \times 10^6 \times 8 = 1.6 \times 10^8\ \mathrm{bits}. \]
Apply $t = \dfrac{\text{size}}{\text{bit rate}}$:
\[ t = \dfrac{1.6 \times 10^8}{10 \times 10^6} = 16\ \mathrm{s}. \]
The theoretical transfer time is $16$ seconds, assuming the full $10\ \mathrm{Mbit/s}$ is available to the file (in reality, VoIP and video share the link, so the time would be longer).

\textbf{Examiner's remark.} Each mark here follows the same pattern: correct term, brief definition, justification drawn from the text. Reproduce this pattern and you secure full credit on synthesis questions of this chapter.
\end{exoresolu}

\begin{aretenir}[Chapter ICT08 -- Key Revision Points]
\begin{itemize}
\item \textbf{Bandwidth vs bit rate:} Bandwidth in Hz, bit rate in bit/s. Convert prefixes (k, M, G) using powers of 10.
\item \textbf{Narrowband:} one signal, low bit rate (voice-grade). \textbf{Broadband:} many signals (multiplexed), high bit rate.
\item \textbf{Transmission modes:} Simplex (1 way), Half-duplex (2 ways alternate), Full-duplex (2 ways simultaneous).
\item \textbf{Serial vs Parallel:} Serial = 1 line, sequential bits (long distance). Parallel = $k$ lines, simultaneous bits (short distance, skew issues).
\item \textbf{Async vs Sync:} Async = start/stop bits per char, independent clocks, $\approx 80\%$ efficiency. Sync = frames, common clock, $\approx 100\%$ efficiency.
\item \textbf{Guided media:} UTP (short, cheap), Coaxial (TV, medium), Fibre (long, high speed, immune to EMI).
\item \textbf{Unguided media:} Infrared (short, LOS), Radio/Wi-Fi (omni, walls), Microwave (directional, LOS, backhaul), Satellite (global, remote).
\item \textbf{Exam technique:} Define $\rightarrow$ Apply $\rightarrow$ Justify. Show formulas and units in calculations.
\end{itemize}
\end{aretenir}

\begin{exercice}[Practice Questions]
\begin{enumerate}
\item A file of $5\ \mathrm{MB}$ is sent over a link of $4\ \mathrm{Mbit/s}$. Calculate the theoretical transfer time.
\item Classify each as narrowband or broadband: (a) Dial-up at $56\ \mathrm{kbit/s}$; (b) 4G mobile at $20\ \mathrm{Mbit/s}$; (c) ADSL at $8\ \mathrm{Mbit/s}$.
\item For each, state the transmission mode: (a) TV broadcast; (b) Walkie-talkie; (c) WhatsApp voice call.
\item A device sends $8$-bit characters asynchronously with $1$ start bit and $2$ stop bits. If the line runs at $9\,600\ \mathrm{bit/s}$, what is the useful data rate in bit/s?
\item A company must link two offices $500\ \mathrm{m}$ apart in an industrial zone with heavy machinery. Choose the best guided medium and justify.
\item A farmer in a remote village needs Internet for market prices. No mobile coverage. Which unguided medium? Justify.
\item \textbf{Synthesis:} A university campus uses a fibre backbone ($1\ \mathrm{Gbit/s}$) between buildings. Inside a lab, PCs connect via UTP Cat6 to a switch. The backbone carries student data, VoIP, and CCTV video simultaneously. Lab PCs send keystrokes asynchronously. Answer: (a) Backbone: broadband or narrowband? (b) Backbone transmission mode? (c) Lab PC to switch: async or sync? (d) Backbone: guided or unguided? (e) Time to send $500\ \mathrm{MB}$ over backbone at full rate?
\end{enumerate}
\end{exercice}

\begin{corrige}[Exercise 1]
$5\ \mathrm{MB} = 5 \times 10^6 \times 8 = 4 \times 10^7\ \mathrm{bits}$. $t = \frac{4 \times 10^7}{4 \times 10^6} = 10\ \mathrm{s}$.
\end{corrige}

\begin{corrige}[Exercise 2]
(a) Narrowband (single voice channel, low bit rate). (b) Broadband (high bit rate, multiple services). (c) Broadband (high bit rate, multiplexed).
\end{corrige}

\begin{corrige}[Exercise 3]
(a) Simplex. (b) Half-duplex. (c) Full-duplex.
\end{corrige}

\begin{corrige}[Exercise 4]
Total bits/char = $1+8+2 = 11$. Efficiency = $8/11 \approx 72.7\%$. Useful rate = $9\,600 \times 8/11 \approx 6\,982\ \mathrm{bit/s}$.
\end{corrige}

\begin{corrige}[Exercise 5]
Optical fibre (single-mode). Justification: $500\ \mathrm{m}$ exceeds UTP $100\ \mathrm{m}$ limit; industrial machinery causes strong EMI, fibre is immune; high bit rate easily supported.
\end{corrige}

\begin{corrige}[Exercise 6]
Satellite (VSAT). Justification: remote area, no terrestrial infrastructure, satellite provides coverage anywhere.
\end{corrige}

\begin{corrige}[Exercise 7]
(a) Broadband (high bit rate, multiple simultaneous services). (b) Full-duplex (data flows both ways simultaneously). (c) Asynchronous (keystrokes sent char by char with start/stop bits). (d) Guided (fibre is a physical cable). (e) $500\ \mathrm{MB} = 4 \times 10^9\ \mathrm{bits}$. $t = \frac{4 \times 10^9}{1 \times 10^9} = 4\ \mathrm{s}$.
\end{corrige}

\newpage

\coursec{Exercises}

\courssub{I check my knowledge}

\begin{exercice}[MCQ: Bandwidth and Transmission Types]
Choose the correct answer for each question.

\begin{enumerate}
\item Which of the following best defines \emph{bandwidth} in data communication?
\begin{itemize}
\item[A.] The physical width of the transmission medium.
\item[B.] The difference between the highest and lowest frequencies that a channel can carry.
\item[C.] The total amount of data transmitted in one second.
\item[D.] The speed of light in the medium.
\end{itemize}

\item Narrowband transmission is typically used for:
\begin{itemize}
\item[A.] High-definition video streaming.
\item[B.] Voice-grade telephone lines.
\item[C.] Fibre-optic backbone networks.
\item[D.] Satellite broadband internet.
\end{itemize}

\item In Cameroon, an example of a broadband service is:
\begin{itemize}
\item[A.] Traditional fixed-line telephone (PSTN).
\item[B.] MTN 4G LTE mobile internet.
\item[C.] A walkie-talkie set.
\item[D.] A simple radio broadcast receiver.
\end{itemize}

\item The unit commonly used to measure bandwidth is:
\begin{itemize}
\item[A.] Metres per second (m/s).
\item[B.] Bits per second (bps) or Hertz (Hz).
\item[C.] Volts (V).
\item[D.] Bytes (B).
\end{itemize}
\end{enumerate}
\end{exercice}

\begin{exercice}[True/False with Justification]
State whether each statement is True or False. Justify your answer briefly.

\begin{enumerate}
\item Broadband transmission always implies a wireless connection.
\item Half-duplex communication allows simultaneous two-way data flow.
\item Asynchronous transmission uses start and stop bits to synchronise the receiver for each character.
\item Optical fibre is immune to electromagnetic interference.
\end{enumerate}
\end{exercice}

\begin{exercice}[Definitions and Direct Calculations]
\begin{enumerate}
\item Define the term \emph{bandwidth} and state its two common units.
\item Distinguish between \emph{narrowband} and \emph{broadband} transmission. Give one Cameroonian example of each.
\item A voice-grade telephone line has a bandwidth of 3.1 kHz. If it uses a modulation scheme that achieves 2 bits per Hz, calculate the maximum theoretical data rate in kbps.
\item A network link operates at 100 Mbps. How many megabytes (MB) can be transferred in one minute? (Assume 1 byte = 8 bits, and ignore overhead.)
\end{enumerate}
\end{exercice}

\begin{exercice}[Transmission Modes Identification]
For each scenario below, identify the transmission mode (simplex, half-duplex, or full-duplex) and justify your choice.

\begin{enumerate}
\item A secretary uses a walkie-talkie to talk to a colleague in another office.
\item The same secretary makes a mobile phone call to a client.
\item A family listens to CRTV radio broadcast on their FM receiver.
\item A computer sends data to a printer via a USB cable (assume bidirectional USB communication).
\end{enumerate}
\end{exercice}

\courssub{I practise}

\begin{exercice}[Comparison Table: Simplex, Half-Duplex, Full-Duplex]
Copy and complete the following table in your answer booklet.

\begin{center}
\begin{tabular}{|l|c|c|c|}
\hline
\textbf{Characteristic} & \textbf{Simplex} & \textbf{Half-Duplex} & \textbf{Full-Duplex} \\ \hline
Direction of data flow & & & \\ \hline
Simultaneity (can both ends transmit at the same time?) & & & \\ \hline
Typical example & & & \\ \hline
Number of communication channels required & & & \\ \hline
\end{tabular}
\end{center}
\end{exercice}

\begin{exercice}[Asynchronous vs Synchronous Transmission]
\begin{enumerate}
\item Explain why asynchronous transmission adds a start bit and one or more stop bits to each character.
\item State one advantage of asynchronous transmission over synchronous transmission.
\item State one disadvantage of asynchronous transmission compared with synchronous transmission.
\item A synchronous link transmits frames of 1000 bytes each. If the frame overhead (header + trailer) is 48 bytes, calculate the overhead percentage.
\end{enumerate}
\end{exercice}

\begin{exercice}[Transmission Media Selection Scenario]
A company in Douala needs to connect two buildings located 5 km apart across a busy highway. The area has high electrical interference from industrial machinery. The required data rate is at least 1 Gbps, and the budget is moderate.

\begin{enumerate}
\item Recommend the most suitable guided transmission medium for this link.
\item Justify your recommendation with at least three technical reasons.
\item Briefly explain why unguided media (e.g., microwave or satellite) would be less appropriate in this specific situation.
\end{enumerate}
\end{exercice}

\begin{exercice}[Guided Media: Coaxial Cable and Optical Fibre]
\begin{enumerate}
\item Draw a labelled cross-sectional diagram of a coaxial cable. Indicate the core, dielectric insulator, braided shield, and outer jacket.
\item Draw a labelled cross-sectional diagram of a step-index optical fibre. Indicate the core, cladding, and protective coating.
\item State two advantages of optical fibre over copper-based media (UTP, STP, coaxial) for long-distance, high-bandwidth links.
\end{enumerate}
\end{exercice}

\courssub{I prepare for the GCE}

\begin{exercice}[Structured Essay: Bandwidth, Transmission Types, and Media] (20 marks)
\begin{enumerate}
\item[a)] Define the term \emph{bandwidth} in the context of data communication. State its unit. (2 marks)
\item[b)] Distinguish clearly between \emph{narrowband} and \emph{broadband} transmission. Give one practical example of each in the Cameroonian context. (4 marks)
\item[c)] Describe the differences between \emph{simplex}, \emph{half-duplex}, and \emph{full-duplex} transmission modes. For each mode, provide a real-life example. (6 marks)
\item[d)] A school in Yaoundé wants to connect its administrative block to the computer laboratory 800 m away. The path passes near a high-voltage power line. The school requires a reliable 100 Mbps link at low cost.
\begin{itemize}
\item[i)] Recommend a suitable guided transmission medium. (1 mark)
\item[ii)] Justify your choice with three reasons, considering the environment and requirements. (3 marks)
\item[iii)] Explain why unshielded twisted pair (UTP) would not be the best choice here. (2 marks)
\end{itemize}
\item[e)] Briefly explain the role of start and stop bits in asynchronous serial transmission. State one advantage and one disadvantage of asynchronous transmission relative to synchronous transmission. (2 marks)
\end{enumerate}
\end{exercice}

\begin{exercice}[Multiple-Choice Quiz: Comprehensive] (10 marks)
Choose the correct option for each question.

\begin{enumerate}
\item Baseband transmission:
\begin{itemize}
\item[A.] Uses multiple frequencies simultaneously.
\item[B.] Sends digital signals directly over the medium without modulation.
\item[C.] Is only used in wireless networks.
\item[D.] Requires a modem for every connection.
\end{itemize}

\item Which of the following is a characteristic of broadband transmission?
\begin{itemize}
\item[A.] It uses the entire bandwidth of the medium for a single signal.
\item[B.] It allows multiple signals to share the medium via frequency division multiplexing.
\item[C.] It is limited to distances under 100 m.
\item[D.] It cannot carry analogue signals.
\end{itemize}

\item In serial transmission:
\begin{itemize}
\item[A.] Multiple bits are sent simultaneously over parallel wires.
\item[B.] Bits are sent one after another over a single line.
\item[C.] It is always slower than parallel transmission.
\item[D.] It requires a separate clock line for each bit.
\end{itemize}

\item Unshielded Twisted Pair (UTP) Category 6 cable:
\begin{itemize}
\item[A.] Has a maximum recommended segment length of 100 m for Gigabit Ethernet.
\item[B.] Is immune to electromagnetic interference.
\item[C.] Uses a braided metal shield around each pair.
\item[D.] Is the preferred medium for outdoor long-distance links.
\end{itemize}

\item Shielded Twisted Pair (STP) differs from UTP by:
\begin{itemize}
\item[A.] Having a metal foil or braid around the pairs to reduce EMI.
\item[B.] Using optical fibres instead of copper.
\item[C.] Supporting only half-duplex communication.
\item[D.] Having a larger maximum segment length (200 m).
\end{itemize}

\item Microwave line-of-sight links:
\begin{itemize}
\item[A.] Require the transmitting and receiving antennas to be in direct visual contact.
\item[B.] Are unaffected by rain or atmospheric conditions.
\item[C.] Have very low latency, similar to fibre optics.
\item[D.] Use frequencies below 30 MHz.
\end{itemize}

\item Optical fibre transmits data using:
\begin{itemize}
\item[A.] Electrical signals over copper.
\item[B.] Light pulses through a glass or plastic core.
\item[C.] Radio waves in free space.
\item[D.] Magnetic fields.
\end{itemize}

\item A key advantage of single-mode fibre over multi-mode fibre is:
\begin{itemize}
\item[A.] Lower cost per metre.
\item[B.] Larger core diameter for easier splicing.
\item[C.] Much higher bandwidth over longer distances.
\item[D.] Compatibility with LED light sources.
\end{itemize}

\item Satellite communication links typically suffer from:
\begin{itemize}
\item[A.] Very low latency (under 10 ms).
\item[B.] High latency (around 250–500 ms for geostationary satellites).
\item[C.] Complete immunity to weather.
\item[D.] Unlimited bandwidth.
\end{itemize}

\item Which statement about transmission media is FALSE?
\begin{itemize}
\item[A.] Coaxial cable has better shielding than UTP.
\item[B.] Optical fibre is immune to electromagnetic interference.
\item[C.] Twisted pair cables use differential signalling to reduce noise.
\item[D.] Wireless media provide higher security than guided media because signals are confined.
\end{itemize}
\end{enumerate}
\end{exercice}

\begin{exercice}[True/False with Justification – GCE Style] (10 marks)
For each statement, write TRUE or FALSE and give a concise justification.

\begin{enumerate}
\item Broadband transmission means the connection is wireless.
\item Half-duplex communication allows both devices to transmit and receive simultaneously.
\item Satellite links have very low latency, making them ideal for real-time gaming.
\item Asynchronous transmission is more efficient than synchronous transmission for large blocks of data.
\item Optical fibre suffers from signal attenuation due to electromagnetic interference.
\item Coaxial cable has a higher bandwidth than twisted pair but lower than optical fibre.
\item Microwave transmission requires line-of-sight between antennas.
\item UTP Category 5e cable can support 1 Gbps over 100 m.
\item In full-duplex mode, only one communication channel is needed.
\item The start bit in asynchronous transmission is always a logic 1 (mark).
\end{enumerate}
\end{exercice}



\newpage

\coursec{Solutions to the exercises}

\courssub{I check my knowledge}

\begin{corrige}[Exercise 1 — MCQ: Bandwidth and Transmission Types]
\begin{enumerate}
\item \textbf{B.} Bandwidth is the \emph{difference between the highest and lowest frequencies} a channel can carry, $B = f_{max} - f_{min}$, measured in Hertz. Option C describes the \emph{data rate} (throughput), which is a consequence of bandwidth, not its definition.
\item \textbf{B.} Voice-grade telephone lines (PSTN) use a narrow channel of about $3.1\ \mathrm{kHz}$ (300 Hz to 3.4 kHz) — the classic narrowband application. Options A, C and D all require broadband capacity.
\item \textbf{B.} MTN 4G LTE offers high data rates (several Mbps), which is broadband. PSTN is narrowband; a walkie-talkie and a simple FM receiver use narrow radio channels.
\item \textbf{B.} Analogue bandwidth is measured in \textbf{Hertz (Hz)}; digital channel capacity is expressed in \textbf{bits per second (bps)}. Bytes measure storage, volts measure electrical potential, m/s measures speed.
\end{enumerate}
\end{corrige}

\begin{corrige}[Exercise 2 — True/False with Justification]
\begin{enumerate}
\item \textbf{False.} Broadband describes a \emph{wide bandwidth / high capacity} (often with several channels multiplexed by FDM), independent of the medium. Broadband exists in wired form (fibre-to-the-home, DSL, cable TV) and in wireless form (4G LTE, WiMAX).
\item \textbf{False.} Half-duplex allows communication in \emph{both directions but alternately} — only one party transmits at a time (e.g., walkie-talkie). Simultaneous two-way flow is the definition of \emph{full-duplex}.
\item \textbf{True.} In asynchronous transmission there is no shared clock, so each character is framed by a \emph{start bit} (logic 0) which triggers the receiver's clock, and one or more \emph{stop bits} (logic 1) which return the line to idle and allow resynchronisation for the next character.
\item \textbf{True.} Optical fibre carries \emph{light}, not electric current, through a dielectric (glass/plastic). External electromagnetic fields therefore cannot induce noise in the signal — a major advantage over copper media in industrial environments.
\end{enumerate}
\end{corrige}

\begin{corrige}[Exercise 3 — Definitions and Direct Calculations]
\begin{enumerate}
\item \textbf{Bandwidth} is the range of frequencies that a transmission channel can carry, i.e., the difference between the highest and the lowest usable frequency: $B = f_{max} - f_{min}$. \textbf{Units:} Hertz (Hz) for analogue channels; bits per second (bps) when expressing digital capacity.
\item \textbf{Narrowband:} transmission over a small frequency range, supporting low data rates (voice-grade, typically $\le 64\ \mathrm{kbps}$). \emph{Cameroonian example:} CAMTEL fixed-line (PSTN) voice service. \textbf{Broadband:} transmission over a wide frequency range, supporting high data rates and/or several simultaneous channels. \emph{Cameroonian example:} MTN or Orange 4G LTE mobile internet, or CAMTEL fibre-to-the-home.
\item Maximum data rate $=$ bandwidth $\times$ spectral efficiency:
\[ D = 3100\ \mathrm{Hz} \times 2\ \mathrm{bits/Hz} = 6200\ \mathrm{bps} = \mathbf{6.2\ kbps}. \]
\item Bits transferred in one minute:
\[ 100 \times 10^{6}\ \mathrm{bps} \times 60\ \mathrm{s} = 6 \times 10^{9}\ \mathrm{bits}. \]
Converting to megabytes ($1\ \mathrm{MB} = 8 \times 10^{6}$ bits):
\[ \frac{6 \times 10^{9}}{8 \times 10^{6}} = \mathbf{750\ MB}. \]
\end{enumerate}
\end{corrige}

\begin{corrige}[Exercise 4 — Transmission Modes Identification]
\begin{enumerate}
\item \textbf{Half-duplex.} A walkie-talkie works in ``push-to-talk'' mode: the same frequency channel is shared, so only one person can speak at a time while the other listens.
\item \textbf{Full-duplex.} During a mobile phone call both parties can speak and hear each other at the same time; the network provides independent paths (separate frequencies or time slots) for each direction.
\item \textbf{Simplex.} CRTV's FM broadcast flows in one direction only — from the transmitter to the receivers; listeners cannot send anything back on that channel.
\item \textbf{Full-duplex.} USB provides separate differential pairs for the two directions, so the computer and the printer can exchange data (print data one way, status/acknowledgements the other) simultaneously.
\end{enumerate}
\end{corrige}

\courssub{I practise}

\begin{corrige}[Exercise 5 — Comparison Table: Simplex, Half-Duplex, Full-Duplex]
\begin{center}
\begin{tabular}{|l|c|c|c|}
\hline
\textbf{Characteristic} & \textbf{Simplex} & \textbf{Half-Duplex} & \textbf{Full-Duplex} \\ \hline
Direction of data flow & One direction only & Both directions, alternately & Both directions, simultaneously \\ \hline
Simultaneity & No & No & Yes \\ \hline
Typical example & FM radio broadcast, keyboard $\to$ CPU & Walkie-talkie, legacy Ethernet hub & Telephone call, switched Ethernet \\ \hline
Number of channels required & 1 & 1 (shared) & 2 (or 1 with FDM/TDD) \\ \hline
\end{tabular}
\end{center}
\end{corrige}

\begin{corrige}[Exercise 6 — Asynchronous vs Synchronous Transmission]
\begin{enumerate}
\item In asynchronous transmission the sender and receiver do \emph{not} share a common clock. The \textbf{start bit} (logic 0) warns the receiver that a character is arriving and starts its local clock for sampling that character. The \textbf{stop bit(s)} (logic 1) return the line to the idle state and guarantee a minimum gap so the receiver can resynchronise independently on the next character.
\item \textbf{Advantage:} the hardware is simple and cheap (no clock line or complex clock-recovery circuitry), and it copes well with irregular, bursty data such as keystrokes.
\item \textbf{Disadvantage:} the framing overhead is high — typically 2 to 3 extra bits per 8-bit character, i.e., about 20--30\% of the line capacity is wasted, so effective throughput is lower than with synchronous transmission.
\item Overhead percentage:
\[ \text{Overhead} = \frac{48}{1000 + 48} \times 100 = \frac{48}{1048} \times 100 \approx \mathbf{4.58\%}. \]
This low overhead illustrates why synchronous transmission is preferred for large blocks of data.
\end{enumerate}
\end{corrige}

\begin{corrige}[Exercise 7 — Transmission Media Selection Scenario]
\begin{enumerate}
\item \textbf{Recommended medium: single-mode optical fibre} (e.g., OS2, 9/125 $\mu$m), installed in ducts or on aerial poles between the two buildings.
\item \textbf{Justification (three technical reasons):}
\begin{itemize}
\item \textbf{Immunity to electromagnetic interference:} fibre is a dielectric and carries light, so the heavy industrial EMI in the area cannot corrupt the signal.
\item \textbf{Bandwidth and distance:} single-mode fibre carries 1 Gbps (and far more, up to 100 Gbps) over 5 km with very low attenuation (about 0.4 dB/km at 1310 nm) and no repeater.
\item \textbf{Security and reliability:} fibre radiates nothing (difficult to tap), suffers no crosstalk, and provides galvanic isolation between the two buildings' electrical earths.
\end{itemize}
\item \textbf{Why unguided media are less appropriate here:}
\begin{itemize}
\item \textbf{Microwave:} needs strict line-of-sight (the busy highway, buildings and vegetation can obstruct the Fresnel zone), requires masts on both roofs, suffers rain fade, and frequency licensing through the regulator (ANTIC) adds cost and delay.
\item \textbf{Satellite:} introduces high latency (about 500 ms round-trip for geostationary links), high recurring subscription costs, and is technically absurd for a mere 5 km link.
\end{itemize}
\end{enumerate}
\end{corrige}

\begin{corrige}[Exercise 8 — Guided Media: Coaxial Cable and Optical Fibre]
\begin{enumerate}
\item \textbf{Coaxial cable cross-section:}
\begin{popfigure}
\begin{tikzpicture}[scale=1.1]
\fill[popblueL] (0,0) circle (2.0);
\fill[gray] (0,0) circle (1.7);
\fill[poporange!30] (0,0) circle (1.2);
\fill[popgold] (0,0) circle (0.5);
\draw[<-, popdark] (0.35,0.35) -- (2.6,1.6) node[right, text width=3cm, align=left] {Inner conductor (copper core)};
\draw[<-, popdark] (0.9,-0.8) -- (2.6,-1.4) node[right, text width=3cm, align=left] {Dielectric insulator (PE/PTFE)};
\draw[<-, popdark] (-1.2,1.2) -- (-2.6,1.8) node[left, text width=3cm, align=right] {Braided shield (copper)};
\draw[<-, popdark] (-1.4,-1.4) -- (-2.6,-1.9) node[left, text width=3cm, align=right] {Outer jacket (PVC)};
\end{tikzpicture}
\legende{Cross-section of a coaxial cable: core, dielectric, braided shield and jacket.}
\end{popfigure}
\item \textbf{Step-index optical fibre cross-section:}
\begin{popfigure}
\begin{tikzpicture}[scale=2.2]
\fill[popblueL] (0,0) circle (0.8);
\fill[popgreenL] (0,0) circle (0.5);
\fill[poppurple] (0,0) circle (0.2);
\draw[<-, popdark] (0.14,0.14) -- (1.1,0.8) node[right, text width=3cm, align=left] {Core (glass, high refractive index)};
\draw[<-, popdark] (-0.35,0.35) -- (-1.1,0.9) node[left, text width=3cm, align=right] {Cladding (glass, lower index)};
\draw[<-, popdark] (-0.55,-0.55) -- (-1.1,-1.0) node[left, text width=3cm, align=right] {Protective coating (acrylate)};
\end{tikzpicture}
\legende{Cross-section of a step-index optical fibre: core, cladding and coating.}
\end{popfigure}
\item \textbf{Two advantages of optical fibre over copper media:}
\begin{enumerate}
\item \textbf{Far higher bandwidth over far longer distances:} single-mode fibre carries tens of Gbps over tens of kilometres without repeaters, whereas copper (e.g., Cat6 UTP) is limited to about 100 m at 1 Gbps.
\item \textbf{Total immunity to electromagnetic interference and crosstalk:} fibre carries light, not current, so it can run beside power lines or industrial machinery without degradation, and it emits no signal that could be intercepted.
\end{enumerate}
\end{enumerate}
\end{corrige}

\courssub{I prepare for the GCE}

\begin{corrige}[Exercise 9 — Structured Essay: Bandwidth, Transmission Types, and Media]
\textbf{Indicative mark scheme and model answer (20 marks).}

\begin{enumerate}
\item[a)] \textbf{Bandwidth} is the range of frequencies a communication channel can transmit, i.e., the difference between the highest and lowest frequencies: $B = f_{max} - f_{min}$. \textbf{Unit:} Hertz (Hz) — equivalently bits per second (bps) for digital capacity. \emph{(1 mark for the definition, 1 mark for the unit.)}
\item[b)] \textbf{Narrowband:} transmission over a small frequency range, supporting low data rates (voice-grade, typically $\le 64\ \mathrm{kbps}$); a single channel occupies the medium. \emph{Example:} CAMTEL PSTN fixed-line telephone. \textbf{Broadband:} transmission over a wide frequency range, supporting high data rates and/or several simultaneous channels (by frequency-division multiplexing). \emph{Example:} MTN/Orange 4G LTE mobile internet or CAMTEL FTTH. \emph{(1 mark per correct distinction element, 1 mark per valid Cameroonian example — 4 marks.)}
\item[c)] \textbf{Simplex:} data flows in one direction only; the receiver cannot reply on that channel. \emph{Example:} CRTV FM radio broadcast. \textbf{Half-duplex:} data flows in both directions but alternately — only one party transmits at a time. \emph{Example:} walkie-talkie (push-to-talk). \textbf{Full-duplex:} data flows in both directions simultaneously, using two channels (or FDM/TDD on one medium). \emph{Example:} an ordinary mobile phone call. \emph{(1 mark per correct description + 1 mark per valid example — 6 marks.)}
\item[d)]
\begin{itemize}
\item[i)] \textbf{Recommended: optical fibre} (multi-mode OM3/OM4 or single-mode) run in a protective duct. \emph{(1 mark.)}
\item[ii)] \textbf{Three reasons:} (1) \emph{EMI immunity:} the link passes near a high-voltage line; fibre, being dielectric, is completely unaffected by the induced electromagnetic noise. (2) \emph{Distance and speed:} 800 m far exceeds the 100 m copper Ethernet limit, while fibre easily carries 100 Mbps (and much more) over 800 m without any repeater. (3) \emph{Safety/isolation:} fibre provides galvanic isolation between the two buildings, eliminating ground-loop and surge risks. \emph{(3 marks — 1 per justified reason.)}
\item[iii)] \textbf{UTP is unsuitable because:} (1) it has no shielding, so the high-voltage line would induce severe noise, causing errors and retransmissions; (2) its maximum Ethernet segment length is 100 m, so 800 m would require several powered intermediate switches/repeaters — raising cost, power needs and failure points. \emph{(2 marks.)}
\end{itemize}
\item[e)] The \textbf{start bit} (logic 0) announces the arrival of a character and synchronises the receiver's clock; the \textbf{stop bit(s)} (logic 1) return the line to idle and guarantee a gap for resynchronisation. \textbf{Advantage:} simple, cheap hardware, no shared clock, ideal for bursty data. \textbf{Disadvantage:} high overhead (about 20--30\%), so lower effective throughput than synchronous transmission. \emph{(1 mark for the role of the bits, 1 mark for a valid advantage + disadvantage — 2 marks.)}
\end{enumerate}

\textbf{Examiner expectations:} precise definitions with correct units; examples that are genuinely Cameroonian and correctly classified; justifications in (d) that explicitly refer to the EMI environment and the 800 m distance; clear, technical vocabulary throughout.
\end{corrige}

\begin{corrige}[Exercise 10 — Multiple-Choice Quiz: Comprehensive]
\begin{enumerate}
\item \textbf{B.} Baseband sends digital signals directly on the medium without modulation, using the whole channel for one signal (e.g., Ethernet).
\item \textbf{B.} Broadband divides the medium's capacity into multiple channels by frequency-division multiplexing (e.g., cable TV, DSL, 4G).
\item \textbf{B.} Serial transmission sends bits one after another over a single line. (It is not always slower than parallel — modern high-speed serial links outperform parallel buses, which suffer skew.)
\item \textbf{A.} Cat6 UTP supports Gigabit Ethernet up to 100 m. It has no shielding (C is false) and is not immune to EMI (B is false).
\item \textbf{A.} STP adds a metallic foil or braid around the pairs (or the whole cable) to reduce electromagnetic interference; its segment length remains 100 m.
\item \textbf{A.} Microwave links require clear line-of-sight (with Fresnel-zone clearance) between the antennas; they use GHz frequencies and are affected by rain fade.
\item \textbf{B.} Optical fibre carries data as light pulses (typically at 850, 1310 or 1550 nm) propagating in a glass or plastic core by total internal reflection.
\item \textbf{C.} Single-mode fibre has a very small core (about 9 $\mu$m) and uses laser sources, eliminating modal dispersion — hence much higher bandwidth over much longer distances than multi-mode fibre.
\item \textbf{B.} Geostationary satellites orbit at about 35,786 km, giving a round-trip latency of roughly 480--500 ms; they are also affected by rain fade.
\item \textbf{D.} This statement is false: wireless signals propagate openly through space and are easily intercepted, so wireless media are \emph{less} secure than guided media, where signals are confined inside cables.
\end{enumerate}
\emph{Mark scheme: 1 mark per correct option — 10 marks.}
\end{corrige}

\begin{corrige}[Exercise 11 — True/False with Justification (GCE Style)]
\begin{enumerate}
\item \textbf{False.} Broadband describes wide bandwidth / high capacity (often multi-channel via FDM); it exists in wired form (fibre, DSL, cable) as well as wireless (4G).
\item \textbf{False.} Half-duplex allows two-way communication but only \emph{alternately}; simultaneous transmission and reception is full-duplex.
\item \textbf{False.} Geostationary satellite links have high latency (about 250--500 ms round-trip), which is poor for real-time gaming; fibre offers a few milliseconds.
\item \textbf{False.} For large blocks, synchronous transmission is more efficient: one clock synchronisation serves the whole block, whereas asynchronous adds 2--3 framing bits to every character (20--30\% overhead).
\item \textbf{False.} Fibre is immune to EMI. Its attenuation comes from absorption, scattering and bending of light in the glass — not from electromagnetic interference.
\item \textbf{True.} Typical usable bandwidths: twisted pair up to a few hundred MHz, coaxial cable up to a few GHz, optical fibre in the tens of THz range.
\item \textbf{True.} Microwave beams are highly directional and cannot diffract around obstacles, so the two antennas must see each other (with Fresnel-zone clearance).
\item \textbf{True.} Category 5e UTP is specified for 1000BASE-T (1 Gbps) over channels up to 100 m.
\item \textbf{False.} Full-duplex needs two independent transmission paths — two wire pairs/fibres, two frequency bands (FDD) or two time slots (TDD) — one for each direction.
\item \textbf{False.} The start bit is a logic \textbf{0} (space); the idle line and the stop bit(s) are logic 1 (mark).
\end{enumerate}
\emph{Mark scheme: 0.5 mark for the correct verdict + 0.5 mark for a correct, concise justification per item — 10 marks. Examiner expectation: a verdict without justification earns only half the marks.}
\end{corrige}

\newpage

\coursec{Summary sheet}

\begin{popfigure}
\begin{tikzpicture}[
  central/.style={circle, draw=popdark, fill=popblue, text=white, font=\bfseries, minimum size=1.3cm},
  branch/.style={rectangle, rounded corners, draw=popdark, fill=popblueL, text=popdark, font=\small, minimum width=2.6cm, minimum height=1.1cm, align=center},
  line/.style={thick, draw=popline}
]
\node[central, align=center] (C){ICT08:\\Transmission Types};
% Branches
\node[branch, fill=poporange, align=center] (B1) at (0:3.8cm){Fundamentals\\ \& Bandwidth};
\node[branch, fill=popgreen, align=center] (B2) at (60:3.8cm){Broadband vs\\ Narrowband};
\node[branch, fill=poppurple, align=center] (B3) at (120:3.8cm){Data Transmission\\ Types};
\node[branch, fill=poppink] (B4) at (180:3.8cm) {Guided Media};
\node[branch, fill=popgold] (B5) at (240:3.8cm) {Unguided Media};
\node[branch, fill=popmuted] (B6) at (300:3.8cm) {Key Formulas};
\node[branch, fill=popblueL] (B7) at (360:3.8cm) {GCE Tips};
% Connections
\foreach \n in {B1,B2,B3,B4,B5,B6,B7}
  \draw[line] (C) -- (\n);
\end{tikzpicture}
\legende{Mind map of Chapter ICT08 -- Transmission Types}
\end{popfigure}

\begin{bilan}
\textbf{Key Definitions}
\begin{itemize}
  \item \cle{Bandwidth}: range of frequencies a channel can carry, measured in hertz (Hz). Determines maximum data rate (Nyquist/Shannon).
  \item \cle{Broadband}: transmission using a wide frequency range, allowing multiple simultaneous channels (e.g. MTN 4G, CAMTEL fibre).
  \item \cle{Narrowband}: transmission confined to a narrow frequency band, typically a single channel (e.g. traditional PSTN voice line).
  \item \cle{Simplex}, \cle{Half‑duplex}, \cle{Full‑duplex}: directionality of data flow.
  \item \cle{Serial} vs \cle{Parallel} transmission: bits sent one after another vs several bits simultaneously.
  \item \cle{Synchronous} vs \cle{Asynchronous} transmission: clock‑controlled continuous stream vs start/stop bits per character.
  \item \cle{Guided media}: physical cables (twisted pair, coaxial, fibre optic).
  \item \cle{Unguided media}: wireless propagation (radio, microwave, satellite, infrared).
\end{itemize}

\textbf{Laws / Formulas to Know}
\begin{itemize}
  \item \cle{Nyquist bit rate}: $C = 2B\log_2 M$ (noise‑free channel, $M$ signal levels).
  \item \cle{Shannon capacity}: $C = B\log_2(1+\mathrm{SNR})$ (real channel, SNR in linear scale).
  \item \cle{Attenuation}: $A_{\mathrm{dB}} = 10\log_{10}\frac{P_{\text{in}}}{P_{\text{out}}}$.
  \item \cle{Propagation speed}: $v = f\lambda$; in fibre $v \approx 2\times10^8\ \mathrm{m/s}$.
  \item \cle{Data rate vs bandwidth}: for binary signalling, max baud rate $\approx 2B$ (Nyquist).
\end{itemize}

\textbf{Methods / Procedures}
\begin{itemize}
  \item \textbf{Choosing a medium}: compare bandwidth, distance, cost, immunity to EMI, installation constraints. Example: fibre for backbone (high bandwidth, low loss), twisted pair for LAN (cost‑effective), microwave for rural links (no trenching).
  \item \textbf{Calculating capacity}: identify $B$ and SNR, apply Shannon formula; convert SNR from dB: $\mathrm{SNR}_{\text{lin}} = 10^{\mathrm{SNR}_{\text{dB}}/10}$.
  \item \textbf{Identifying transmission mode}: check if both ends transmit simultaneously (full‑duplex), alternate (half‑duplex), or one‑way only (simplex).
  \item \textbf{Serial vs parallel}: parallel used for short distances (e.g. PCI bus), serial for long distances (USB, Ethernet, fibre).
\end{itemize}

\textbf{Common Mistakes (Traps)}
\begin{itemize}
  \item Confusing \emph{bandwidth} (Hz) with \emph{data rate} (bit/s); they are related but not equal.
  \item Using Nyquist formula on a noisy channel — must use Shannon.
  \item Forgetting to convert SNR from dB to linear before Shannon.
  \item Thinking broadband always means “high speed”; it means \emph{wide frequency range} (multiple channels).
  \item Mixing up \emph{guided} and \emph{unguided} examples: satellite is unguided, coaxial is guided.
  \item Assuming full‑duplex requires two physical cables; a single fibre with WDM can be full‑duplex.
\end{itemize}

\textbf{GCE Examination Tips}
\begin{itemize}
  \item \textbf{Define} terms precisely (bandwidth, broadband, narrowband, simplex, half‑duplex, full‑duplex, serial, parallel, synchronous, asynchronous).
  \item \textbf{Calculate} channel capacity: show conversion of SNR dB $\to$ linear, then apply Shannon.
  \item \textbf{Compare} guided vs unguided media in a table: bandwidth, distance, cost, interference, typical use (Cameroon context: fibre backbone, MTN/Orange 4G/5G, CAMTEL copper, VSAT for remote areas).
  \item \textbf{Draw} a simple block diagram of a communication system (source $\to$ transmitter $\to$ medium $\to$ receiver $\to$ destination) and label where bandwidth limitation occurs.
  \item \textbf{Explain} why fibre is preferred for long‑haul: low attenuation ($\approx0.2\ \mathrm{dB/km}$), immunity to EMI, huge bandwidth.
  \item \textbf{Watch} for “explain the difference between broadband and baseband” — baseband uses the whole medium for one signal (e.g. Ethernet), broadband multiplexes many signals (e.g. cable TV, 4G).
  \item \textbf{Practice} past GCE questions: 2019 Q3 (bandwidth vs data rate), 2021 Q5 (Shannon capacity), 2023 Q2 (guided/unguided media comparison).
\end{itemize}
\end{bilan}

\findecours

\end{document}
